\documentclass[11pt,a4paper]{article}

% ============================================================
% PACKAGES
% ============================================================

\usepackage[utf8]{inputenc}
\usepackage[T1]{fontenc}
\usepackage[
    a4paper,
    top=2.5cm,
    bottom=2.3cm,
    left=2.2cm,
    right=2.2cm,
    headheight=24pt,
    headsep=10mm
]{geometry}

\usepackage{lmodern}
\usepackage{xcolor}
\usepackage{hyperref}
\usepackage{enumitem}
\usepackage{parskip}
\usepackage{fancyhdr}
\usepackage{lastpage}
\usepackage{titlesec}
\usepackage{booktabs}
\usepackage{tabularx}
\usepackage{array}
\usepackage[most]{tcolorbox}
\usepackage{microtype}
\usepackage{listings}

% ============================================================
% COLORS
% ============================================================

\definecolor{AaltoNavy}{HTML}{0B1F4D}
\definecolor{AaltoBlue}{HTML}{0067B9}
\definecolor{LightBlue}{HTML}{EAF4FB}
\definecolor{SoftGray}{HTML}{F4F6F8}
\definecolor{MediumGray}{HTML}{667085}
\definecolor{DarkText}{HTML}{1D2939}
\definecolor{SuccessGreen}{HTML}{16866B}
\definecolor{WarningOrange}{HTML}{D97706}
\definecolor{BorderGray}{HTML}{D0D5DD}
\definecolor{CodeBackground}{HTML}{F7F9FC}

% ============================================================
% PDF METADATA
% ============================================================

\hypersetup{
    colorlinks=true,
    linkcolor=AaltoBlue,
    urlcolor=AaltoBlue,
    pdftitle={
        Aalto EE Net Sales Forecasting Case --
        Implementation Progress Report
    },
    pdfauthor={Arman Golbidi},
    pdfsubject={
        Power Query Transformation,
        Actual Sales Validation and
        Accrual and CRM Opportunity Validation,
        Current-Year Weighted Controls and
        Cross-Year Allocation Validation
    },
    pdfkeywords={
        Aalto EE,
        Net Sales Forecast,
        Power BI,
        Power Query,
        Actual Sales,
        Accruals,
        Data Quality,
        Opportunity Reconciliation,
        Current-Year Weighted Value,
        Allocation Validation,
        Semantic Model
    }
}

% ============================================================
% GENERAL STYLE
% ============================================================

\setlength{\parindent}{0pt}
\setlength{\parskip}{6pt}
\setlength{\emergencystretch}{2em}
\renewcommand{\arraystretch}{1.25}
\setcounter{tocdepth}{2}

\titleformat{\section}
    {\Large\bfseries\color{AaltoNavy}}
    {}
    {0pt}
    {}

\titleformat{\subsection}
    {\large\bfseries\color{AaltoBlue}}
    {}
    {0pt}
    {}

\titlespacing*{\section}{0pt}{12pt}{5pt}
\titlespacing*{\subsection}{0pt}{9pt}{4pt}

% ============================================================
% HEADER AND FOOTER
% ============================================================

\pagestyle{fancy}
\fancyhf{}

\fancyhead[L]{
    \textcolor{AaltoNavy}{
        \textbf{AALTO EE FORECASTING CASE}
    }
}

\fancyhead[R]{
    \textcolor{MediumGray}{
        Implementation Progress Report
    }
}

\fancyfoot[L]{
    \small\textcolor{MediumGray}{
        Arman Golbidi
    }
}

\fancyfoot[C]{
    \small\textcolor{MediumGray}{
        Data Analyst/Engineer Interview
    }
}

\fancyfoot[R]{
    \small\textcolor{MediumGray}{
        Page \thepage\ of \pageref*{LastPage}
    }
}

\renewcommand{\headrulewidth}{0.6pt}
\renewcommand{\footrulewidth}{0.4pt}

% ============================================================
% CUSTOM BOXES
% ============================================================

\newtcolorbox{summarybox}{
    colback=LightBlue,
    colframe=AaltoBlue,
    boxrule=0.8pt,
    arc=2mm,
    left=4mm,
    right=4mm,
    top=3mm,
    bottom=3mm
}

\newtcolorbox{successbox}{
    colback=white,
    colframe=SuccessGreen,
    boxrule=0.8pt,
    arc=2mm,
    left=4mm,
    right=4mm,
    top=3mm,
    bottom=3mm
}

\newtcolorbox{warningbox}{
    colback=SoftGray,
    colframe=WarningOrange,
    boxrule=0.8pt,
    arc=2mm,
    left=4mm,
    right=4mm,
    top=3mm,
    bottom=3mm
}

\newtcolorbox{inferencebox}{
    colback=SoftGray,
    colframe=AaltoNavy,
    boxrule=0.8pt,
    arc=2mm,
    left=4mm,
    right=4mm,
    top=3mm,
    bottom=3mm
}

% ============================================================
% CODE STYLE
% ============================================================

\lstdefinestyle{powerquery}{
    backgroundcolor=\color{CodeBackground},
    basicstyle=\ttfamily\small,
    frame=single,
    rulecolor=\color{BorderGray},
    breaklines=true,
    columns=fullflexible,
    showstringspaces=false,
    xleftmargin=3mm,
    xrightmargin=3mm,
    aboveskip=6pt,
    belowskip=6pt
}

% ============================================================
% DOCUMENT
% ============================================================

\begin{document}

% ============================================================
% TITLE
% ============================================================

{
    \Huge\bfseries\color{AaltoNavy}
    Net Sales Forecasting Case
    \par
}

\vspace{2mm}

{
    \Large\color{AaltoBlue}
    Implementation Progress Report
    \par
}

\vspace{2mm}

{
    \large\color{MediumGray}
    Actual Sales, Accrual and CRM Financial Controls Completed
    \par
}

\vspace{7mm}

\begin{summarybox}

\begin{tabularx}{\textwidth}{
    >{\bfseries\color{AaltoNavy}}p{4.3cm}
    X
}

Candidate
&
Arman Golbidi
\\

Position
&
Data Analyst/Engineer
\\

Organization
&
Aalto University Executive Education and Professional Development
\\

Solution Platform
&
Microsoft Power BI Desktop and Power Query
\\

Reporting Period
&
Financial year 2026
\\

Phase 1 Status
&
Actual net sales transformation and quality validation completed
\\

Phase 2 Status
&
Accrual reconciliation, quality and duplicate controls completed
\\

Phase 3 Status
&
CRM preparation, quality classification, weighted-value reconciliation,
opportunity-value reconciliation, current-year weighted-value controls and
cross-year allocation validation completed
\\

Current-Year Status Summary
&
Completed: 85 Match, 39 Not Testable -- Missing Current-Year Allocation,
and 2 Not Testable -- Missing Probability
\\

Allocation Validation
&
Completed: 126 records classified into 6 allocation statuses with
0 errors and 0 empty classifications
\\

Allocation Status Summary
&
Completed: 6 grouped statuses reconcile exactly to 126 records
\\

Allocation Exceptions
&
Completed: 4 non-100\% records isolated for targeted review
\\

Missing-Allocation Horizon Analysis
&
Completed: 52 records reconciled; 41 (78.8\%) close in 2026
\\

2026 Missing-Allocation Detail
&
Completed: 41 near-term records isolated for direct review
\\

2026 Missing-Allocation Exposure
&
Completed: 41 records represent EUR 4.209M estimated value and
EUR 2.318M probability-weighted value
\\

2026 Current-Year Allocation Risk
&
Completed: 30 records represent EUR 3.572M estimated value and
EUR 1.989M probability-weighted value
\\

2026 Forecast Readiness
&
Completed: 106 opportunities classified; 75 Ready, 28 Missing Current-Year
Allocation, 2 Missing Probability and 1 Cross-Year Allocation Review
\\

2026 Forecast Readiness Status Summary
&
Completed: Ready current-year contribution EUR 807,900.00; one review
amount EUR 13,500.00; unresolved weighted pipeline EUR 1,989,178.99
\\

2026 CRM Forecast Bridge
&
Completed: 4 governed buckets keep EUR 807,900 Ready forecast,
EUR 13,500 Review value, EUR 1,989,178.99 weighted risk and
EUR 50,000 missing-Probability exposure separate
\\

Business Unit Source Audit
&
Completed: Actual Sales uses 3 shorter labels; CRM and Accrual share
3 canonical BU labels; 4 Accrual rows remain unassigned
\\

Business Unit Dimension
&
Completed: \texttt{DimBusinessUnit} contains 3 canonical Business Units plus
1 explicit \texttt{Unassigned / Missing BU} technical member
\\

Actual Sales BU Key Mapping
&
Completed: all 795 Actual Sales rows mapped to BusinessUnitKey 1--3 with
100\% valid values, 0 errors and 0 empty keys
\\

Project Accrual BU Key Mapping
&
Completed: all 313 Accrual rows mapped with 100\% valid values,
0 errors, 0 empty keys and 4 distinct keys including key 0
\\

CRM Opportunity BU Key Mapping
&
Completed: all 126 Opportunity rows mapped with 100\% valid values,
0 errors, 0 empty keys and 3 distinct keys
\\

Semantic Model Load Scope
&
Completed: staging and DQ queries remain Power Query-only; the visible semantic
model now contains governed Business Unit, Date, CRM Date and Project dimensions,
three analytical facts and one intentionally disconnected Forecast Bucket table
\\

Business Unit Relationships
&
Completed: 3 active one-to-many relationships from \texttt{DimBusinessUnit}
to the three facts using single-direction filtering
\\

Business Unit Relationship Validation
&
Completed: ExEd 346/162/92, Qualification 120/43/7,
University 329/104/27 and Unassigned 0/4/0 across Actual/Accrual/CRM
\\

Date Dimension
&
Completed: 365 unique daily dates for 2026 with 12 months,
12 YearMonth values, 12 MonthStart values and 4 quarters
\\

Date Relationships
&
Completed: two active many-to-one relationships connect Actual Sales
\texttt{MonthStart} and Accrual \texttt{ForecastMonth} to \texttt{DimDate[Date]}
with Single cross-filter direction
\\

Date Relationship Validation
&
Completed: January--August contains 795 Actual rows and no Accrual rows;
September--December contains 313 Accrual rows and no Actual rows
\\

Month Sorting
&
Completed: \texttt{DimDate[MonthName]} now displays chronologically from
January through December using \texttt{MonthNumber}
\\

Base Financial Measures
&
Completed: Actual Net Sales EUR 14,635,067.81; Project Accrual Forecast
EUR 6,616,025.78; Base Full-Year Forecast EUR 21,251,093.59
\\

Base Forecast Monthly Validation
&
Completed: Actual Net Sales populates January--August and Project Accrual
Forecast populates September--December on one shared 2026 Date axis
\\

Date-Table Designation
&
Completed: \texttt{DimDate} is formally marked as the model Date table using
\texttt{DimDate[Date]}
\\

Business Unit Financial Validation
&
Completed: the three base financial measures reconcile by Business Unit to
EUR 21,251,093.59, including EUR 42,550.00 in the explicit Unassigned member
\\

CRM Ready Measures
&
Completed: 75 Ready opportunities contribute EUR 807,900.00 of reconciled
current-year CRM value; BU split is 60/2/13 across ExEd/Qualification/University
\\

Forecast Including Ready CRM
&
Completed: Base Full-Year Forecast EUR 21,251,093.59 plus Ready CRM
EUR 807,900.00 equals EUR 22,058,993.59
\\

Semantic CRM Governance Measures
&
Completed: Review EUR 13,500.00; At-Risk weighted pipeline
EUR 1,989,178.99; Missing-Probability exposure EUR 50,000.00; review-approved
scenario EUR 22,072,493.59
\\

Forecast Bucket Dimension
&
Completed: disconnected \texttt{DimForecastBucket} presents five governed
forecast treatments without adding another fact-filter path
\\

CRM Calendar
&
Completed: \texttt{DimCRMDate} covers 1 January 2024 through
31 December 2028 with 1,827 unique daily dates and 0 errors
\\

CRM Date-Role Relationships
&
Completed: active \texttt{ClosingDate} relationship and inactive
\texttt{SimulatedProgramStartDate} relationship, both Single direction
\\

CRM Date-Role DAX Validation
&
Completed: both date roles reconcile to all 126 opportunities; yearly distributions
are 1/1, 1/1, 106/56, 15/64 and 3/4 for Closing/Program Start across 2024--2028
\\

Project Dimension Profile
&
Completed in Power Query: \texttt{DimProject} contains 252 rows with
252 distinct and 252 unique ProjectIDs, 0 errors and 0 empty values
\\

Current Project-Dimension Status
&
Completed: \texttt{DimProject} is governed under \texttt{04\_Dimensions} and
filters both project-based facts through active Single-direction relationships
\\

Project Semantic Relationships
&
Completed: \texttt{DimProject[ProjectID]} is on the one side of active
many-to-one relationships from \texttt{Fact\_ActualSales[ProjectID]} and
\texttt{Fact\_ProjectAccrual[ProjectID]}
\\

Project Relationship Validation
&
Completed: 795 Actual rows and 313 Accrual rows reconcile to EUR 14,635,067.81
Actual, EUR 6,616,025.78 Accrual and EUR 21,251,093.59 Base Full-Year Forecast
\\

Next Control
&
Perform one final interactive Project filter test, then move into the final
executive dashboard refinement and Project-level forecast ranking visuals
\\


Project Filter Interaction Validation
&
Completed: selecting \texttt{O7901009} filters the project-based Actual and Accrual
measures to EUR 439,132.89 and EUR 308,913.00 respectively, producing
EUR 748,045.89 Base Full-Year Forecast; CRM remains independent by design
\\

Top-10 Project Forecast Ranking
&
Completed: an executive horizontal bar chart now ranks the ten largest projects by
\texttt{Base Full-Year Forecast} using \texttt{DimProject[ProjectID]}
\\

Current Dashboard Next Control
&
Build the Top-10 Project forecast-composition visual so the ranking also shows the
split between recognized Actual Net Sales and remaining Project Accrual Forecast
\\


Top-10 Project Forecast Composition
&
Completed: the companion stacked horizontal chart now shows
\texttt{Actual Net Sales} versus \texttt{Project Accrual Forecast} for the
Top-10 projects selected by \texttt{Base Full-Year Forecast}
\\

Executive Dashboard KPI Layer
&
Completed: eight EUR-formatted KPI cards now surface Actual, Accrual, Base Forecast,
Forecast Including Ready CRM, Ready CRM, Review CRM, At-Risk Pipeline and
Missing-Probability exposure
\\

Executive Dashboard Business Unit Slicer
&
Completed: \texttt{DimBusinessUnit[BusinessUnitName]} is exposed as a governed
four-member report slicer
\\

Executive Monthly Forecast Visual
&
Added and month order is correct, but one presentation correction remains:
the current visual is a 100\% stacked column chart and must be changed to a
standard stacked column chart so bar height represents monetary value rather
than percentage share
\\

Current Dashboard Next Control -- Updated
&
Correct the monthly chart type, validate the Business Unit slicer interaction,
then add the Business Unit forecast-composition visual to the Executive Dashboard
\\


Executive Monthly Forecast Correction -- Status Update
&
Completed: the monthly visual now uses a standard stacked column chart with a
monetary EUR axis; January--August shows Actual Net Sales and
September--December shows Project Accrual Forecast
\\

Business Unit Slicer Interactive Validation
&
Completed: all four governed Business Unit members were tested and the financial,
CRM-governance and monthly visuals respond consistently through the shared
\texttt{DimBusinessUnit} filter path
\\

Current Dashboard Next Control -- Latest
&
Build the Business Unit forecast-composition visual using the validated
Actual, Accrual and Base Full-Year Forecast measures
\\



Business Unit Forecast Composition
&
Completed: horizontal stacked bar compares Actual Net Sales and remaining Project
Accrual Forecast across the four governed Business Unit members
\\

Forecast Governance Table
&
Completed: five governed forecast buckets are displayed in the intended business
order with amount and treatment kept explicit
\\

Executive Dashboard Composition Checkpoint
&
Completed: eight KPI cards, Business Unit slicer, Forecast Governance table,
Business Unit composition, monthly forecast and Top-10 Project composition are
present on the executive page
\\

Current Dashboard Next Control -- Final Interaction Validation
&
Re-run the Business Unit slicer against the newly added Governance, Business Unit
and Top-10 Project visuals, then complete final presentation cleanup
\\



Final Semantic Relationship Audit
&
Completed: 9 model relationships were re-checked in \texttt{Manage relationships}.
Eight are active; the \texttt{SimulatedProgramStartDate} CRM role is intentionally
inactive. All dimensional filter paths remain many-to-one from fact to dimension
with Single-direction propagation
\\

Executive Dashboard End-to-End BU Interaction Validation
&
Completed: ExEd, Qualification Programs, University Programs and
\texttt{Unassigned / Missing BU} were re-tested against KPI cards, Forecast Governance,
Business Unit composition, monthly forecast and Top-10 Project composition visuals.
All visuals respond consistently to the shared Business Unit slicer
\\

Technical Field Visibility Cleanup
&
Completed: technical keys, sort helpers, row-quality fields and diagnostic columns
needed for model governance remain in the semantic model but are hidden from the
business-facing field list where appropriate
\\

Current Next Control -- Final Financial Reconciliation Sign-Off
&
Create the final reconciliation control showing that
\texttt{Base Full-Year Forecast = Actual Net Sales + Project Accrual Forecast} and
\texttt{Forecast Including Ready CRM = Base Full-Year Forecast + CRM Ready Current-Year Value},
then complete presentation polish and the interview narrative
\\

\end{tabularx}

\end{summarybox}

\clearpage

\tableofcontents

\clearpage

% ============================================================
% OBJECTIVE
% ============================================================

\section{1. Objective}

The objective is to establish a controlled, repeatable and auditable
foundation for the Aalto EE monthly net sales forecasting solution.

The implementation completed so far includes:

\begin{itemize}[leftmargin=7mm,itemsep=3pt]
    \item importing the three provided Excel data sources;
    \item separating raw staging queries from analytical fact queries;
    \item preparing and validating the actual net sales dataset;
    \item creating valid monthly date fields for actuals and accruals;
    \item implementing row-level data-quality controls;
    \item validating the project-month grain of actual sales;
    \item separating the accrual source Total from analytical detail;
    \item retaining the Total in a dedicated reconciliation query;
    \item classifying the quality of all 313 accrual detail records;
    \item validating the accrual detail sum against the source Total;
    \item identifying and quantifying exact accrual duplicates;
    \item identifying repeated project-month business keys;
    \item counting distinct forecast amounts within every repeated key;
    \item classifying repeated keys into three interpretable patterns;
    \item profiling all 126 CRM opportunity records across 14 source columns;
    \item assigning explicit Power Query data types to all 14 CRM fields;
    \item confirming zero CRM data-type conversion errors;
    \item standardizing the four CRM text fields using Trim and Clean;
    \item assigning clear analytical names to all 14 CRM source fields;
    \item classifying the quality of all 126 CRM opportunity records;
    \item confirming zero errors and zero empty CRM quality classifications;
    \item independently recalculating probability-weighted opportunity values;
    \item reconciling calculated and source weighted values using a EUR 0.01
    tolerance;
    \item identifying one material weighted-value mismatch for investigation;
    \item reconciling \texttt{OpportunityValue} against \texttt{EstimatedValue};
    \item confirming that the opportunity-value check has one distinct result,
    \texttt{Match}, across all 126 CRM records;
    \item independently calculating the current-year probability-weighted amount
    using estimated value, probability and the current-year allocation percentage;
    \item calculating the source-minus-calculated
    \texttt{CurrentYearWeightedVariance};
    \item correcting the Power Query step order so that the variance is created
    before the final \texttt{CurrentYearWeightedCheck};
    \item confirming zero errors and zero empty values in
    \texttt{CurrentYearWeightedCheck}; and
    \item confirming that the current-year check contains three distinct statuses;
    \item creating \texttt{DQ\_Opportunity\_CurrentYearWeightedStatus} as a dedicated
    summary query without changing \texttt{Fact\_Opportunities};
    \item confirming the exact current-year status distribution: 85 \texttt{Match},
    39 \texttt{Not Testable - Missing Current-Year Allocation}, and
    2 \texttt{Not Testable - Missing Probability}; and
    \item confirming that all 85 fully testable current-year records reconcile within
    the EUR 0.01 tolerance, with no \texttt{Mismatch} status present;
    \item creating \texttt{DQ\_Opportunity\_AllocationValidation} as a separate
    reference query without changing \texttt{Fact\_Opportunities};
    \item calculating \texttt{AllocationTotalPct} only when both annual allocation
    inputs are present;
    \item classifying all 126 opportunity records into six observed allocation
    validation statuses with 0 errors and 0 empty classifications;
    \item confirming that 74 records have both annual allocation percentages present,
    while 52 records are missing one or both allocation inputs; and
    \item identifying 70 records with a 100\% two-year allocation total, 3 records
    below 100\%, and 1 record above 100\% for targeted business review;
    \item creating \texttt{DQ\_Opportunity\_AllocationStatusSummary} and independently
    confirming the six allocation-status counts sum to all 126 source records;
    \item creating \texttt{DQ\_Opportunity\_AllocationExceptions} to isolate the four
    opportunities whose two-year allocation total is not 100\%;
    \item identifying one 110\% over-allocation for \texttt{Open custom opportunity 111}
    in University Programs BU; and
    \item isolating three below-100\% cases with totals of 50\%, 50\% and 30\% for
    date-aware business interpretation rather than automatic correction;
    \item creating \texttt{DQ\_Opportunity\_MissingAllocationByYear} and reconciling
    all 52 missing-allocation records by allocation status, closing year and
    simulated program-start year;
    \item confirming that 41 of the 52 missing-allocation records (78.8\%) have
    \texttt{ClosingYear = 2026}, so missing allocation is not primarily a distant-horizon issue;
    \item creating \texttt{DQ\_Opportunity\_2026MissingAllocationDetail} to isolate
    the 41 near-term records without changing \texttt{Fact\_Opportunities}; and
    \item confirming the 2026 closing-year composition as 21 records with both
    allocations missing, 9 missing only the current-year allocation and
    11 missing only the next-year allocation;
    \item creating \texttt{DQ\_Opportunity\_2026MissingAllocationExposure} to quantify
    the financial value represented by the 41 near-term allocation gaps;
    \item confirming EUR 4,208,699.23 of estimated value and EUR 2,317,982.74 of
    probability-weighted value across those 41 records;
    \item creating \texttt{DQ\_Opportunity\_2026CurrentYearAllocationRisk} to isolate
    the 30 2026-closing opportunities whose current-year allocation is missing;
    \item confirming EUR 3,572,367.98 of estimated value and EUR 1,989,178.99 of
    probability-weighted value in that current-year allocation risk set; and
    \item confirming that ExEd Programs BU contains the largest share of this
    weighted risk, EUR 1,229,613.75 across 18 records;
    \item creating \texttt{DQ\_Opportunity\_2026ForecastReadinessSummary} to classify
    all 106 CRM opportunities closing in 2026 by forecast readiness and Business Unit;
    \item confirming 75 \texttt{Ready - Reconciled Current-Year Value} records,
    28 \texttt{Not Ready - Missing Current-Year Allocation} records,
    2 \texttt{Not Ready - Missing Probability} records and
    1 \texttt{Review - Cross-Year Allocation Above 100\%} record;
    \item confirming that the 75 Ready opportunities contribute
    EUR 807,900.00 of reconciled current-year CRM value;
    \item creating \texttt{DQ\_Opportunity\_2026ForecastReadinessStatusSummary} to
    collapse the Business Unit detail into four auditable forecast-readiness buckets; and
    \item preserving EUR 1,989,178.99 of unresolved weighted pipeline and
    EUR 50,000.00 of missing-probability estimated value outside the Ready forecast
    instead of silently converting them into current-year revenue;
    \item creating \texttt{DQ\_Opportunity\_2026ForecastBridge} and separating the
    four forecast treatments into dedicated financial columns;
    \item confirming the governed bridge values as EUR 807,900.00 Ready current-year
    forecast, EUR 13,500.00 held for review, EUR 1,989,178.99 at-risk weighted
    pipeline and EUR 50,000.00 estimated pipeline with missing Probability;
    \item creating \texttt{DQ\_BusinessUnit\_SourceValues} to compare Business Unit
    vocabularies across Actual Sales, Project Accrual and CRM Opportunities;
    \item confirming that Actual Sales uses \texttt{ExEd},
    \texttt{Qualification Programs} and \texttt{University Programs}, while CRM and
    Accrual use the corresponding \texttt{... Programs BU} labels; and
    \item confirming that 4 accrual records still have a missing Business Unit and
    must remain explicitly unassigned until business validation;
    \item creating the \texttt{04\_Dimensions} query group and the canonical
    \texttt{DimBusinessUnit} dimension with four governed members;
    \item assigning Business Unit keys 1, 2 and 3 to the three populated canonical
    Business Units and key 0 to \texttt{Unassigned / Missing BU};
    \item adding \texttt{BusinessUnitKey} to \texttt{Fact\_ActualSales} while
    preserving the original \texttt{BusinessUnitSource} field for lineage; and
    \item confirming that all 795 Actual Sales rows have a valid non-null key:
    346 rows map to key 1, 120 rows to key 2 and 329 rows to key 3;
    \item adding \texttt{BusinessUnitKey} to \texttt{Fact\_ProjectAccrual} while
    preserving \texttt{BusinessUnitSource};
    \item confirming that all 313 Accrual rows have valid keys, including the
    4 missing-BU rows assigned explicitly to key 0;
    \item adding \texttt{BusinessUnitKey} to \texttt{Fact\_Opportunities} while
    preserving the original CRM Business Unit label;
    \item confirming that all 126 Opportunity rows have valid keys across the three
    canonical Business Units with no key-0 Opportunity records;
    \item disabling model load for staging and data-quality queries while retaining
    their Power Query logic and refresh lineage;
    \item reducing the visible semantic model to \texttt{DimBusinessUnit},
    \texttt{Fact\_ActualSales}, \texttt{Fact\_ProjectAccrual} and
    \texttt{Fact\_Opportunities}; and
    \item creating three active single-direction Business Unit relationships with
    the dimension on the one side and each fact on the many side;
    \item validating Business Unit filter propagation with three DAX row-count
    measures and confirming exact reconciliation to 795 Actual, 313 Accrual and
    126 CRM Opportunity rows;
    \item confirming the Business Unit split as 346/162/92 for ExEd,
    120/43/7 for Qualification Programs, 329/104/27 for University Programs and
    0/4/0 for the Unassigned member across Actual/Accrual/CRM respectively; and
    \item creating \texttt{DimDate} for the full 2026 calendar with 365 unique
    daily dates, 12 months, 12 YearMonth values, 12 MonthStart values and
    4 quarters, with zero errors and zero empty dates;
    \item moving \texttt{DimDate} into the governed dimension layer;
    \item creating active Single-direction Date relationships from
    \texttt{DimDate[Date]} to \texttt{Fact\_ActualSales[MonthStart]} and
    \texttt{Fact\_ProjectAccrual[ForecastMonth]}; and
    \item validating monthly Date-filter propagation across both facts, confirming
    795 Actual rows in January--August and 313 Accrual rows in
    September--December with no cross-period leakage;
    \item sorting \texttt{DimDate[MonthName]} chronologically by
    \texttt{MonthNumber};
    \item creating the DAX measures \texttt{Actual Net Sales},
    \texttt{Project Accrual Forecast} and \texttt{Base Full-Year Forecast}; and
    \item reconciling the completed base full-year forecast to
    EUR 14,635,067.81 of Actual Net Sales plus EUR 6,616,025.78 of Project
    Accrual Forecast, producing EUR 21,251,093.59 before CRM contribution;
    \item formally marking \texttt{DimDate} as the model Date table using
    \texttt{DimDate[Date]};
    \item validating the three base financial measures by Business Unit and
    reconciling the BU rows exactly to EUR 21,251,093.59;
    \item creating \texttt{CRM Ready Opportunity Count} and
    \texttt{CRM Ready Current-Year Value} measures and confirming 75 Ready
    opportunities with EUR 807,900.00 of current-year value; and
    \item creating \texttt{Forecast Including Ready CRM} and reconciling the
    current governed annual forecast to EUR 22,058,993.59;
    \item creating and validating the semantic CRM governance measures for
    EUR 13,500.00 of Review value, EUR 1,989,178.99 of At-Risk weighted pipeline
    and EUR 50,000.00 of Missing-Probability estimated exposure;
    \item creating the optional \texttt{Forecast If Review Approved} scenario and
    reconciling it to EUR 22,072,493.59;
    \item creating the disconnected \texttt{DimForecastBucket} presentation
    dimension to expose governed forecast treatments without adding an ambiguous
    relationship path;
    \item creating \texttt{DimCRMDate} for the complete 2024--2028 CRM horizon
    with 1,827 unique daily dates and zero errors;
    \item creating one active CRM Date relationship on \texttt{ClosingDate} and
    one inactive role-playing relationship on \texttt{SimulatedProgramStartDate},
    both using Single cross-filter direction;
    \item validating both CRM Date roles with DAX and confirming that each role
    reconciles to all 126 opportunity records while producing different yearly
    distributions;
    \item creating \texttt{DimProject} from Actual Sales and Project Accrual
    ProjectIDs; and
    \item confirming that \texttt{DimProject} contains 252 rows, 252 distinct and
    252 unique ProjectIDs, with zero errors and zero empty values;
    \item moving \texttt{DimProject} into the governed \texttt{04\_Dimensions}
    query group with model load retained;
    \item creating two active Single-direction Project relationships directly on
    the clean natural key \texttt{ProjectID}, with \texttt{DimProject} on the one side
    and Actual Sales / Project Accrual on the many side;
    \item validating Project-level filter propagation and reconciling the relationship
    layer to 795 Actual rows, 313 Accrual rows, EUR 14,635,067.81 of Actual Net Sales,
    EUR 6,616,025.78 of Project Accrual Forecast and EUR 21,251,093.59 of Base
    Full-Year Forecast; and
    \item retaining \texttt{ProjectKey} in \texttt{DimProject} as a technical surrogate
    identifier while deliberately avoiding an unnecessary fact-table merge because
    the supplied fact \texttt{ProjectID} values already reconcile cleanly to the unique
    dimension key.
    \item completing the final interactive Project slicer validation with
    \texttt{O7901009}, confirming EUR 439,132.89 Actual Net Sales,
    EUR 308,913.00 Project Accrual Forecast and EUR 748,045.89 Base Full-Year Forecast;
    \item confirming that the Project slicer does not create an unsupported CRM-to-Project
    filter path and that CRM forecast-governance measures remain independent; and
    \item creating the Top-10 Project ranking chart using \texttt{DimProject[ProjectID]}
    and \texttt{Base Full-Year Forecast}, sorted from highest to lowest forecast value.

    \item completing the companion Top-10 Project composition chart using the same
    governed Project dimension and the existing \texttt{Actual Net Sales} and
    \texttt{Project Accrual Forecast} measures;
    \item creating a dedicated \texttt{Executive Dashboard} report page with eight
    EUR-formatted forecast and CRM-governance KPI cards;
    \item adding a governed Business Unit slicer based on
    \texttt{DimBusinessUnit[BusinessUnitName]}; and
    \item adding the 12-month Actual-versus-Accrual forecast visual and identifying
    the remaining presentation correction: replace the current 100\% stacked column
    chart with a standard stacked column chart so the y-axis represents financial
    magnitude rather than percentage composition.
    \item completing that presentation correction by changing the monthly visual to
    a standard stacked column chart with a monetary EUR axis;
    \item interactively validating the Executive Dashboard with each governed
    Business Unit member and confirming consistent filtering across Actual Sales,
    Project Accrual and CRM governance measures; and
    \item confirming that the explicit \texttt{Unassigned / Missing BU} member
    isolates EUR 42,550.00 of Accrual value without fabricating Actual or CRM value.

    \item completing the Business Unit forecast-composition chart using the governed
    \texttt{DimBusinessUnit} dimension and the validated \texttt{Actual Net Sales}
    and \texttt{Project Accrual Forecast} measures;
    \item adding the \texttt{Forecast Governance -- Included vs At Risk} table with
    five explicit governed forecast buckets and their treatment descriptions; and
    \item bringing the Executive Dashboard to a near-final analytical composition
    containing annual KPIs, organizational composition, monthly timing, Project
    concentration and forecast-governance context.
    \item completing a final semantic-model relationship audit covering all nine
    relationships and confirming eight active relationships plus the intentionally
    inactive CRM program-start date role;
    \item completing an end-to-end Business Unit slicer validation across the newly
    added Governance, Business Unit and Top-10 Project visuals for all four governed
    Business Unit members; and
    \item hiding technical keys, sort helpers and selected data-quality diagnostics
    from the business-facing field list while preserving them in the model for audit
    and maintenance.
\end{itemize}

% ============================================================
% SOURCE DATA
% ============================================================

\section{2. Source Data}

The following source datasets were imported into Power BI:

\begin{tabularx}{\textwidth}{
    >{\bfseries\color{AaltoNavy}}p{4.7cm}
    X
}

\toprule

Dataset
&
Purpose
\\

\midrule

Net Sales Actuals
&
Actual net sales by Business Unit, project and accounting month.
\\

Accruals
&
Forecasted accrued net sales for ongoing projects.
\\

Open Custom Opportunities
&
Open CRM sales opportunities, including estimated value,
probability and expected annual allocation.
\\


\bottomrule

\end{tabularx}

\begin{inferencebox}

\textbf{Inference -- Source Scope}

The three datasets represent complementary forecast layers rather than
interchangeable inputs. Actual sales provide recognized revenue, accruals
provide the remaining forecast for ongoing projects, and CRM opportunities
provide probability-adjusted future potential. They must therefore be
validated independently before being combined in one forecast model.

\end{inferencebox}

% ============================================================
% QUERY ARCHITECTURE
% ============================================================

\section{3. Power Query Architecture}

A staging-and-fact architecture was implemented to separate source
ingestion from analytical transformation and validation.

\subsection{3.1 Staging Layer}

The raw staging queries are:

\begin{itemize}[leftmargin=7mm,itemsep=2pt]
    \item \texttt{stg\_Actuals\_Raw}
    \item \texttt{stg\_Accruals\_Raw}
    \item \texttt{stg\_Opportunities\_Raw}
\end{itemize}

Loading is disabled for these staging queries. They remain part of the
refresh process and preserve traceability to the original Excel sources.

\subsection{3.2 Analytical Fact Layer}

Reference queries were created from the staging layer:

\begin{itemize}[leftmargin=7mm,itemsep=2pt]
    \item \texttt{Fact\_ActualSales}
    \item \texttt{Fact\_ProjectAccrual}
    \item \texttt{Fact\_Opportunities}
\end{itemize}

For example, the actual-sales and accrual fact queries reference their
corresponding staging queries:

\begin{lstlisting}[style=powerquery]
= stg_Actuals_Raw
\end{lstlisting}

\begin{lstlisting}[style=powerquery]
= stg_Accruals_Raw
\end{lstlisting}

\subsection{3.3 Data-Quality Layer}

Dedicated control queries were created to keep validation logic separate
from the analytical facts:

\begin{itemize}[leftmargin=7mm,itemsep=2pt]
    \item \texttt{DQ\_Actual\_DuplicateKeys};
    \item \texttt{DQ\_Accrual\_SourceTotal};
    \item \texttt{DQ\_Accrual\_ExactDuplicates};
    \item \texttt{DQ\_Accrual\_RepeatedBusinessKeys};
    \item \texttt{DQ\_Opportunity\_WeightedValueMismatch};
    \item \texttt{DQ\_Opportunity\_CurrentYearWeightedStatus};
    \item \texttt{DQ\_Opportunity\_AllocationValidation};
    \item \texttt{DQ\_Opportunity\_AllocationStatusSummary};
    \item \texttt{DQ\_Opportunity\_AllocationExceptions};
    \item \texttt{DQ\_Opportunity\_MissingAllocationByYear};
    \item \texttt{DQ\_Opportunity\_2026MissingAllocationDetail};
    \item \texttt{DQ\_Opportunity\_2026MissingAllocationExposure};
    \item \texttt{DQ\_Opportunity\_2026CurrentYearAllocationRisk};
    \item \texttt{DQ\_Opportunity\_2026ForecastReadinessSummary};
    \item \texttt{DQ\_Opportunity\_2026ForecastReadinessStatusSummary};
    \item \texttt{DQ\_Opportunity\_2026ForecastBridge}; and
    \item \texttt{DQ\_BusinessUnit\_SourceValues}.
\end{itemize}

Loading has now been disabled for the data-quality queries before semantic-model
construction. The queries remain available in Power Query with all applied
steps intact, but they no longer appear as report-model tables.

This layered architecture provides:

\begin{itemize}[leftmargin=7mm,itemsep=2pt]
    \item transparent source-to-report lineage;
    \item separation between raw, analytical and control data;
    \item repeatable refresh operations;
    \item auditable transformation steps;
    \item easier future replacement of source systems; and
    \item cleaner tables in the semantic model.
\end{itemize}

\begin{inferencebox}

\textbf{Inference -- Query Architecture}

The separation of staging, analytical facts and data-quality controls creates
a traceable refresh path. Raw source evidence remains unchanged while business
transformations and exception logic can evolve without obscuring lineage.

\end{inferencebox}

\subsection{3.4 Dimension Layer}

A dedicated dimension group has now been added:

\begin{lstlisting}[style=powerquery]
04_Dimensions
\end{lstlisting}

The first shared dimension is:

\begin{lstlisting}[style=powerquery]
DimBusinessUnit
\end{lstlisting}

It contains three canonical Business Units and one explicit technical member
for records whose Business Unit is missing. The dimension is intentionally
separate from the facts so source labels can be preserved while relationships
use stable numeric keys.

A second shared dimension has now been created:

\begin{lstlisting}[style=powerquery]
DimDate
\end{lstlisting}

It covers every day from 1 January 2026 through 31 December 2026 and provides
daily, monthly and quarterly attributes for common time filtering. The query has
now been positioned in the governed dimension layer before Date relationships
were finalized.

A separate CRM calendar has subsequently been added:

\begin{lstlisting}[style=powerquery]
DimCRMDate
\end{lstlisting}

It covers the complete 2024--2028 opportunity horizon and supports explicit
role-playing analysis of CRM closing dates and simulated program-start dates.

A disconnected presentation dimension has also been added:

\begin{lstlisting}[style=powerquery]
DimForecastBucket
\end{lstlisting}

It is intentionally not related to the facts and is used only to present the
governed forecast bridge and scenario-treatment buckets.

The Project dimension has now been constructed in Power Query as:

\begin{lstlisting}[style=powerquery]
DimProject
\end{lstlisting}

Its profile contains 252 ProjectIDs, all distinct and unique, with zero
errors and zero empty values. The query has now been moved into
\texttt{04\_Dimensions} and retained as a loaded semantic-model dimension.
Although the table also contains a technical \texttt{ProjectKey}, the final
relationships use the clean natural key \texttt{ProjectID} directly because that
field is already present in both project-based facts and is unique on the
Dimension side.

\begin{inferencebox}

\textbf{Inference -- Dimension Layer}

Introducing a shared Business Unit dimension removes dependence on source-
specific text labels. Actual Sales can retain its shorter source vocabulary,
CRM and Accrual can retain their existing labels, and all three facts can still
filter consistently through the same governed Business Unit keys.

\end{inferencebox}

\subsection{3.5 Semantic-Model Load Scope and Business Unit Relationships}

Before relationship construction, model load was disabled for the staging and
data-quality queries. Their transformation and validation logic remains in
Power Query. At that stage the visible semantic model was reduced to
\texttt{DimBusinessUnit} plus the three analytical facts. After the shared Date
dimension was created and loaded, the current visible model contains:

\begin{itemize}[leftmargin=7mm,itemsep=2pt]
    \item \texttt{DimBusinessUnit};
    \item \texttt{DimDate};
    \item \texttt{Fact\_ActualSales};
    \item \texttt{Fact\_ProjectAccrual}; and
    \item \texttt{Fact\_Opportunities}.
\end{itemize}

Subsequent modelling phases added \texttt{DimCRMDate}, the disconnected
\texttt{DimForecastBucket} presentation table and \texttt{DimProject}. The current
core star-schema dimensions therefore include Business Unit, 2026 Date, CRM Date
roles and Project, while \texttt{DimForecastBucket} remains intentionally disconnected.

Three Business Unit relationships were then created:

\begin{tabularx}{\textwidth}{
    >{\ttfamily\raggedright\arraybackslash}p{5.2cm}
    >{\ttfamily\raggedright\arraybackslash}p{5.2cm}
    >{\centering\arraybackslash}p{1.8cm}
    >{\centering\arraybackslash}X
}

\toprule

Dimension Side
&
Fact Side
&
Cardinality
&
Filter
\\

\midrule

DimBusinessUnit[BusinessUnitKey]
&
Fact\_ActualSales[BusinessUnitKey]
&
1:*
&
Single
\\

DimBusinessUnit[BusinessUnitKey]
&
Fact\_ProjectAccrual[BusinessUnitKey]
&
1:*
&
Single
\\

DimBusinessUnit[BusinessUnitKey]
&
Fact\_Opportunities[BusinessUnitKey]
&
1:*
&
Single
\\

\bottomrule

\end{tabularx}

All three relationship lines are active. In the Power BI relationship editor,
the same structure may be displayed from the fact-table perspective as
\texttt{Many to one (*:1)}, with \texttt{DimBusinessUnit} on the one side.
The selected Actual Sales relationship is explicitly confirmed as active with
\texttt{Cross-filter direction = Single}.

\begin{inferencebox}

\textbf{Inference -- Business Unit Semantic Model}

The Business Unit portion of the model now follows a controlled star-schema
pattern: one governed dimension filters three separate facts, while fact-to-fact
relationships are avoided. Disabling load for staging and DQ queries removes
technical tables from the report model without deleting their audit logic.
Single-direction relationships reduce ambiguity and establish a clean basis for
shared Business Unit slicing across actuals, accruals and CRM opportunities.

\end{inferencebox}

\subsection{3.6 Shared 2026 Date Dimension}

A daily Date dimension was created to provide one governed calendar for
Actual Sales and Project Accrual reporting:

\begin{lstlisting}[style=powerquery]
let
    StartDate = #date(2026, 1, 1),
    EndDate = #date(2026, 12, 31),

    DateList =
        List.Dates(
            StartDate,
            Duration.Days(EndDate - StartDate) + 1,
            #duration(1, 0, 0, 0)
        ),

    ToTable =
        Table.FromList(
            DateList,
            Splitter.SplitByNothing(),
            {"Date"}
        ),

    ChangedType =
        Table.TransformColumnTypes(
            ToTable,
            {
                {"Date", type date}
            }
        ),

    AddedYear =
        Table.AddColumn(
            ChangedType,
            "Year",
            each Date.Year([Date]),
            Int64.Type
        ),

    AddedMonthNumber =
        Table.AddColumn(
            AddedYear,
            "MonthNumber",
            each Date.Month([Date]),
            Int64.Type
        ),

    AddedMonthName =
        Table.AddColumn(
            AddedMonthNumber,
            "MonthName",
            each Date.MonthName([Date]),
            type text
        ),

    AddedYearMonth =
        Table.AddColumn(
            AddedMonthName,
            "YearMonth",
            each Date.Year([Date]) * 100 + Date.Month([Date]),
            Int64.Type
        ),

    AddedMonthStart =
        Table.AddColumn(
            AddedYearMonth,
            "MonthStart",
            each Date.StartOfMonth([Date]),
            type date
        ),

    AddedQuarter =
        Table.AddColumn(
            AddedMonthStart,
            "Quarter",
            each "Q" & Text.From(Date.QuarterOfYear([Date])),
            type text
        )

in
    AddedQuarter
\end{lstlisting}

The completed profile confirms:

\begin{tabularx}{\textwidth}{
    >{\raggedright\arraybackslash\bfseries\color{AaltoNavy}}p{7.0cm}
    X
}

\toprule

Date-Dimension Metric
&
Result
\\

\midrule

Calendar period
&
1 January 2026 -- 31 December 2026
\\

Rows
&
365
\\

Date valid values
&
100\%
\\

Date errors
&
0
\\

Date empty values
&
0
\\

Distinct / unique Date values
&
365 / 365
\\

Distinct Year values
&
1
\\

Distinct MonthNumber / MonthName values
&
12 / 12
\\

Distinct YearMonth values
&
12
\\

Distinct MonthStart values
&
12
\\

Distinct Quarter values
&
4
\\

\bottomrule

\end{tabularx}

\begin{successbox}

\textbf{\color{SuccessGreen}{Date-Dimension Relationship Result}}

\texttt{DimDate} is now used as a governed semantic-model dimension. Two active
Date relationships connect the calendar to Actual Sales and Project Accrual:

\begin{itemize}[leftmargin=7mm,itemsep=2pt]
    \item \texttt{Fact\_ActualSales[MonthStart]} many-to-one
    \texttt{DimDate[Date]}; and
    \item \texttt{Fact\_ProjectAccrual[ForecastMonth]} many-to-one
    \texttt{DimDate[Date]}.
\end{itemize}

Both relationships use \texttt{Cross-filter direction = Single} and are active.

\end{successbox}

\begin{inferencebox}

\textbf{Inference -- Shared Date Dimension}

The 2026 calendar is complete, unique and error-free and now functions as the
shared time dimension for the two monthly financial facts. Actual Sales and
Project Accrual remain separate facts, but both can be filtered through one
calendar. CRM Opportunity dates are intentionally not connected at this stage
because the CRM table contains multiple date roles and records outside 2026;
forcing one active CRM Date relationship now would create an unnecessary
modelling constraint.

\end{inferencebox}

% ============================================================
% ACTUALS TRANSFORMATION
% ============================================================

\section{4. Actual Net Sales Transformation}

The \texttt{Fact\_ActualSales} query was prepared using a controlled
transformation sequence.

\subsection{4.1 Source Structure and Data Types}

\begin{tabularx}{\textwidth}{
    >{\bfseries}p{4.5cm}
    p{4cm}
    X
}

\toprule

Source Column
&
Data Type
&
Business Meaning
\\

\midrule

BU
&
Text
&
Source Business Unit label
\\

Cost center (prj)
&
Text
&
Project-related cost center
\\

Project
&
Text
&
Project identifier
\\

Year-Month
&
Whole Number
&
Accounting year and month
\\

Net sales
&
Decimal Number
&
Actual recognized net sales
\\

\bottomrule

\end{tabularx}

\subsection{4.2 Text Standardization}

Trim and Clean transformations were applied to the Business Unit, cost
center and project columns. This removes leading or trailing spaces and
non-printing characters without changing valid business values.

The corresponding Power Query steps were documented as:

\begin{lstlisting}[style=powerquery]
Trimmed Text
Cleaned Text
\end{lstlisting}

\subsection{4.3 MonthStart Field}

The numeric \texttt{Year-Month} field was converted into a valid
month-level date:

\begin{lstlisting}[style=powerquery]
let
    ym = Text.PadStart(Text.From([#"Year-Month"]), 6, "0")
in
    #date(
        Number.FromText(Text.Start(ym, 4)),
        Number.FromText(Text.End(ym, 2)),
        1
    )
\end{lstlisting}

The resulting column was named \texttt{MonthStart} and its data type was
set to \texttt{Date}. The actual-sales period covers January through
August 2026.

\begin{inferencebox}

\textbf{Inference -- Actual-Sales Transformation}

The actual-sales source has been converted into a consistent monthly fact
structure. The valid \texttt{MonthStart} field enables reliable time filtering,
date relationships and comparison with the September--December forecast
period.

\end{inferencebox}

% ============================================================
% ACTUAL QUALITY
% ============================================================

\section{5. Actual Sales Data-Quality Classification}

A row-level quality classification was implemented to distinguish
critical data problems from financially meaningful values requiring
review.

\begin{lstlisting}[style=powerquery]
if [BU] = null or Text.Trim(Text.From([BU])) = "" then
    "Critical - Missing BU"
else if [Project] = null
    or Text.Trim(Text.From([Project])) = "" then
    "Critical - Missing Project"
else if [MonthStart] = null then
    "Critical - Invalid Month"
else if [#"Net sales"] = null then
    "Critical - Missing Net Sales"
else if [#"Net sales"] < 0 then
    "Review - Negative Amount"
else if [#"Net sales"] = 0 then
    "Review - Zero Amount"
else
    "Valid"
\end{lstlisting}

The resulting column was named \texttt{ActualQualityStatus}.

\begin{warningbox}

Zero and negative net sales values were not automatically removed.

Negative values may represent credit notes, accounting corrections or
revenue reversals. These records were retained and classified as review
items for Finance validation.

\end{warningbox}

\begin{inferencebox}

\textbf{Inference -- Actual-Sales Quality}

The absence of critical issues means all 795 actual-sales records contain the
minimum identifiers and values required for analysis. The 118 zero or negative
records are accounting review items rather than technical failures and should
remain visible until their financial meaning is confirmed.

\end{inferencebox}

% ============================================================
% ACTUAL GRAIN VALIDATION
% ============================================================

\section{6. Actual Project-Month Grain Validation}

The intended analytical grain of the actual-sales fact is:

\begin{center}
\textbf{One row per project per accounting month}
\end{center}

A composite project-month key was created:

\begin{lstlisting}[style=powerquery]
if [Project] = null or [MonthStart] = null then
    null
else
    Text.Upper(Text.Trim(Text.From([Project])))
    & "|"
    & Date.ToText([MonthStart], "yyyyMM")
\end{lstlisting}

The resulting field was named \texttt{ProjectMonthKey}.

The dedicated duplicate-control query grouped the records by this field.
Only keys with a record count greater than one were retained. The query
was named \texttt{DQ\_Actual\_DuplicateKeys}.

The result was:

\begin{center}
\textbf{\color{SuccessGreen}{0 duplicate project-month keys}}
\end{center}

This confirms that the actual-sales table currently complies with the
intended grain.

\begin{inferencebox}

\textbf{Inference -- Actual-Sales Grain}

No project-month duplication was detected. Consequently, actual net sales can
be aggregated by project and month without a duplicate-driven overstatement at
the validated source grain.

\end{inferencebox}

% ============================================================
% ACTUAL COLUMNS AND RESULTS
% ============================================================

\section{7. Actual Sales Analytical Structure and Results}

\subsection{7.1 Standardized Column Names}

\begin{tabularx}{\textwidth}{
    >{\ttfamily}p{5cm}
    >{\ttfamily}X
}

\toprule

Source Name
&
Analytical Name
\\

\midrule

BU
&
BusinessUnitSource
\\

Cost center (prj)
&
CostCenter
\\

Project
&
ProjectID
\\

Year-Month
&
YearMonth
\\

Net sales
&
ActualNetSales
\\

\bottomrule

\end{tabularx}

The final \texttt{Fact\_ActualSales} structure now contains:

\begin{itemize}[leftmargin=7mm,itemsep=2pt]
    \item \texttt{BusinessUnitSource};
    \item \texttt{CostCenter};
    \item \texttt{ProjectID};
    \item \texttt{YearMonth};
    \item \texttt{ActualNetSales};
    \item \texttt{MonthStart};
    \item \texttt{ActualQualityStatus};
    \item \texttt{ProjectMonthKey}; and
    \item \texttt{BusinessUnitKey}.
\end{itemize}

\subsection{7.2 Validation Results}

\begin{tabularx}{\textwidth}{
    >{\bfseries\color{AaltoNavy}}p{6cm}
    X
}

\toprule

Validation Metric
&
Result
\\

\midrule

Total source records
&
795
\\

Valid records
&
677
\\

Zero-value review records
&
102
\\

Negative-value review records
&
16
\\

Critical data-quality issues
&
0
\\

Missing quality classifications
&
0
\\

Duplicate project-month keys
&
0
\\

Available accounting period
&
January--August 2026
\\

Total actual net sales
&
EUR 14,635,067.81
\\

\bottomrule

\end{tabularx}

\begin{successbox}

\textbf{\color{SuccessGreen}{Actual Sales Result}}

The actual net sales dataset is structurally valid, traceable and ready
for integration into the forecasting semantic model.

All 795 records were retained. Zero and negative accounting values remain
visible for Finance review.

\end{successbox}

\begin{inferencebox}

\textbf{Inference -- Actual-Sales Readiness}

The actual-sales fact is structurally ready for the semantic model. Its total
of EUR 14,635,067.81 is supported by complete row-level classifications and a
validated project-month key, providing a dependable year-to-date baseline for
the forecast.

\end{inferencebox}

\subsection{7.3 Canonical Business Unit Key Mapping}

After \texttt{DimBusinessUnit} was created, a numeric
\texttt{BusinessUnitKey} was added to \texttt{Fact\_ActualSales}. The original
\texttt{BusinessUnitSource} text field was retained for lineage.

The implemented mapping is:

\begin{lstlisting}[style=powerquery]
if [BusinessUnitSource] = "ExEd" then
    1
else if [BusinessUnitSource] = "Qualification Programs" then
    2
else if [BusinessUnitSource] = "University Programs" then
    3
else
    0
\end{lstlisting}

The new field was explicitly typed as Whole Number / \texttt{Int64.Type}. The
completed profile shows:

\begin{tabularx}{\textwidth}{
    >{\raggedright\arraybackslash\bfseries\color{AaltoNavy}}p{7cm}
    X
}

\toprule

Actual Sales BU-Key Metric
&
Result
\\

\midrule

Total Actual Sales rows
&
795
\\

BusinessUnitKey valid values
&
100\%
\\

BusinessUnitKey errors
&
0
\\

BusinessUnitKey empty values
&
0
\\

Distinct BusinessUnitKey values
&
3
\\

Key 1 -- ExEd Programs BU
&
346 rows
\\

Key 2 -- Qualification Programs BU
&
120 rows
\\

Key 3 -- University Programs BU
&
329 rows
\\

Key 0 -- Unassigned / Missing BU
&
0 rows
\\

\bottomrule

\end{tabularx}

\begin{inferencebox}

\textbf{Inference -- Actual Sales Business Unit Mapping}

All 795 Actual Sales records now have a stable numeric Business Unit key while
the source label remains available for traceability. The mapping reproduces the
earlier source-value audit exactly: 346 ExEd rows, 120 Qualification Programs
rows and 329 University Programs rows. No Actual Sales record falls into the
Unassigned member, so the first fact is ready for a one-to-many relationship
from \texttt{DimBusinessUnit}.

\end{inferencebox}

% ============================================================
% ACCRUAL PREPARATION
% ============================================================

\section{8. Ongoing-Project Accrual Transformation}

The second implementation phase prepared the ongoing-project accrual
forecast contained in \texttt{Fact\_ProjectAccrual}.

\subsection{8.1 Source Structure and Data Types}

\begin{tabularx}{\textwidth}{
    >{\bfseries}p{5cm}
    p{3.5cm}
    X
}

\toprule

Source Column
&
Data Type
&
Business Meaning
\\

\midrule

BU
&
Text
&
Source Business Unit label
\\

Accrual Basis
&
Text
&
Basis or rule used for the project accrual
\\

PNo.
&
Text
&
Project number
\\

Net Sales Year
&
Whole Number
&
Forecast financial year
\\

Net Sales Month
&
Whole Number
&
Forecast financial month
\\

Accrued Net Sales Forecast
&
Decimal Number
&
Forecasted net sales for the project-month
\\

\bottomrule

\end{tabularx}

\subsection{8.2 Text Standardization}

Trim and Clean transformations were applied to \texttt{BU},
\texttt{Accrual Basis} and \texttt{PNo.}. This prevents false differences
caused by whitespace or non-printing characters.

\subsection{8.3 ForecastMonth Field}

The forecast year and month were combined into a valid monthly date:

\begin{lstlisting}[style=powerquery]
try
    #date(
        Int64.From([#"Net Sales Year"]),
        Int64.From([#"Net Sales Month"]),
        1
    )
otherwise
    null
\end{lstlisting}

The resulting column was named \texttt{ForecastMonth} and its data type
was set to \texttt{Date}. The source forecast period covers September
through December 2026.

\begin{inferencebox}

\textbf{Inference -- Accrual Transformation}

The accrual source now has a reliable monthly date at the same calendar level
as actual sales. Its September--December coverage complements the
January--August actual-sales period and supports a full-year 2026 forecast when
the two fact layers are combined through a shared Date dimension.

\end{inferencebox}

% ============================================================
% ACCRUAL TOTAL CONTROL
% ============================================================

\section{9. Accrual Source Total Separation}

The accrual source contains a Total row. Including this row in the
analytical fact would double-count the forecast. A Boolean helper field
was therefore created:

\begin{lstlisting}[style=powerquery]
if [BU] = null then
    false
else
    Text.Upper(Text.Trim(Text.From([BU]))) = "TOTAL"
\end{lstlisting}

The resulting field was named \texttt{IsTotalRow}.

The source Total was retained in the dedicated control query:

\begin{lstlisting}[style=powerquery]
DQ_Accrual_SourceTotal
\end{lstlisting}

This control query contains only records where:

\begin{lstlisting}[style=powerquery]
IsTotalRow = true
\end{lstlisting}

The analytical fact was filtered using:

\begin{lstlisting}[style=powerquery]
IsTotalRow = false
\end{lstlisting}

The validation results are:

\begin{tabularx}{\textwidth}{
    >{\bfseries\color{AaltoNavy}}p{6cm}
    X
}

\toprule

Accrual Control Metric
&
Result
\\

\midrule

Original source records
&
314
\\

Source Total rows identified
&
1
\\

Detail records retained in fact
&
313
\\

Source control total
&
EUR 6,616,025.78
\\

Sum of all 313 detail records
&
EUR 6,616,025.78
\\

Reconciliation variance
&
EUR 0.00
\\

Classification errors
&
0
\\

Empty classifications
&
0
\\

Forecast period
&
September--December 2026
\\

\bottomrule

\end{tabularx}

\begin{successbox}

\textbf{\color{SuccessGreen}{Total-Row Control Result}}

The source Total has been separated from the analytical detail. It
remains available for reconciliation while only the 313 project-month
detail records are retained in \texttt{Fact\_ProjectAccrual}.

The helper field \texttt{IsTotalRow} can now be removed from the final
analytical fact structure.

\end{successbox}

\begin{inferencebox}

\textbf{Inference -- Source-Total Reconciliation}

The detail sum reconciles exactly to the source control total, with a variance
of EUR 0.00. This proves that filtering the Total row prevents double-counting
without losing financial value or breaking source-to-report reconciliation.

\end{inferencebox}

% ============================================================
% ACCRUAL QUALITY
% ============================================================

\section{10. Accrual Data-Quality Classification}

A row-level data-quality field was created in
\texttt{Fact\_ProjectAccrual}:

\begin{lstlisting}[style=powerquery]
if [BU] = null
    or Text.Trim(Text.From([BU])) = "" then
    "Critical - Missing BU"
else if [#"PNo."] = null
    or Text.Trim(Text.From([#"PNo."])) = "" then
    "Critical - Missing Project"
else if [ForecastMonth] = null then
    "Critical - Invalid Forecast Month"
else if [#"Accrued Net Sales Forecast"] = null then
    "Critical - Missing Forecast"
else if [#"Accrued Net Sales Forecast"] < 0 then
    "Review - Negative Forecast"
else if [#"Accrued Net Sales Forecast"] = 0 then
    "Review - Zero Forecast"
else
    "Valid"
\end{lstlisting}

The resulting column was named \texttt{AccrualQualityStatus}.

\begin{tabularx}{\textwidth}{
    >{\bfseries\color{AaltoNavy}}p{6cm}
    X
}

\toprule

Accrual Quality Status
&
Record Count
\\

\midrule

Valid
&
308
\\

Critical -- Missing BU
&
4
\\

Review -- Zero Forecast
&
1
\\

Critical -- Missing Project
&
0
\\

Critical -- Invalid Forecast Month
&
0
\\

Critical -- Missing Forecast
&
0
\\

Review -- Negative Forecast
&
0
\\

Total detail records
&
313
\\

\bottomrule

\end{tabularx}

\begin{warningbox}

Four accrual records have no Business Unit classification. Their combined
forecast value is EUR 42,550.00.

These records were not removed or assigned automatically. They remain
visible as controlled exceptions and require validation against a project
master or confirmation from Finance.

One zero-value forecast was also retained because it may represent a
valid project-month with no expected revenue.

\end{warningbox}

\begin{inferencebox}

\textbf{Inference -- Accrual Quality}

Most accrual records are analytically usable, but Business Unit completeness
is not yet perfect. The four unclassified records represent EUR 42,550.00 and
must remain visible as controlled exceptions until they can be matched to a
project master or confirmed by Finance.

\end{inferencebox}

% ============================================================
% ACCRUAL ANALYTICAL STRUCTURE
% ============================================================

\section{11. Accrual Analytical Structure}

The following standardized analytical names are used in the accrual fact
table:

\begin{tabularx}{\textwidth}{
    >{\ttfamily}p{5.6cm}
    >{\ttfamily}X
}

\toprule

Source Name
&
Analytical Name
\\

\midrule

BU
&
BusinessUnitSource
\\

Accrual Basis
&
AccrualBasis
\\

PNo.
&
ProjectID
\\

Net Sales Year
&
ForecastYear
\\

Net Sales Month
&
ForecastMonthNumber
\\

Accrued Net Sales Forecast
&
ProjectAccrualForecast
\\

ForecastMonth
&
ForecastMonth
\\

AccrualQualityStatus
&
AccrualQualityStatus
\\

\bottomrule

\end{tabularx}

After renaming the source fields and removing the temporary
\texttt{IsTotalRow} helper, the final \texttt{Fact\_ProjectAccrual}
structure now contains:

\begin{itemize}[leftmargin=7mm,itemsep=2pt]
    \item \texttt{BusinessUnitSource};
    \item \texttt{AccrualBasis};
    \item \texttt{ProjectID};
    \item \texttt{ForecastYear};
    \item \texttt{ForecastMonthNumber};
    \item \texttt{ProjectAccrualForecast};
    \item \texttt{ForecastMonth};
    \item \texttt{AccrualQualityStatus}; and
    \item \texttt{BusinessUnitKey}.
\end{itemize}

\begin{inferencebox}

\textbf{Inference -- Accrual Analytical Structure}

The standardized fields provide clear semantic names and a stable grain for
modelling. Source-specific labels remain available for lineage, while the
canonical Business Unit key supports controlled dimensional filtering.

\end{inferencebox}

\subsection{11.1 Canonical Business Unit Key Mapping}

A numeric \texttt{BusinessUnitKey} was added to
\texttt{Fact\_ProjectAccrual} using the canonical Business Unit vocabulary:

\begin{lstlisting}[style=powerquery]
if [BusinessUnitSource] = null
    or Text.Trim(Text.From([BusinessUnitSource])) = "" then
    0
else if [BusinessUnitSource] = "ExEd Programs BU" then
    1
else if [BusinessUnitSource] = "Qualification Programs BU" then
    2
else if [BusinessUnitSource] = "University Programs BU" then
    3
else
    null
\end{lstlisting}

The field was typed as Whole Number / \texttt{Int64.Type}. The completed
column profile contains 313 rows, 100\% valid values, 0 errors, 0 empty values
and 4 distinct keys.

The source-value audit reconciles the key distribution to:

\begin{itemize}[leftmargin=7mm,itemsep=2pt]
    \item key 1: 162 ExEd Programs BU records;
    \item key 2: 43 Qualification Programs BU records;
    \item key 3: 104 University Programs BU records; and
    \item key 0: 4 records whose Business Unit is missing.
\end{itemize}

\begin{inferencebox}

\textbf{Inference -- Project Accrual Business Unit Mapping}

All 313 accrual records can now participate in dimensional Business Unit
filtering without fabricating a classification for incomplete source rows.
The four missing-BU records are deliberately isolated under key 0, preserving
their EUR 42,550.00 forecast exposure as an explicit unassigned exception rather
than hiding it inside a real Business Unit.

\end{inferencebox}

% ============================================================
% ACCRUAL DUPLICATE CONTROL
% ============================================================

\section{12. Accrual Exact-Duplicate Control}

A dedicated data-quality query was created from
\texttt{Fact\_ProjectAccrual}:

\begin{lstlisting}[style=powerquery]
DQ_Accrual_ExactDuplicates
\end{lstlisting}

The query grouped the accrual records using all six fields that define an
exact duplicate:

\begin{itemize}[leftmargin=7mm,itemsep=2pt]
    \item \texttt{BusinessUnitSource};
    \item \texttt{AccrualBasis};
    \item \texttt{ProjectID};
    \item \texttt{ForecastYear};
    \item \texttt{ForecastMonthNumber}; and
    \item \texttt{ProjectAccrualForecast}.
\end{itemize}

The grouped records were counted in \texttt{RecordCount}. Only groups
with more than one source record were retained:

\begin{lstlisting}[style=powerquery]
= Table.SelectRows(
    #"Grouped Rows",
    each [RecordCount] > 1
)
\end{lstlisting}

Two additional diagnostic fields were created:

\begin{lstlisting}[style=powerquery]
DuplicateExcessCount = [RecordCount] - 1
\end{lstlisting}

\begin{lstlisting}[style=powerquery]
PotentialOverstatement =
    [DuplicateExcessCount] * [ProjectAccrualForecast]
\end{lstlisting}

\clearpage

\subsection{12.1 Duplicate-Control Results}

The completed control produced the following results:

\begin{tabularx}{\textwidth}{
    >{\raggedright\arraybackslash\bfseries\color{AaltoNavy}}p{6.3cm}
    X
}

\toprule

Duplicate-Control Metric
&
Result
\\

\midrule

Exact duplicate groups
&
8
\\

Source records involved
&
20
\\

Excess duplicate records
&
12
\\

Minimum records per duplicate group
&
2
\\

Maximum records per duplicate group
&
3
\\

Projects represented
&
2
\\

Business Unit represented
&
University Programs BU
\\

Accrual basis represented
&
Actual Close Date
\\

Affected forecast period
&
September--December 2026
\\

Potential duplicate overstatement
&
EUR 175,318.63
\\

\bottomrule

\end{tabularx}

Four duplicate groups relate to project \texttt{L9109051}. Each monthly
record appears three times. The other four groups relate to project
\texttt{L9109067}. The exact duplicated amount appears twice in each
affected month.

\begin{warningbox}

The potential overstatement of EUR 175,318.63 is a diagnostic amount. It
must not be treated as a confirmed correction before Finance validates
the business meaning of the repeated source records.

No records were deleted automatically from
\texttt{Fact\_ProjectAccrual}. The source-aligned amount remains available
for reconciliation and transparent review.

\end{warningbox}

\begin{inferencebox}

\textbf{Inference -- Exact-Duplicate Control}

The control identifies a concentrated duplication risk: 8 exact groups,
20 involved rows and 12 excess rows across only two projects. The EUR
175,318.63 potential overstatement is a risk estimate, not an approved
correction, because Finance must confirm the business meaning before any
source-aligned value is removed.

\end{inferencebox}

% ============================================================
% REPEATED BUSINESS-KEY ANALYSIS
% ============================================================

\section{13. Repeated Accrual Business-Key Analysis}

A second accrual control was implemented to distinguish repeated
project-month business keys from exact duplicate records. The dedicated
control query was named:

\begin{lstlisting}[style=powerquery]
DQ_Accrual_RepeatedBusinessKeys
\end{lstlisting}

\subsection{13.1 Candidate Business Key}

The accrual records were grouped using the following candidate business
key:

\begin{itemize}[leftmargin=7mm,itemsep=2pt]
    \item \texttt{BusinessUnitSource};
    \item \texttt{AccrualBasis};
    \item \texttt{ProjectID};
    \item \texttt{ForecastYear}; and
    \item \texttt{ForecastMonthNumber}.
\end{itemize}

The forecast amount was intentionally excluded from the grouping key.
This made it possible to determine whether repeated project-month records
contain the same forecast amount, different forecast components or a
combination of both patterns.

For every group, the query calculated:

\begin{itemize}[leftmargin=7mm,itemsep=2pt]
    \item \texttt{RecordCount}, representing the number of source rows;
    \item \texttt{TotalForecast}, representing the sum of the group;
    \item \texttt{AllRows}, retaining the underlying records for audit; and
    \item \texttt{DistinctAmountCount}, representing the number of distinct
    non-null forecast amounts.
\end{itemize}

Only groups satisfying the following condition were retained:

\begin{lstlisting}[style=powerquery]
RecordCount > 1
\end{lstlisting}

\subsection{13.2 Power Query Implementation}

The following implementation counts the distinct forecast amounts inside
the nested \texttt{AllRows} table of every repeated business key:

\begin{lstlisting}[style=powerquery]
let
    Source = Fact_ProjectAccrual,

    GroupedRows =
        Table.Group(
            Source,
            {
                "BusinessUnitSource",
                "AccrualBasis",
                "ProjectID",
                "ForecastYear",
                "ForecastMonthNumber"
            },
            {
                {
                    "RecordCount",
                    each Table.RowCount(_),
                    Int64.Type
                },
                {
                    "TotalForecast",
                    each List.Sum(
                        List.RemoveNulls(
                            [ProjectAccrualForecast]
                        )
                    ),
                    type nullable number
                },
                {
                    "AllRows",
                    each _,
                    type table
                }
            }
        ),

    FilteredRows =
        Table.SelectRows(
            GroupedRows,
            each [RecordCount] > 1
        ),

    AddedDistinctAmountCount =
        Table.AddColumn(
            FilteredRows,
            "DistinctAmountCount",
            each
                List.Count(
                    List.Distinct(
                        List.RemoveNulls(
                            Table.Column(
                                [AllRows],
                                "ProjectAccrualForecast"
                            )
                        )
                    )
                ),
            Int64.Type
        )

in
    AddedDistinctAmountCount
\end{lstlisting}

\subsection{13.3 Repeated Business-Key Results}

The completed control produced the following results:

\begin{tabularx}{\textwidth}{
    >{\raggedright\arraybackslash\bfseries\color{AaltoNavy}}p{7cm}
    X
}

\toprule

Repeated Business-Key Metric
&
Result
\\

\midrule

Repeated project-month business-key groups
&
10
\\

Source records represented
&
32
\\

Records above a one-row-per-key grain
&
22
\\

Groups containing one distinct amount
&
4
\\

Groups containing multiple distinct amounts
&
6
\\

Exact duplicate excess rows within the repeated keys
&
12
\\

Errors in distinct-amount calculation
&
0
\\

\bottomrule

\end{tabularx}

The value of 22 represents the number of records above a theoretical
one-row-per-business-key grain. It is not a confirmed duplicate count,
because some repeated keys contain different forecast amounts that may
represent legitimate forecast components.

The 12 exact duplicate excess rows can also be reconciled within this
control by summing:

\begin{lstlisting}[style=powerquery]
RecordCount - DistinctAmountCount
\end{lstlisting}

This result agrees with the independent exact-duplicate control described
in Section 12.

\clearpage

\subsection{13.4 Detailed Findings}

\begin{tabularx}{\textwidth}{
    >{\ttfamily}p{2.6cm}
    p{2.8cm}
    p{2.2cm}
    p{2.7cm}
    X
}

\toprule

ProjectID
&
Forecast Period
&
Record Count
&
Distinct Amounts
&
Interpretation
\\

\midrule

L9109051
&
September--December
&
3 per month
&
1 per month
&
The same forecast amount is repeated three times in each month.
\\

L9109067
&
September--December
&
4 per month
&
3 per month
&
Each group contains multiple forecast components and one repeated amount.
\\

O7401158
&
September
&
2
&
2
&
Two different forecast amounts are recorded for the same project-month.
\\

P8202082
&
September
&
2
&
2
&
Two different forecast amounts are recorded for the same project-month.
\\

\bottomrule

\end{tabularx}

\begin{warningbox}

A repeated business key does not automatically represent a duplicate.

A \texttt{DistinctAmountCount} of one indicates that the same amount is
repeated within the project-month group. A value greater than one indicates
that the group contains multiple forecast amounts.

The \texttt{L9109067} groups require particular attention because each group
contains four records and three distinct amounts. These groups contain both
multiple forecast components and an exact repeated amount.

No records were removed automatically. All repeated records remain available
for Finance validation and reconciliation with the original forecasting
logic.

\end{warningbox}

\subsection{13.5 Repeated-Key Status Classification}

An explicit status field was added to make the repeated-key findings easier
to review and communicate. The new field was named:

\begin{lstlisting}[style=powerquery]
RepeatedKeyStatus
\end{lstlisting}

The implemented classification logic was:

\begin{lstlisting}[style=powerquery]
if [DistinctAmountCount] = 1 then
    "Same Amount Repeated"
else if [RecordCount] > [DistinctAmountCount] then
    "Mixed - Components and Exact Duplicate"
else
    "Multiple Distinct Components"
\end{lstlisting}

The completed classification produced the following results:

\begin{tabularx}{\textwidth}{
    >{\raggedright\arraybackslash\bfseries\color{AaltoNavy}}X
    >{\centering\arraybackslash}p{2.2cm}
    >{\centering\arraybackslash}p{2.4cm}
    >{\centering\arraybackslash}p{2.5cm}
}

\toprule

Repeated-Key Status
&
Groups
&
Source Rows
&
Exact Excess Rows
\\

\midrule

Same Amount Repeated
&
4
&
12
&
8
\\

Mixed -- Components and Exact Duplicate
&
4
&
16
&
4
\\

Multiple Distinct Components
&
2
&
4
&
0
\\

\midrule

Total
&
10
&
32
&
12
\\

\bottomrule

\end{tabularx}

The four \texttt{Same Amount Repeated} groups relate to project
\texttt{L9109051}. Each month contains three identical amounts and therefore
two excess exact-repeat rows.

The four \texttt{Mixed} groups relate to project \texttt{L9109067}. Each month
contains four records but only three distinct amounts. Consequently, each
group contains multiple financial components and one exact-repeat row.

The remaining two groups, relating to projects \texttt{O7401158} and
\texttt{P8202082}, contain two records and two distinct amounts. They contain
no exact duplicate amount and are more likely to represent separate forecast
components, subject to business confirmation.

\begin{successbox}

\textbf{\color{SuccessGreen}{Repeated-Key Control Result}}

The exact-duplicate control identified 8 exact duplicate groups involving
20 source records and 12 excess duplicate records.

The broader business-key control identified 10 repeated project-month groups
involving 32 source records. Six groups contain multiple forecast amounts and
therefore may represent legitimate forecast components rather than simple
duplication.

The completed status classification divides these findings into 4 same-amount
groups, 4 mixed groups and 2 multiple-distinct-component groups. This confirms
that the 12 exact excess rows are concentrated in the first two categories.

\end{successbox}

\begin{inferencebox}

\textbf{Inference -- Repeated-Key Classification}

The repeated-key analysis proves that repeated project-month records do not
all have the same business meaning. Four groups are pure same-amount repeats,
four are mixed groups containing both distinct components and an exact repeat,
and two contain only distinct components. This evidence supports targeted
Finance review instead of deleting all 22 rows above a one-row-per-key grain.

\end{inferencebox}

% ============================================================
% CRM OPPORTUNITY SOURCE PROFILING
% ============================================================

\clearpage

\section{14. CRM Opportunity Preparation and Quality Controls}

The third source-preparation phase started with a full-column profile of the
open CRM opportunities in \texttt{Fact\_Opportunities}. The fact query
continues to reference the raw staging query:

\begin{lstlisting}[style=powerquery]
= stg_Opportunities_Raw
\end{lstlisting}

No business transformations, row deletions or automatic replacements were
applied during this profiling step.

\subsection{14.1 Profiling Configuration}

The following Power Query profiling controls were enabled:

\begin{itemize}[leftmargin=7mm,itemsep=2pt]
    \item Column quality;
    \item Column distribution; and
    \item Column profile.
\end{itemize}

Profiling was explicitly changed from the default first 1,000 rows to:

\begin{lstlisting}[style=powerquery]
Column profiling based on entire data set
\end{lstlisting}

The source contains 126 records and 14 columns.

\subsection{14.2 Observed CRM Source Structure}

\begin{tabularx}{\textwidth}{
    >{\ttfamily}p{5.8cm}
    p{3.4cm}
    X
}

\toprule

Source Column
&
Observed Type
&
Required Analytical Role
\\

\midrule

Status
&
Text
&
Opportunity status
\\

Est. Closing Date
&
Any / date-like
&
Estimated closing date
\\

Closing Date
&
Any / date-like
&
Current closing date
\\

Program start (simulated)
&
Any / date-like
&
Simulated program start date
\\

Interest / Opportunity
&
Text
&
Opportunity identifier or description
\\

Expected Accrued Net Sales for This Year \%
&
Whole Number
&
Current-year allocation percentage
\\

Expected Accrued Net Sales for Next Year \%
&
Whole Number
&
Next-year allocation percentage
\\

\% Probability
&
Percentage
&
Sales probability
\\

Value in EUR
&
Decimal Number
&
Source opportunity value
\\

Est. Value
&
Decimal Number
&
Estimated opportunity value
\\

Weighted Est. Value\_
&
Decimal Number
&
Probability-weighted value
\\

Weighted Est. Value this year
&
Decimal Number
&
Current-year weighted value
\\

Business Area
&
Text
&
Business-area classification
\\

P\&LS BU
&
Text
&
Business Unit classification
\\

\bottomrule

\end{tabularx}

\subsection{14.3 Profiling Results}

\begin{tabularx}{\textwidth}{
    >{\raggedright\arraybackslash\bfseries\color{AaltoNavy}}p{7cm}
    X
}

\toprule

CRM Profiling Metric
&
Observed Result
\\

\midrule

Total records
&
126
\\

Total source columns
&
14
\\

Profiling scope
&
Entire dataset
\\

Status
&
100\% valid, 0\% error, 0\% empty and 1 distinct value
\\

Status value represented
&
Open
\\

Interest / Opportunity
&
100\% valid, 126 distinct and 126 unique values
\\

Date-like source fields
&
100\% populated in the profile but not yet explicitly typed as Date
\\

Current-year allocation percentage
&
67\% valid and 33\% empty
\\

Next-year allocation percentage
&
68\% valid and 32\% empty
\\

Probability
&
98\% valid, 2\% empty and 11 distinct values
\\

Opportunity and weighted-value fields
&
100\% valid, 0\% error and 0\% empty
\\

Business Area
&
100\% valid and 1 distinct value
\\

P\&LS BU
&
100\% valid and 3 distinct values
\\

\bottomrule

\end{tabularx}

The preview includes closing and program-start dates extending into 2027 and
2028. Consequently, the financial-year allocation logic cannot be inferred
from the closing date alone and must be validated against the explicit
current-year and next-year percentage fields.

The source fields \texttt{Value in EUR} and \texttt{Est. Value} appear aligned
in the preview, but their equivalence must be tested across all 126 records
before either field is treated as redundant. Likewise, the weighted-value
fields must be reconciled against estimated value and probability rather than
accepted without calculation control.

\begin{inferencebox}

\textbf{Inference -- CRM Source Profiling}

The CRM source is structurally compact and contains no visible technical error
values, but it is not yet ready for the semantic model. Date-like fields require
explicit Date types, Probability contains missing values, and annual allocation
percentages are intentionally sparse. The unique opportunity descriptions
provide a candidate row identifier, while the three Business Unit values can
support cross-source BU standardization. Missing allocation values must be
classified through business rules rather than replaced with zero automatically.

\end{inferencebox}

\subsection{14.4 CRM Data-Type Standardization}

After the full-dataset profile was reviewed, explicit Power Query data types
were assigned to all 14 columns in \texttt{Fact\_Opportunities}. This step was
completed before text cleaning, analytical renaming or row-level quality
classification so that subsequent calculations operate on stable types.

\begin{tabularx}{\textwidth}{
    >{\ttfamily}p{7.2cm}
    X
}

\toprule

CRM Column
&
Assigned Data Type
\\

\midrule

Status
&
Text
\\

Est. Closing Date
&
Date
\\

Closing Date
&
Date
\\

Program start (simulated)
&
Date
\\

Interest / Opportunity
&
Text
\\

Expected Accrued Net Sales for This Year \%
&
Whole Number
\\

Expected Accrued Net Sales for Next Year \%
&
Whole Number
\\

\% Probability
&
Percentage
\\

Value in EUR
&
Decimal Number
\\

Est. Value
&
Decimal Number
\\

Weighted Est. Value\_
&
Decimal Number
\\

Weighted Est. Value this year
&
Decimal Number
\\

Business Area
&
Text
\\

P\&LS BU
&
Text
\\

\bottomrule

\end{tabularx}

\subsubsection*{Data-Type Validation Results}

\begin{tabularx}{\textwidth}{
    >{\raggedright\arraybackslash\bfseries\color{AaltoNavy}}p{7cm}
    X
}

\toprule

Data-Type Control Metric
&
Result
\\

\midrule

Source records retained
&
126
\\

Columns assigned explicit types
&
14
\\

Text columns
&
4
\\

Date columns
&
3
\\

Whole-number allocation columns
&
2
\\

Percentage columns
&
1
\\

Decimal-number value columns
&
4
\\

Type-conversion errors
&
0
\\

Rows removed or nulls replaced
&
None
\\

\bottomrule

\end{tabularx}

The data types were reviewed after conversion and no errors were observed.
The existing missing probability and annual-allocation values were deliberately
preserved as nulls. They were not converted to zero because a missing value is
not necessarily equivalent to a confirmed zero allocation or probability.

\begin{inferencebox}

\textbf{Inference -- CRM Data-Type Standardization}

All 126 CRM opportunity records can be converted to the required analytical
types without technical errors or record loss. The three date fields can now
support reliable period comparisons, while the probability and financial-value
fields can support controlled calculations. Because nulls were preserved, the
next quality rules can distinguish genuinely missing business information from
valid zero values rather than masking the difference during type conversion.

\end{inferencebox}

\subsection{14.5 CRM Text Standardization and Analytical Renaming}

The four CRM text columns were standardized after data-type validation. The
Power Query \texttt{Trim} transformation was applied first, followed by
\texttt{Clean}, to remove leading or trailing spaces and non-printing
characters:

\begin{itemize}[leftmargin=7mm,itemsep=2pt]
    \item \texttt{Status};
    \item \texttt{Interest / Opportunity};
    \item \texttt{Business Area}; and
    \item \texttt{P\&LS BU}.
\end{itemize}

The corresponding applied steps were retained as \texttt{Trimmed Text} and
\texttt{Cleaned Text}. No rows were filtered and no null values were replaced.

All 14 source fields were then renamed to clear analytical names:

\begin{tabularx}{\textwidth}{
    >{\ttfamily}p{7.1cm}
    >{\ttfamily}X
}

\toprule

Source Name
&
Analytical Name
\\

\midrule

Status
&
OpportunityStatus
\\

Est. Closing Date
&
EstimatedClosingDate
\\

Closing Date
&
ClosingDate
\\

Program start (simulated)
&
SimulatedProgramStartDate
\\

Interest / Opportunity
&
OpportunityName
\\

Expected Accrued Net Sales for This Year \%
&
CurrentYearAllocationPct
\\

Expected Accrued Net Sales for Next Year \%
&
NextYearAllocationPct
\\

\% Probability
&
Probability
\\

Value in EUR
&
OpportunityValue
\\

Est. Value
&
EstimatedValue
\\

Weighted Est. Value\_
&
WeightedEstimatedValue
\\

Weighted Est. Value this year
&
CurrentYearWeightedEstimatedValue
\\

Business Area
&
BusinessArea
\\

P\&LS BU
&
BusinessUnitSource
\\

\bottomrule

\end{tabularx}

The renaming was recorded in the \texttt{Renamed Columns} step. The operation
changed only metadata: all 126 rows and all 14 source fields remained present.

\begin{inferencebox}

\textbf{Inference -- CRM Text Standardization and Renaming}

Text cleaning reduces false category differences caused by invisible characters
or extra spaces. The analytical naming convention also aligns the CRM fact with
the existing facts, particularly through the shared
\texttt{BusinessUnitSource} field. The CRM source is therefore easier to audit,
join and explain without changing any source financial value.

\end{inferencebox}

\subsection{14.6 CRM Opportunity Data-Quality Classification}

A row-level quality field named \texttt{OpportunityQualityStatus} was added
after the analytical renaming step. The rule gives priority to missing core
identifiers and invalid ranges, followed by financial and completeness review
conditions:

\begin{lstlisting}[style=powerquery]
if [OpportunityStatus] = null
    or Text.Trim(Text.From([OpportunityStatus])) = "" then
    "Critical - Missing Status"

else if [OpportunityName] = null
    or Text.Trim(Text.From([OpportunityName])) = "" then
    "Critical - Missing Opportunity"

else if [BusinessUnitSource] = null
    or Text.Trim(Text.From([BusinessUnitSource])) = "" then
    "Critical - Missing BU"

else if [ClosingDate] = null then
    "Critical - Missing Closing Date"

else if [OpportunityValue] = null then
    "Critical - Missing Opportunity Value"

else if [EstimatedValue] = null then
    "Critical - Missing Estimated Value"

else if [Probability] <> null
    and ([Probability] < 0 or [Probability] > 1) then
    "Critical - Invalid Probability"

else if [CurrentYearAllocationPct] <> null
    and (
        [CurrentYearAllocationPct] < 0
        or [CurrentYearAllocationPct] > 100
    ) then
    "Critical - Invalid Current-Year Allocation"

else if [NextYearAllocationPct] <> null
    and (
        [NextYearAllocationPct] < 0
        or [NextYearAllocationPct] > 100
    ) then
    "Critical - Invalid Next-Year Allocation"

else if [CurrentYearAllocationPct] <> null
    and [NextYearAllocationPct] <> null
    and (
        [CurrentYearAllocationPct]
        + [NextYearAllocationPct]
    ) > 100 then
    "Critical - Allocation Above 100%"

else if [OpportunityValue] < 0
    or [EstimatedValue] < 0 then
    "Review - Negative Value"

else if [OpportunityValue] = 0
    or [EstimatedValue] = 0 then
    "Review - Zero Value"

else if [Probability] = null then
    "Review - Missing Probability"

else if [CurrentYearAllocationPct] = null
    and [NextYearAllocationPct] = null then
    "Review - Allocation Not Provided"

else
    "Valid"
\end{lstlisting}

The resulting field was explicitly typed as Text. Profiling was performed on
the entire dataset after the classification.

\begin{tabularx}{\textwidth}{
    >{\raggedright\arraybackslash\bfseries\color{AaltoNavy}}X
    >{\centering\arraybackslash}p{2.6cm}
    >{\centering\arraybackslash}p{2.6cm}
}

\toprule

Opportunity Quality Status
&
Record Count
&
Share of Records
\\

\midrule

Valid
&
95
&
75.4\%
\\

Review -- Allocation Not Provided
&
26
&
20.6\%
\\

Review -- Zero Value
&
2
&
1.6\%
\\

Review -- Missing Probability
&
2
&
1.6\%
\\

Critical -- Allocation Above 100\%
&
1
&
0.8\%
\\

\midrule

Total
&
126
&
100.0\%
\\

\bottomrule

\end{tabularx}

The completed classification contains five distinct statuses, zero errors,
zero empty classifications and one status occurring in only one record. No
record was deleted automatically. Since the rule returns one status per row,
the ordering of conditions deliberately exposes the highest-priority issue
first; secondary conditions can still be examined through dedicated control
queries.

\begin{warningbox}

The single \texttt{Critical - Allocation Above 100\%} record requires source
or Finance validation before it contributes to the forecast. The 26 records
without current-year or next-year allocations are review items rather than
confirmed errors, because opportunities extending beyond the immediate
forecast horizon may legitimately have no allocation yet.

\end{warningbox}

\begin{inferencebox}

\textbf{Inference -- CRM Opportunity Quality Classification}

The control confirms that 95 of 126 opportunities are valid under the current
rules and that every record receives an auditable classification. Most
exceptions relate to allocation completeness rather than structural source
failure. Only one record currently violates an explicit allocation rule, while
the remaining review records are retained for business confirmation. The CRM
fact is therefore technically controlled, but financial calculation checks are
still required before it is used in the final net-sales forecast.

\end{inferencebox}

\subsection{14.7 CRM Weighted-Value Reconciliation}

The source probability-weighted opportunity amount was independently
recalculated to test whether \texttt{WeightedEstimatedValue} agrees with
\texttt{EstimatedValue} multiplied by \texttt{Probability}. Three diagnostic
fields were added without overwriting the source values.

The independent weighted amount was calculated as:

\begin{lstlisting}[style=powerquery]
if [EstimatedValue] = null
    or [Probability] = null then
    null
else
    [EstimatedValue] * [Probability]
\end{lstlisting}

The resulting field was named
\texttt{CalculatedWeightedEstimatedValue} and typed as Decimal Number.

The source-to-calculated variance was then created:

\begin{lstlisting}[style=powerquery]
if [WeightedEstimatedValue] = null
    or [CalculatedWeightedEstimatedValue] = null then
    null
else
    [WeightedEstimatedValue]
        - [CalculatedWeightedEstimatedValue]
\end{lstlisting}

The resulting field was named \texttt{WeightedValueVariance} and typed as
Decimal Number. Finally, the tolerance-based reconciliation status was added:

\begin{lstlisting}[style=powerquery]
if [Probability] = null then
    "Not Testable - Missing Probability"
else if [WeightedEstimatedValue] = null then
    "Not Testable - Missing Weighted Value"
else if Number.Abs([WeightedValueVariance]) <= 0.01 then
    "Match"
else
    "Mismatch"
\end{lstlisting}

The resulting field was named \texttt{WeightedValueCheck} and typed as Text.
The EUR 0.01 tolerance prevents harmless binary floating-point residuals from
being reported as financial mismatches.

\begin{tabularx}{\textwidth}{
    >{\raggedright\arraybackslash\bfseries\color{AaltoNavy}}p{7cm}
    X
}

\toprule

Weighted-Value Control Metric
&
Result
\\

\midrule

Source records evaluated
&
126
\\

Columns after adding the three diagnostics
&
18
\\

WeightedValueCheck errors
&
0
\\

WeightedValueCheck empty values
&
0
\\

Distinct reconciliation statuses
&
3
\\

Match
&
123
\\

Not Testable -- Missing Probability
&
2
\\

Mismatch
&
1
\\

Exact zero variances
&
118
\\

Non-zero residuals accepted within EUR 0.01 tolerance
&
5
\\

Material mismatch variance
&
EUR -468,750.00
\\

\bottomrule

\end{tabularx}

The variance definition is source weighted value minus independently calculated
weighted value. Therefore, the EUR -468,750.00 result indicates that the source
weighted amount is EUR 468,750 lower than the probability-based calculation for
the affected opportunity. Two other records cannot be tested because their
probability is missing; this is a business-data limitation rather than a Power
Query error.

\begin{warningbox}

The single mismatch must not be corrected automatically. Its opportunity name,
Business Unit, estimated value, probability and source weighted value must first
be reviewed together. The difference could reflect a source formula issue, a
manual override or a business rule that is not represented by the simple
\texttt{EstimatedValue * Probability} calculation.

\end{warningbox}

\begin{inferencebox}

\textbf{Inference -- CRM Weighted-Value Reconciliation}

The control demonstrates strong calculation consistency: 123 of 126 records
match the independent calculation within one cent. Of these, 118 are exact zero
variances and five contain only negligible floating-point residuals. Two rows
remain explicitly not testable because Probability is missing, and one material
mismatch has been isolated for targeted investigation. This approach prevents a
single anomaly from being hidden inside an otherwise reliable aggregate.

\end{inferencebox}

\subsection{14.8 Dedicated CRM Weighted-Value Mismatch Control}

The earlier reconciliation identified one material
\texttt{WeightedValueCheck = "Mismatch"} record. The dedicated data-quality
reference query originally planned for this exception has now been created and is
retained in the \texttt{03\_Data\_Quality} group:

\begin{lstlisting}[style=powerquery]
DQ_Opportunity_WeightedValueMismatch
\end{lstlisting}

The query isolates the exception without changing
\texttt{Fact\_Opportunities}. The source-minus-calculated variance remains
EUR -468,750.00 as documented in Section 14.7. The record therefore remains a
business-validation item rather than an automatically corrected source value.

\begin{inferencebox}

\textbf{Inference -- Weighted-Value Exception Isolation}

Separating the material mismatch into a dedicated control query keeps the
analytical fact source-aligned while making the exception easy to review,
demonstrate and audit. The technical control is complete; the business meaning
of the mismatch still requires confirmation.

\end{inferencebox}

\subsection{14.9 Opportunity Value versus Estimated Value Reconciliation}

The next CRM control compared the source fields
\texttt{OpportunityValue} and \texttt{EstimatedValue}. A diagnostic variance
field named \texttt{OpportunityValueVariance} was added:

\begin{lstlisting}[style=powerquery]
if [OpportunityValue] = null
    or [EstimatedValue] = null then
    null
else
    [OpportunityValue] - [EstimatedValue]
\end{lstlisting}

The resulting field was typed as Decimal Number. A second field,
\texttt{OpportunityValueCheck}, applies the same EUR 0.01 financial tolerance
used in the weighted-value control:

\begin{lstlisting}[style=powerquery]
if [OpportunityValue] = null then
    "Not Testable - Missing Opportunity Value"
else if [EstimatedValue] = null then
    "Not Testable - Missing Estimated Value"
else if [OpportunityValueVariance] = null then
    "Not Testable - Missing Variance"
else if Number.Abs([OpportunityValueVariance]) <= 0.01 then
    "Match"
else
    "Mismatch"
\end{lstlisting}

The full-dataset profile shows that both diagnostic fields are 100\% valid with
zero errors and zero empty values. \texttt{OpportunityValueCheck} has one
distinct status and the value shown across the dataset is \texttt{Match}.
Consequently, all 126 CRM records pass this reconciliation within the one-cent
tolerance.

\begin{tabularx}{\textwidth}{
    >{\raggedright\arraybackslash\bfseries\color{AaltoNavy}}p{7cm}
    X
}

\toprule

Opportunity-Value Control Metric
&
Result
\\

\midrule

Source records evaluated
&
126
\\

OpportunityValueVariance valid values
&
100\%
\\

OpportunityValueVariance errors
&
0
\\

OpportunityValueVariance empty values
&
0
\\

OpportunityValueCheck valid values
&
100\%
\\

OpportunityValueCheck errors
&
0
\\

OpportunityValueCheck empty values
&
0
\\

Distinct OpportunityValueCheck statuses
&
1
\\

Observed reconciliation status
&
Match
\\

\bottomrule

\end{tabularx}

\begin{successbox}

\textbf{\color{SuccessGreen}{Opportunity-Value Reconciliation Result}}

The CRM source fields \texttt{OpportunityValue} and
\texttt{EstimatedValue} reconcile across all 126 records within the EUR 0.01
tolerance. The earlier profiling observation that the two fields appeared
aligned has therefore been converted into an explicit row-level control.

\end{successbox}

\begin{inferencebox}

\textbf{Inference -- Opportunity Value versus Estimated Value}

This result removes one potential ambiguity from the source structure. The two
fields currently carry equivalent analytical values, but both are retained at
this stage to preserve source lineage. Any later decision to hide or remove one
from the semantic model can now be based on evidence rather than assumption.

\end{inferencebox}

\subsection{14.10 Current-Year Weighted-Value Reconciliation}

A further control was added to validate the CRM source field
\texttt{CurrentYearWeightedEstimatedValue}. The purpose is to test whether the
current-year weighted amount is consistent with estimated value, probability
and the explicit current-year allocation percentage.

\subsubsection*{Independent Current-Year Weighted Calculation}

The independently calculated amount was added as
\texttt{CalculatedCurrentYearWeightedValue}:

\begin{lstlisting}[style=powerquery]
if [EstimatedValue] = null
    or [Probability] = null
    or [CurrentYearAllocationPct] = null then
    null
else
    [EstimatedValue]
        * [Probability]
        * ([CurrentYearAllocationPct] / 100)
\end{lstlisting}

The field was typed as Decimal Number. The full-dataset profile shows 67\%
valid values, 0\% errors and 33\% empty values. The empty values are preserved
because the calculation is intentionally not forced when a required business
input is missing.

\subsubsection*{Current-Year Variance}

The source-to-calculated variance was then created as
\texttt{CurrentYearWeightedVariance}:

\begin{lstlisting}[style=powerquery]
if [CurrentYearWeightedEstimatedValue] = null
    or [CalculatedCurrentYearWeightedValue] = null then
    null
else
    [CurrentYearWeightedEstimatedValue]
        - [CalculatedCurrentYearWeightedValue]
\end{lstlisting}

The variance was typed as Decimal Number. The current profile reports:

\begin{itemize}[leftmargin=7mm,itemsep=2pt]
    \item 67\% valid values;
    \item 0\% errors; and
    \item 33\% empty values.
\end{itemize}

The empty variance values are expected wherever the source or calculated
current-year amount cannot be evaluated.

\subsubsection*{Null-Safe Current-Year Reconciliation Status}

The final row-level control is named
\texttt{CurrentYearWeightedCheck}. The implemented logic is:

\begin{lstlisting}[style=powerquery]
if [Probability] = null then
    "Not Testable - Missing Probability"
else if [CurrentYearAllocationPct] = null then
    "Not Testable - Missing Current-Year Allocation"
else if [CurrentYearWeightedEstimatedValue] = null then
    "Not Testable - Missing Source Value"
else if [CalculatedCurrentYearWeightedValue] = null
    or [CurrentYearWeightedVariance] = null then
    "Not Testable - Missing Calculated Value"
else if Number.Abs([CurrentYearWeightedVariance]) <= 0.01 then
    "Match"
else
    "Mismatch"
\end{lstlisting}

The field was explicitly typed as Text.

\subsubsection*{Power Query Step-Order Correction}

The first version of \texttt{CurrentYearWeightedCheck} was created before
\texttt{CurrentYearWeightedVariance}. Because Power Query evaluates applied
steps sequentially, the check could not reliably reference a field that did not
yet exist at that point in the query.

The control was corrected without changing the source data:

\begin{enumerate}[leftmargin=9mm,itemsep=2pt]
    \item create and type \texttt{CalculatedCurrentYearWeightedValue};
    \item create and type \texttt{CurrentYearWeightedVariance};
    \item remove the earlier misplaced \texttt{CurrentYearWeightedCheck};
    \item recreate \texttt{CurrentYearWeightedCheck} after the variance exists;
    \item reorder the diagnostic columns for readability; and
    \item explicitly type the recreated check as Text.
\end{enumerate}

The applied-step sequence now includes the later
\texttt{Removed Columns}, \texttt{Added Custom9},
\texttt{Reordered Columns} and \texttt{Changed Type11} steps that implement
this correction.

\subsubsection*{Current-Year Control Results at the Present Stage}

The completed profile now shows:

\begin{tabularx}{\textwidth}{
    >{\raggedright\arraybackslash\bfseries\color{AaltoNavy}}p{7.2cm}
    X
}

\toprule

Current-Year Weighted Control Metric
&
Result
\\

\midrule

Source records retained
&
126
\\

Columns in Fact\_Opportunities after diagnostics
&
23
\\

CalculatedCurrentYearWeightedValue valid values
&
67\%
\\

CalculatedCurrentYearWeightedValue errors
&
0
\\

CalculatedCurrentYearWeightedValue empty values
&
33\%
\\

CurrentYearWeightedVariance valid values
&
67\%
\\

CurrentYearWeightedVariance errors
&
0
\\

CurrentYearWeightedVariance empty values
&
33\%
\\

CurrentYearWeightedCheck valid values
&
100\%
\\

CurrentYearWeightedCheck errors
&
0
\\

CurrentYearWeightedCheck empty values
&
0
\\

Distinct CurrentYearWeightedCheck statuses
&
3
\\

Exact status counts
&
85 Match / 39 Missing Current-Year Allocation / 2 Missing Probability
\\

\bottomrule

\end{tabularx}

The three statuses have now been quantified in a dedicated summary query.
Across all 126 CRM opportunity records, 85 are \texttt{Match}, 39 are
\texttt{Not Testable - Missing Current-Year Allocation}, and 2 are
\texttt{Not Testable - Missing Probability}. No \texttt{Mismatch} row appears
in the status summary. Therefore, all 85 records for which the required
current-year inputs are available reconcile within the EUR 0.01 tolerance.

\begin{successbox}

\textbf{\color{SuccessGreen}{Current-Year Weighted Control Result}}

The current-year reconciliation logic is now technically stable. The variance
contains no errors, and the final \texttt{CurrentYearWeightedCheck} is
100\% valid with zero errors and zero empty classifications across all
126 records. The completed status summary confirms 85 matching records,
39 records that cannot be tested because the current-year allocation is
missing, and 2 records that cannot be tested because Probability is missing.

\end{successbox}

\begin{inferencebox}

\textbf{Inference -- Current-Year Weighted Reconciliation}

The control distinguishes a true numerical mismatch from a record that cannot be
tested because a business input is missing. This is important for forecast
governance: null allocation or probability values remain visible as explicit
``Not Testable'' categories rather than being silently converted to zero.
The summary shows that 85 of 126 records (67.5\%) are fully testable and all
85 reconcile, while 41 records (32.5\%) remain not testable because one of the
required business inputs is missing. The next technical focus can therefore
move to broader allocation completeness and cross-year allocation validation.

\end{inferencebox}

\subsection{14.11 Current-Year Weighted Status Summary Query}

A dedicated reference query was created from \texttt{Fact\_Opportunities} to
summarize the final row-level status without changing the analytical fact. The
query is named:

\begin{lstlisting}[style=powerquery]
DQ_Opportunity_CurrentYearWeightedStatus
\end{lstlisting}

The implemented Power Query logic is:

\begin{lstlisting}[style=powerquery]
let
    Source = Fact_Opportunities,

    GroupedRows =
        Table.Group(
            Source,
            {"CurrentYearWeightedCheck"},
            {
                {
                    "RecordCount",
                    each Table.RowCount(_),
                    Int64.Type
                }
            }
        ),

    SortedRows =
        Table.Sort(
            GroupedRows,
            {
                {"RecordCount", Order.Descending}
            }
        )

in
    SortedRows
\end{lstlisting}

The completed query returns exactly three rows:

\begin{tabularx}{\textwidth}{
    >{\raggedright\arraybackslash\ttfamily}X
    >{\centering\arraybackslash}p{2.5cm}
    >{\centering\arraybackslash}p{2.5cm}
}

\toprule

CurrentYearWeightedCheck
&
Record Count
&
Share
\\

\midrule

Match
&
85
&
67.5\%
\\

Not Testable -- Missing Current-Year Allocation
&
39
&
31.0\%
\\

Not Testable -- Missing Probability
&
2
&
1.6\%
\\

\midrule

Total
&
126
&
100.0\%
\\

\bottomrule

\end{tabularx}

\begin{successbox}

\textbf{\color{SuccessGreen}{Current-Year Status Summary Result}}

The control reconciles exactly to the 126 source opportunity records. There are
85 fully testable records and every one of them is a \texttt{Match}. No
\texttt{Mismatch} status is present. The remaining 41 records are not
calculation failures: 39 are missing the current-year allocation and 2 are
missing Probability.

\end{successbox}

\begin{inferencebox}

\textbf{Inference -- Current-Year Testability}

The distinction between a failed reconciliation and an untestable record is now
quantified explicitly. The calculation itself is consistent for every record
that has the required inputs. The remaining issue is data completeness, not a
detected numerical discrepancy. This provides a clear business message for the
interview: the model does not convert missing allocation or probability into
zero and therefore does not create false forecast certainty.

\end{inferencebox}

\subsection{14.12 Current-Year and Next-Year Allocation Validation}

After the current-year weighted-value controls were completed, a dedicated
reference query was created to validate the relationship between
\texttt{CurrentYearAllocationPct} and \texttt{NextYearAllocationPct}. The query
was named:

\begin{lstlisting}[style=powerquery]
DQ_Opportunity_AllocationValidation
\end{lstlisting}

The purpose of the control is to distinguish valid two-year allocations from
missing, incomplete or logically inconsistent allocation patterns without
modifying \texttt{Fact\_Opportunities}.

\subsubsection*{Allocation Total Calculation}

A diagnostic field named \texttt{AllocationTotalPct} was added. The total is
calculated only when both annual allocation values are present:

\begin{lstlisting}[style=powerquery]
if [CurrentYearAllocationPct] = null
    or [NextYearAllocationPct] = null then
    null
else
    [CurrentYearAllocationPct]
        + [NextYearAllocationPct]
\end{lstlisting}

The completed profile shows that \texttt{AllocationTotalPct} is 59\% populated,
41\% empty and contains 0 errors. This is consistent with 74 records having
both annual allocation values available and 52 records missing at least one of
the two inputs.

\subsubsection*{Allocation Validation Status}

A second diagnostic field named \texttt{AllocationValidationStatus} was added
using the following null-safe classification:

\begin{lstlisting}[style=powerquery]
if [CurrentYearAllocationPct] = null
    and [NextYearAllocationPct] = null then
    "Review - Both Allocations Missing"

else if [CurrentYearAllocationPct] = null then
    "Review - Missing Current-Year Allocation"

else if [NextYearAllocationPct] = null then
    "Review - Missing Next-Year Allocation"

else if [CurrentYearAllocationPct] < 0
    or [CurrentYearAllocationPct] > 100
    or [NextYearAllocationPct] < 0
    or [NextYearAllocationPct] > 100 then
    "Critical - Allocation Out of Range"

else if [AllocationTotalPct] > 100 then
    "Critical - Allocation Above 100%"

else if [AllocationTotalPct] = 100 then
    "Valid - Allocation Total 100%"

else
    "Review - Allocation Below 100%"
\end{lstlisting}

The final diagnostic query retains the following review fields:

\begin{itemize}[leftmargin=7mm,itemsep=2pt]
    \item \texttt{OpportunityName};
    \item \texttt{BusinessUnitSource};
    \item \texttt{ClosingDate};
    \item \texttt{SimulatedProgramStartDate};
    \item \texttt{CurrentYearAllocationPct};
    \item \texttt{NextYearAllocationPct};
    \item \texttt{AllocationTotalPct};
    \item \texttt{AllocationValidationStatus}; and
    \item \texttt{OpportunityQualityStatus}.
\end{itemize}

\subsubsection*{Allocation Validation Results}

The full-dataset profile contains 126 records, 0 errors, 0 empty
\texttt{AllocationValidationStatus} values and 6 observed statuses. The
distribution shown by the completed value profile corresponds to:

\begin{tabularx}{\textwidth}{
    >{\raggedright\arraybackslash}X
    >{\centering\arraybackslash}p{2.4cm}
    >{\centering\arraybackslash}p{2.4cm}
}

\toprule

Allocation Validation Status
&
Record Count
&
Share
\\

\midrule

Valid -- Allocation Total 100\%
&
70
&
55.6\%
\\

Review -- Both Allocations Missing
&
29
&
23.0\%
\\

Review -- Missing Current-Year Allocation
&
12
&
9.5\%
\\

Review -- Missing Next-Year Allocation
&
11
&
8.7\%
\\

Review -- Allocation Below 100\%
&
3
&
2.4\%
\\

Critical -- Allocation Above 100\%
&
1
&
0.8\%
\\

\midrule

Total
&
126
&
100.0\%
\\

\bottomrule

\end{tabularx}

The profile is internally consistent with the individual allocation-column
completeness results. The 29 records with both allocations missing plus the
12 records missing only the current-year allocation produce 41 missing
current-year values. Similarly, the same 29 records plus the 11 records missing
only the next-year allocation produce 40 missing next-year values.

Among the 74 records where both allocation percentages are present, 70 sum to
exactly 100\%, 3 sum to less than 100\%, and 1 sums to more than 100\%. No
\texttt{Critical - Allocation Out of Range} status is observed in the completed
profile.

\begin{successbox}

\textbf{\color{SuccessGreen}{Allocation Validation Result}}

The allocation-validation query classifies every CRM opportunity without
changing the source fact. Seventy records have a complete two-year allocation
that totals exactly 100\%. Fifty-two records are missing one or both annual
allocation inputs. Four additional records have both inputs present but do not
sum to 100\%: three are below 100\% and one is above 100\%.

\end{successbox}

\begin{inferencebox}

\textbf{Inference -- Cross-Year Allocation Validation}

The main allocation issue is completeness rather than widespread arithmetic
inconsistency. Most records with both annual allocation values present reconcile
to exactly 100\%. The single allocation-above-100\% record is a clear business
exception, while the three below-100\% records require interpretation because
their remaining share may belong to a later financial year. Missing allocation
values are therefore preserved as missing rather than being replaced by zero.
This keeps the forecast transparent and prevents false certainty.

\end{inferencebox}

\subsection{14.13 Allocation Status Summary Query}

The visual distribution from the allocation-validation query was converted into
an explicit grouped control table using a new reference query:

\begin{lstlisting}[style=powerquery]
DQ_Opportunity_AllocationStatusSummary
\end{lstlisting}

The implemented Power Query logic is:

\begin{lstlisting}[style=powerquery]
let
    Source = DQ_Opportunity_AllocationValidation,

    GroupedRows =
        Table.Group(
            Source,
            {"AllocationValidationStatus"},
            {
                {
                    "RecordCount",
                    each Table.RowCount(_),
                    Int64.Type
                }
            }
        ),

    SortedRows =
        Table.Sort(
            GroupedRows,
            {
                {"RecordCount", Order.Descending}
            }
        )

in
    SortedRows
\end{lstlisting}

The grouped output confirms the following exact distribution:

\begin{tabularx}{\textwidth}{
    >{\raggedright\arraybackslash}X
    >{\centering\arraybackslash}p{2.5cm}
}

\toprule

Allocation Validation Status
&
Record Count
\\

\midrule

Valid -- Allocation Total 100\%
&
70
\\

Review -- Both Allocations Missing
&
29
\\

Review -- Missing Current-Year Allocation
&
12
\\

Review -- Missing Next-Year Allocation
&
11
\\

Review -- Allocation Below 100\%
&
3
\\

Critical -- Allocation Above 100\%
&
1
\\

\midrule

Total
&
126
\\

\bottomrule

\end{tabularx}

\begin{inferencebox}

\textbf{Inference -- Allocation Status Summary}

The grouped control independently reproduces the allocation distribution and
reconciles exactly to all 126 CRM opportunities. This converts a visual profile
into an auditable control table that can later be reused directly in a
data-quality dashboard. No allocation record is lost between the source fact,
the row-level validation query and the grouped summary.

\end{inferencebox}

\subsection{14.14 Detailed Allocation Exceptions}

A dedicated exception query was created from
\texttt{DQ\_Opportunity\_AllocationValidation} to isolate only the four records
whose complete two-year allocation does not total 100\%:

\begin{lstlisting}[style=powerquery]
DQ_Opportunity_AllocationExceptions
\end{lstlisting}

The query retains only:

\begin{lstlisting}[style=powerquery]
AllocationValidationStatus =
    "Review - Allocation Below 100%"
or
AllocationValidationStatus =
    "Critical - Allocation Above 100%"
\end{lstlisting}

The resulting four records are:

\begin{tabularx}{\textwidth}{
    >{\ttfamily\raggedright\arraybackslash}p{3.1cm}
    >{\raggedright\arraybackslash}p{2.7cm}
    >{\centering\arraybackslash}p{1.8cm}
    >{\centering\arraybackslash}p{1.8cm}
    >{\centering\arraybackslash}p{1.4cm}
    >{\centering\arraybackslash}p{1.4cm}
    >{\centering\arraybackslash}p{1.4cm}
}

\toprule

Opportunity
&
Business Unit
&
Closing Date
&
Program Start
&
CY \%
&
NY \%
&
Total \%
\\

\midrule

Open custom opportunity 111
&
University Programs BU
&
2026-08-31
&
2026-11-29
&
30
&
80
&
110
\\

Open custom opportunity 86
&
Qualification Programs BU
&
2026-09-30
&
2026-12-29
&
10
&
40
&
50
\\

Open custom opportunity 3
&
ExEd Programs BU
&
2028-09-30
&
2028-12-29
&
25
&
25
&
50
\\

Open custom opportunity 33
&
ExEd Programs BU
&
2026-12-31
&
2027-03-31
&
0
&
30
&
30
\\

\bottomrule

\end{tabularx}

The first record is the only critical over-allocation: its current-year and
next-year shares sum to 110\%. Its existing
\texttt{OpportunityQualityStatus} is also
\texttt{Critical - Allocation Above 100\%}.

The remaining three records are below 100\%. Their existing quality statuses
show that two are otherwise \texttt{Valid}, while
\texttt{Open custom opportunity 3} is already classified as
\texttt{Review - Zero Value}. The date pattern is not uniform: one below-100\%
record closes and starts in 2028, one closes in late 2026 and starts in early
2027, and one both closes and starts in 2026.

\begin{warningbox}

The three below-100\% cases must not be automatically increased to 100\%.
A remaining share could reflect revenue allocated to a later financial year,
incomplete source planning, or another business rule not represented by the two
available annual fields. Likewise, the 110\% record must be validated with the
source owner or Finance before any correction is made.

\end{warningbox}

\begin{inferencebox}

\textbf{Inference -- Allocation Exception Review}

The four non-100\% cases are heterogeneous rather than one repeated technical
error. The 110\% record is a clear over-allocation exception. The three
below-100\% records require business interpretation because their closing and
program-start dates span different horizons. In particular, the 2028-dated
record supports the possibility that part of the revenue belongs beyond the
current and next financial years, whereas the 2026-dated 50\% case cannot be
explained by date horizon alone. The correct action is targeted review, not
blanket normalization.

\end{inferencebox}

\subsection{14.15 Missing Allocation Horizon Analysis}

The 52 opportunities with one or both annual allocation fields missing were
grouped by allocation status, closing year and simulated program-start year in
a dedicated reference query:

\begin{lstlisting}[style=powerquery]
DQ_Opportunity_MissingAllocationByYear
\end{lstlisting}

The query preserves the three missing-allocation categories and adds
\texttt{ClosingYear} and \texttt{ProgramStartYear} before grouping.

The completed result is:

\begin{tabularx}{\textwidth}{
    >{\raggedright\arraybackslash}p{5.7cm}
    >{\centering\arraybackslash}p{2cm}
    >{\centering\arraybackslash}p{2.3cm}
    >{\centering\arraybackslash}p{2cm}
}

\toprule

Allocation Status
&
Closing Year
&
Program Start Year
&
Record Count
\\

\midrule

Review -- Both Allocations Missing
&
2024
&
2024
&
1
\\

Review -- Both Allocations Missing
&
2025
&
2025
&
1
\\

Review -- Both Allocations Missing
&
2026
&
2026
&
15
\\

Review -- Both Allocations Missing
&
2026
&
2027
&
6
\\

Review -- Both Allocations Missing
&
2027
&
2027
&
4
\\

Review -- Both Allocations Missing
&
2027
&
2028
&
1
\\

Review -- Both Allocations Missing
&
2028
&
2028
&
1
\\

Review -- Missing Current-Year Allocation
&
2026
&
2026
&
5
\\

Review -- Missing Current-Year Allocation
&
2026
&
2027
&
4
\\

Review -- Missing Current-Year Allocation
&
2027
&
2027
&
3
\\

Review -- Missing Next-Year Allocation
&
2026
&
2026
&
8
\\

Review -- Missing Next-Year Allocation
&
2026
&
2027
&
3
\\

\midrule

Total
&
&
&
52
\\

\bottomrule

\end{tabularx}

The closing-year distribution is:

\begin{itemize}[leftmargin=7mm,itemsep=2pt]
    \item 2024: 1 record;
    \item 2025: 1 record;
    \item 2026: 41 records;
    \item 2027: 8 records; and
    \item 2028: 1 record.
\end{itemize}

Therefore, 41 of 52 missing-allocation opportunities, or 78.8\%, close in 2026.
The program-start-year distribution is 1 record in 2024, 1 in 2025,
28 in 2026, 20 in 2027 and 2 in 2028.

\begin{inferencebox}

\textbf{Inference -- Missing Allocation Horizon}

Missing annual allocation is not mainly a long-horizon phenomenon. Nearly four
out of five affected opportunities close in 2026, and 28 of the 52 also start
their simulated program in 2026. The issue therefore includes a material
near-term data-completeness component and cannot be explained simply by the
presence of future-dated opportunities. This makes the 2026 records the highest
priority for targeted review before forecast aggregation.

\end{inferencebox}

\subsection{14.16 Detailed 2026 Missing-Allocation Review}

A second reference query was created to isolate the near-term opportunities
identified in the horizon analysis:

\begin{lstlisting}[style=powerquery]
DQ_Opportunity_2026MissingAllocationDetail
\end{lstlisting}

The query retains records where:

\begin{itemize}[leftmargin=7mm,itemsep=2pt]
    \item \texttt{ClosingDate} is in 2026; and
    \item one or both annual allocation inputs are missing.
\end{itemize}

The completed query contains exactly 41 rows. The status composition is:

\begin{tabularx}{\textwidth}{
    >{\raggedright\arraybackslash}X
    >{\centering\arraybackslash}p{2.5cm}
    >{\centering\arraybackslash}p{2.5cm}
}

\toprule

2026 Missing-Allocation Status
&
Record Count
&
Share of 41
\\

\midrule

Both Allocations Missing
&
21
&
51.2\%
\\

Missing Current-Year Allocation
&
9
&
22.0\%
\\

Missing Next-Year Allocation
&
11
&
26.8\%
\\

\midrule

Total
&
41
&
100.0\%
\\

\bottomrule

\end{tabularx}

Within these 41 near-term records, 28 have a simulated program-start year of
2026 and 13 have a simulated program-start year of 2027. The detail query keeps
the opportunity name, Business Unit, closing date, simulated program-start date,
both allocation fields, the allocation total and the allocation-validation
status visible for direct review.

\begin{inferencebox}

\textbf{Inference -- 2026 Missing-Allocation Detail}

The 2026 detail confirms that the allocation-completeness issue is operationally
relevant to the current forecast horizon. More than half of the 41 near-term
records are missing both allocation fields, while the remaining records are
split between missing current-year and missing next-year allocation. Because
these are not merely distant opportunities, the next useful control is to
measure their financial exposure rather than only their record count.

\end{inferencebox}

\subsection{14.17 Financial Exposure of 2026 Missing Allocations}

The 41 near-term opportunities identified in the 2026 missing-allocation detail
were aggregated by missing-allocation status and Business Unit in a dedicated
control query:

\begin{lstlisting}[style=powerquery]
DQ_Opportunity_2026MissingAllocationExposure
\end{lstlisting}

The completed result is:

\begin{tabularx}{\textwidth}{
    >{\raggedright\arraybackslash}p{4.2cm}
    >{\raggedright\arraybackslash}p{3.1cm}
    >{\centering\arraybackslash}p{1.4cm}
    >{\raggedleft\arraybackslash}p{2.6cm}
    >{\raggedleft\arraybackslash}p{2.6cm}
}

\toprule

Missing Allocation Status
&
Business Unit
&
Records
&
Estimated Value
&
Weighted Value
\\

\midrule

Missing Current-Year Allocation
&
ExEd Programs BU
&
9
&
EUR 1,335,562.50
&
EUR 1,034,593.75
\\

Both Allocations Missing
&
Qualification Programs BU
&
3
&
EUR 1,312,500.00
&
EUR 615,625.00
\\

Both Allocations Missing
&
ExEd Programs BU
&
9
&
EUR 553,925.00
&
EUR 195,020.00
\\

Missing Next-Year Allocation
&
ExEd Programs BU
&
6
&
EUR 348,306.25
&
EUR 171,812.50
\\

Missing Next-Year Allocation
&
University Programs BU
&
5
&
EUR 288,025.00
&
EUR 156,991.25
\\

Both Allocations Missing
&
University Programs BU
&
9
&
EUR 370,380.48
&
EUR 143,940.24
\\

\midrule

Total
&
&
41
&
EUR 4,208,699.23
&
EUR 2,317,982.74
\\

\bottomrule

\end{tabularx}

By missing-allocation status, the same 41 records reconcile to:

\begin{itemize}[leftmargin=7mm,itemsep=2pt]
    \item 21 records with both allocations missing:
    EUR 2,236,805.48 estimated value and EUR 954,585.24 weighted value;
    \item 9 records missing only current-year allocation:
    EUR 1,335,562.50 estimated value and EUR 1,034,593.75 weighted value; and
    \item 11 records missing only next-year allocation:
    EUR 636,331.25 estimated value and EUR 328,803.75 weighted value.
\end{itemize}

\begin{inferencebox}

\textbf{Inference -- 2026 Missing-Allocation Financial Exposure}

The allocation-completeness issue is financially material, not merely a record-
count issue. The 41 near-term records represent more than EUR 4.2 million of
gross estimated pipeline and approximately EUR 2.32 million of probability-
weighted pipeline. The missing current-year allocation category is especially
important because only 9 records account for more than EUR 1.03 million of
weighted value. This supports prioritizing allocation completion by financial
exposure rather than treating all missing records equally.

\end{inferencebox}

\subsection{14.18 2026 Current-Year Allocation Risk by Business Unit}

A further control was created to focus only on 2026-closing opportunities where
\texttt{CurrentYearAllocationPct} is missing:

\begin{lstlisting}[style=powerquery]
DQ_Opportunity_2026CurrentYearAllocationRisk
\end{lstlisting}

The completed Business Unit summary is:

\begin{tabularx}{\textwidth}{
    >{\raggedright\arraybackslash}X
    >{\centering\arraybackslash}p{2.1cm}
    >{\raggedleft\arraybackslash}p{3.1cm}
    >{\raggedleft\arraybackslash}p{3.1cm}
}

\toprule

Business Unit
&
Records
&
Estimated Value
&
Weighted Value
\\

\midrule

ExEd Programs BU
&
18
&
EUR 1,889,487.50
&
EUR 1,229,613.75
\\

Qualification Programs BU
&
3
&
EUR 1,312,500.00
&
EUR 615,625.00
\\

University Programs BU
&
9
&
EUR 370,380.48
&
EUR 143,940.24
\\

\midrule

Total
&
30
&
EUR 3,572,367.98
&
EUR 1,989,178.99
\\

\bottomrule

\end{tabularx}

The 30 records with missing current-year allocation account for 85.8\% of the
EUR 2,317,982.74 weighted exposure identified across all 41 2026 missing-
allocation records. ExEd Programs BU alone represents approximately 61.8\% of
the current-year weighted risk, followed by Qualification Programs BU at
approximately 30.9\% and University Programs BU at approximately 7.2\%.

\begin{inferencebox}

\textbf{Inference -- Current-Year Allocation Risk}

The largest unresolved forecast risk is concentrated rather than evenly
distributed. Thirty opportunities cannot be reliably allocated into the 2026
forecast because their current-year allocation is missing, and they represent
approximately EUR 1.99 million of weighted pipeline. ExEd Programs BU is the
highest-priority review area. The appropriate next step is therefore to classify
2026 opportunities by forecast readiness and keep unresolved weighted pipeline
visible as at-risk rather than silently treating it as zero or fully included.

\end{inferencebox}

\subsection{14.19 2026 CRM Forecast Readiness by Business Unit}

All CRM opportunities with \texttt{ClosingDate} in 2026 were classified according
to whether their current-year contribution is ready for use, requires business
review or cannot yet be calculated reliably. The dedicated control query is:

\begin{lstlisting}[style=powerquery]
DQ_Opportunity_2026ForecastReadinessSummary
\end{lstlisting}

The rule applies the following priority:

\begin{enumerate}[leftmargin=9mm,itemsep=2pt]
    \item missing Probability;
    \item missing current-year allocation;
    \item cross-year allocation above 100\%;
    \item reconciled current-year weighted value; and
    \item any remaining current-year reconciliation review.
\end{enumerate}

The completed Business Unit result is:

\begin{tabularx}{\textwidth}{
    >{\raggedright\arraybackslash}p{4.8cm}
    >{\raggedright\arraybackslash}p{2.7cm}
    >{\centering\arraybackslash}p{1.2cm}
    >{\raggedleft\arraybackslash}p{2.4cm}
    >{\raggedleft\arraybackslash}p{2.4cm}
}

\toprule

Forecast Readiness Status
&
Business Unit
&
Rows
&
Weighted Value
&
Current-Year Value
\\

\midrule

Not Ready -- Missing Current-Year Allocation
&
ExEd Programs BU
&
18
&
EUR 1,229,613.75
&
--
\\

Not Ready -- Missing Current-Year Allocation
&
Qualification Programs BU
&
3
&
EUR 615,625.00
&
--
\\

Not Ready -- Missing Current-Year Allocation
&
University Programs BU
&
7
&
EUR 143,940.24
&
--
\\

Not Ready -- Missing Probability
&
University Programs BU
&
2
&
EUR 0.00
&
--
\\

Ready -- Reconciled Current-Year Value
&
ExEd Programs BU
&
60
&
EUR 1,779,552.50
&
EUR 527,877.50
\\

Ready -- Reconciled Current-Year Value
&
University Programs BU
&
13
&
EUR 365,763.75
&
EUR 244,772.50
\\

Ready -- Reconciled Current-Year Value
&
Qualification Programs BU
&
2
&
EUR 187,125.00
&
EUR 35,250.00
\\

Review -- Cross-Year Allocation Above 100\%
&
University Programs BU
&
1
&
EUR 45,000.00
&
EUR 13,500.00
\\

\bottomrule

\end{tabularx}

The eight Business Unit rows reconcile to 106 opportunities closing in 2026.
The 75 Ready records represent EUR 5,109,031.25 of estimated value,
EUR 2,332,441.25 of probability-weighted value and EUR 807,900.00 of
reconciled current-year value.

The 28 records classified as
\texttt{Not Ready - Missing Current-Year Allocation} represent
EUR 3,522,367.98 of estimated value and EUR 1,989,178.99 of probability-weighted
pipeline. Two additional records are classified first as
\texttt{Not Ready - Missing Probability}; they have EUR 50,000.00 of estimated
value. Their weighted source amount is zero, but this zero must not be interpreted
as a reliable economic valuation because Probability itself is missing.

\begin{inferencebox}

\textbf{Inference -- 2026 Forecast Readiness}

The readiness control separates usable current-year CRM contribution from
pipeline that is financially material but not yet forecastable. Seventy-five of
106 2026-closing opportunities, or 70.8\%, are ready and contribute
EUR 807,900.00 to the current-year CRM forecast. Twenty-eight records carry
EUR 1.989 million of weighted pipeline but cannot yet be allocated to 2026.
The two missing-Probability records are also unresolved even though their source
weighted value is zero. The single 110\% allocation record has a calculated
EUR 13,500.00 current-year amount, but that amount remains under review rather
than being included automatically.

\end{inferencebox}

\subsection{14.20 Forecast Readiness Status Summary}

The Business Unit detail was then collapsed into one row per readiness status in:

\begin{lstlisting}[style=powerquery]
DQ_Opportunity_2026ForecastReadinessStatusSummary
\end{lstlisting}

The completed grouped result is:

\begin{tabularx}{\textwidth}{
    >{\raggedright\arraybackslash}X
    >{\centering\arraybackslash}p{1.7cm}
    >{\raggedleft\arraybackslash}p{2.8cm}
    >{\raggedleft\arraybackslash}p{2.8cm}
}

\toprule

Forecast Readiness Status
&
Rows
&
Weighted Value
&
Current-Year Value
\\

\midrule

Ready -- Reconciled Current-Year Value
&
75
&
EUR 2,332,441.25
&
EUR 807,900.00
\\

Not Ready -- Missing Current-Year Allocation
&
28
&
EUR 1,989,178.99
&
--
\\

Not Ready -- Missing Probability
&
2
&
EUR 0.00
&
--
\\

Review -- Cross-Year Allocation Above 100\%
&
1
&
EUR 45,000.00
&
EUR 13,500.00
\\

\midrule

Total
&
106
&
EUR 4,366,620.24
&
EUR 821,400.00 calculated
\\

\bottomrule

\end{tabularx}

The EUR 821,400.00 calculated current-year total consists of EUR 807,900.00
from Ready opportunities plus EUR 13,500.00 from the one over-allocation review
record. Only the Ready portion is considered cleared for forecast use at the
present technical stage.

\begin{warningbox}

The total weighted value in the readiness table is not a complete measure of
economic potential because the two records with missing Probability have a source
weighted value of zero. Their EUR 50,000.00 estimated value remains visible as an
unresolved exposure. Likewise, the EUR 1,989,178.99 weighted pipeline associated
with missing current-year allocation must not be added directly to the
EUR 807,900.00 current-year forecast because its 2026 allocation is unknown.

\end{warningbox}

\begin{inferencebox}

\textbf{Inference -- Forecast Readiness Status Summary}

The grouped status control provides the first governed bridge between CRM data
quality and forecast use. The current technically cleared CRM contribution is
EUR 807,900.00. A further EUR 13,500.00 has been calculated but is held for
allocation review. EUR 1,989,178.99 remains probability-weighted pipeline at
risk because the current-year split is missing, while EUR 50,000.00 of estimated
pipeline remains unresolved because Probability is missing. These values have
different meanings and must remain in separate forecast buckets.

\end{inferencebox}

\subsection{14.21 Governed 2026 CRM Forecast Bridge}

The four forecast-readiness statuses were converted into explicit forecast
treatments in a dedicated reference query:

\begin{lstlisting}[style=powerquery]
DQ_Opportunity_2026ForecastBridge
\end{lstlisting}

The completed bridge contains four rows and four financially distinct value
columns:

\begin{tabularx}{\textwidth}{
    >{\raggedright\arraybackslash}p{5.3cm}
    >{\centering\arraybackslash}p{1.5cm}
    >{\raggedleft\arraybackslash}p{3.1cm}
    >{\raggedleft\arraybackslash}p{3.1cm}
}

\toprule

Forecast Treatment
&
Rows
&
Primary Governed Value
&
Amount
\\

\midrule

Include -- Ready Current-Year Forecast
&
75
&
Ready Current-Year Forecast
&
EUR 807,900.00
\\

At Risk -- Weighted Pipeline Not Allocated to 2026
&
28
&
At-Risk Weighted Pipeline
&
EUR 1,989,178.99
\\

At Risk -- Estimated Pipeline Missing Probability
&
2
&
Missing-Probability Estimated Value
&
EUR 50,000.00
\\

Hold -- Current-Year Value Under Review
&
1
&
Review Current-Year Value
&
EUR 13,500.00
\\

\bottomrule

\end{tabularx}

The query deliberately stores the four values in separate columns:
\texttt{ReadyCurrentYearForecastValue},
\texttt{ReviewCurrentYearValue},
\texttt{AtRiskWeightedPipeline} and
\texttt{MissingProbabilityEstimatedValue}. Nulls are retained in the
non-applicable columns.

\begin{inferencebox}

\textbf{Inference -- Governed CRM Forecast Bridge}

The CRM pipeline now has an explicit forecast-treatment layer. EUR 807,900.00 is
technically cleared for inclusion in the current-year CRM forecast. EUR 13,500.00
is calculated but held pending allocation review. EUR 1,989,178.99 remains
probability-weighted pipeline that cannot yet be allocated into 2026, while
EUR 50,000.00 of estimated opportunity value remains unresolved because
Probability is missing. These values are intentionally not added together as
one forecast measure because they represent different levels of readiness and
different financial meanings.

\end{inferencebox}

\subsection{14.22 Cross-Source Business Unit Vocabulary Audit}

Before creating the shared Business Unit dimension, the observed Business Unit
labels from all three analytical facts were combined and counted in:

\begin{lstlisting}[style=powerquery]
DQ_BusinessUnit_SourceValues
\end{lstlisting}

The completed result is:

\begin{tabularx}{\textwidth}{
    >{\raggedright\arraybackslash}p{3.4cm}
    >{\raggedright\arraybackslash}X
    >{\centering\arraybackslash}p{2.1cm}
}

\toprule

Source Dataset
&
BusinessUnitSource
&
Record Count
\\

\midrule

Actual Sales
&
ExEd
&
346
\\

Actual Sales
&
Qualification Programs
&
120
\\

Actual Sales
&
University Programs
&
329
\\

CRM Opportunities
&
ExEd Programs BU
&
92
\\

CRM Opportunities
&
Qualification Programs BU
&
7
\\

CRM Opportunities
&
University Programs BU
&
27
\\

Project Accrual
&
Missing / null
&
4
\\

Project Accrual
&
ExEd Programs BU
&
162
\\

Project Accrual
&
Qualification Programs BU
&
43
\\

Project Accrual
&
University Programs BU
&
104
\\

\bottomrule

\end{tabularx}

The grouped counts reconcile exactly to the fact-table record counts:
795 Actual Sales rows, 126 CRM Opportunity rows and 313 Project Accrual rows.

The audit shows that CRM Opportunities and Project Accrual already use the same
three populated labels:

\begin{itemize}[leftmargin=7mm,itemsep=2pt]
    \item \texttt{ExEd Programs BU};
    \item \texttt{Qualification Programs BU}; and
    \item \texttt{University Programs BU}.
\end{itemize}

Actual Sales represents the same three business areas with shorter labels:
\texttt{ExEd}, \texttt{Qualification Programs} and
\texttt{University Programs}. The four null accrual records are not assigned to
any valid Business Unit automatically.

\begin{inferencebox}

\textbf{Inference -- Business Unit Vocabulary Audit}

The cross-source Business Unit mismatch is semantic rather than structural.
CRM and Accrual already share one common vocabulary, while Actual Sales uses
shorter names for the same three categories. This can be standardized safely
through an explicit mapping layer. The four accrual records with missing
Business Unit must remain mapped to an \texttt{Unassigned / Missing BU}
technical member rather than being guessed into one of the three real Business
Units. The next modelling step is therefore to create a canonical
\texttt{DimBusinessUnit} table before modifying fact-table labels.

\end{inferencebox}

\subsection{14.23 Canonical CRM Business Unit Key Mapping}

A numeric \texttt{BusinessUnitKey} was added to
\texttt{Fact\_Opportunities} while retaining the source
\texttt{BusinessUnitSource} field:

\begin{lstlisting}[style=powerquery]
if [BusinessUnitSource] = null
    or Text.Trim(Text.From([BusinessUnitSource])) = "" then
    0
else if [BusinessUnitSource] = "ExEd Programs BU" then
    1
else if [BusinessUnitSource] = "Qualification Programs BU" then
    2
else if [BusinessUnitSource] = "University Programs BU" then
    3
else
    null
\end{lstlisting}

The field was typed as Whole Number / \texttt{Int64.Type}. The completed
profile contains 126 rows, 100\% valid values, 0 errors, 0 empty values and
3 distinct keys. The source audit reconciles the distribution to 92 records on
key 1, 7 records on key 2 and 27 records on key 3. No CRM Opportunity maps to
key 0.

\begin{inferencebox}

\textbf{Inference -- CRM Business Unit Mapping}

The CRM fact is now aligned to the same governed Business Unit key structure as
Actual Sales and Project Accrual. Because every Opportunity has one of the three
known Business Units, no CRM record falls into the technical Unassigned member.
The original CRM label remains available for lineage while the numeric key
provides a stable relationship field.

\end{inferencebox}

% ============================================================
% CURRENT STATUS
% ============================================================

\section{15. Current Implementation Status}

\begin{tabularx}{\textwidth}{
    >{\bfseries}p{5.8cm}
    p{2.8cm}
    X
}

\toprule

Implementation Item
&
Status
&
Result
\\

\midrule

Source files imported
&
Completed
&
Three Excel sources available
\\

Staging queries created
&
Completed
&
Loading disabled
\\

Fact references created
&
Completed
&
Three analytical facts available
\\

Actual sales transformation
&
Completed
&
795 records prepared
\\

Actual quality classification
&
Completed
&
No critical issues
\\

Actual project-month validation
&
Completed
&
No duplicate keys
\\

Accrual date preparation
&
Completed
&
September--December 2026
\\

Accrual Total identification
&
Completed
&
One source Total identified
\\

Accrual Total separation
&
Completed
&
Control query created
\\

Accrual detail filtering
&
Completed
&
313 detail rows retained
\\

Accrual quality classification
&
Completed
&
308 valid and 5 exception records
\\

Accrual column standardization
&
Completed
&
Eight final analytical columns
\\

Accrual reconciliation
&
Completed
&
EUR 0.00 variance
\\

Accrual duplicate investigation
&
Completed
&
8 groups and 12 excess records
\\

Repeated business-key investigation
&
Completed
&
10 groups containing 32 source rows
\\

Distinct forecast-amount analysis
&
Completed
&
4 same-amount and 6 multiple-amount groups
\\

Repeated-key status classification
&
Completed
&
4 same, 4 mixed and 2 component groups
\\

CRM opportunity preparation
&
Completed technically
&
Transformation, reconciliation, allocation, forecast-readiness and BU-key controls completed
\\

CRM source profiling
&
Completed
&
Entire dataset evaluated
\\

CRM data-type standardization
&
Completed
&
14 columns typed with 0 conversion errors
\\

CRM text standardization
&
Completed
&
Four text columns received Trim and Clean
\\

CRM analytical renaming
&
Completed
&
14 source fields renamed consistently
\\

CRM quality classification
&
Completed
&
126 classified rows, 0 errors and 0 empty statuses
\\

CRM weighted-value reconciliation
&
Completed
&
123 Match, 2 Not Testable and 1 Mismatch
\\

CRM weighted-value mismatch control
&
Completed technically
&
Dedicated DQ query created; one material exception remains for business review
\\

Opportunity value reconciliation
&
Completed
&
126 records reconcile; one distinct check status: Match
\\

Current-year weighted calculation
&
Completed
&
67\% populated, 33\% null, 0 errors
\\

Current-year weighted variance
&
Completed
&
67\% populated, 33\% null, 0 errors
\\

Current-year weighted check
&
Completed
&
100\% valid, 0 errors, 0 empty values and 3 distinct statuses
\\

Current-year status-count summary
&
Completed
&
85 Match, 39 Missing Current-Year Allocation and 2 Missing Probability
\\

Cross-year allocation validation
&
Completed
&
126 records classified; 70 total 100\%, 52 missing input, 3 below 100\%, 1 above 100\%
\\

Allocation-status summary query
&
Completed
&
6 grouped statuses reconcile exactly to 126 records
\\

Allocation exception isolation
&
Completed
&
4 records isolated: totals 110\%, 50\%, 50\% and 30\%
\\

Missing-allocation date-horizon analysis
&
Completed
&
52 records reconciled; 41 (78.8\%) close in 2026
\\

2026 missing-allocation detail
&
Completed
&
41 rows isolated: 21 both missing, 9 missing current-year, 11 missing next-year
\\

2026 missing-allocation financial exposure
&
Completed
&
41 records: EUR 4.209M estimated and EUR 2.318M weighted exposure
\\

2026 current-year allocation risk
&
Completed
&
30 records: EUR 3.572M estimated and EUR 1.989M weighted exposure
\\

2026 opportunity forecast-readiness summary
&
Completed
&
106 records: 75 Ready, 28 missing current-year allocation, 2 missing Probability, 1 Review
\\

2026 forecast-readiness status summary
&
Completed
&
Ready current-year value EUR 807,900.00; review value EUR 13,500.00
\\

2026 CRM forecast bridge
&
Completed
&
4 governed treatments: EUR 807.9k Ready, EUR 13.5k Hold, EUR 1.989M weighted risk, EUR 50k missing-Probability exposure
\\

Cross-source Business Unit vocabulary audit
&
Completed
&
Actuals use 3 shorter labels; CRM and Accrual share 3 canonical labels; 4 Accrual rows are missing BU
\\

Business Unit dimension
&
Completed
&
4 members: 3 canonical Business Units plus key 0 Unassigned / Missing BU
\\

Actual Sales Business Unit key
&
Completed
&
795 rows mapped; 100\% valid, 0 errors, 0 empty keys and 3 distinct keys
\\

Project Accrual Business Unit key
&
Completed
&
313 rows mapped; 100\% valid, 0 errors, 0 empty keys and 4 distinct keys
\\

CRM Opportunity Business Unit key
&
Completed
&
126 rows mapped; 100\% valid, 0 errors, 0 empty keys and 3 distinct keys
\\

Staging and DQ model-load cleanup
&
Completed
&
Staging/DQ tables remain hidden from the model; current visible model has 2 dimensions and 3 facts
\\

Business Unit semantic relationships
&
Completed
&
3 active 1:* relationships from DimBusinessUnit with Single cross-filter direction
\\

Business Unit relationship validation
&
Completed
&
Exact BU row counts reconcile: 795 Actual, 313 Accrual and 126 Opportunities
\\

2026 Date dimension
&
Completed technically
&
365 unique dates, 0 errors, 0 empty dates, 12 months and 4 quarters
\\

DimDate query-group placement
&
Completed
&
DimDate positioned in the governed dimension layer
\\

Date semantic relationships
&
Completed
&
Two active *:1 relationships to DimDate with Single cross-filter direction
\\

Date relationship validation
&
Completed
&
795 Actual rows map to January--August and 313 Accrual rows to September--December
\\

Month-name chronological sort
&
Completed
&
MonthName displays January--December using MonthNumber
\\

Date-table designation
&
Completed
&
DimDate formally marked as the model Date table using DimDate[Date]
\\

Base financial measures
&
Completed
&
Actual EUR 14.635M + Accrual EUR 6.616M = Base Full-Year Forecast EUR 21.251M
\\

Base financial monthly validation
&
Completed
&
Actual value appears Jan--Aug; Accrual value appears Sep--Dec; 12-month base series is continuous
\\

Business Unit financial validation
&
Completed
&
BU rows reconcile exactly to EUR 14.635M Actual, EUR 6.616M Accrual and EUR 21.251M Base Forecast
\\

CRM Ready measures
&
Completed
&
75 Ready opportunities; EUR 807,900.00 reconciled current-year CRM value
\\

CRM Ready Business Unit validation
&
Completed
&
ExEd 60 / EUR 527,877.50; Qualification 2 / EUR 35,250.00; University 13 / EUR 244,772.50
\\

Forecast Including Ready CRM
&
Completed
&
EUR 21,251,093.59 Base Forecast + EUR 807,900.00 Ready CRM = EUR 22,058,993.59
\\

Semantic model and DAX
&
In progress
&
Base forecast and Ready CRM inclusion completed; Review and At-Risk CRM measures are next
\\

Semantic CRM governance measures
&
Completed
&
Review EUR 13,500.00; At-Risk EUR 1,989,178.99; Missing Probability EUR 50,000.00;
review-approved forecast EUR 22,072,493.59
\\

Disconnected forecast-bucket dimension
&
Completed
&
\texttt{DimForecastBucket} exposes five governed forecast treatments without
creating a new fact relationship
\\

CRM Date dimension
&
Completed
&
\texttt{DimCRMDate}: 1,827 unique dates covering 2024--2028 with 0 errors
\\

CRM Date-role relationships
&
Completed
&
Active ClosingDate and inactive SimulatedProgramStartDate relationships to
\texttt{DimCRMDate}; both Single direction
\\

CRM Date-role DAX validation
&
Completed
&
Both roles reconcile to 126 opportunities; 2026 split is 106 by Closing Date
versus 56 by Program Start
\\

Project dimension construction
&
Completed
&
\texttt{DimProject}: 252 rows, 252 distinct/unique ProjectIDs, 0 errors and 0 empty values
\\

Project dimension governance placement
&
Completed
&
\texttt{DimProject} moved to \texttt{04\_Dimensions} with model load retained
\\

Project semantic relationships
&
Completed
&
Two active *:1 relationships from Actual Sales and Project Accrual to
\texttt{DimProject[ProjectID]}, both with Single cross-filter direction
\\

Project relationship validation
&
Completed
&
795 Actual rows and 313 Accrual rows reconcile to EUR 14.635M Actual,
EUR 6.616M Accrual and EUR 21.251M Base Full-Year Forecast by Project
\\

Project slicer functional validation
&
Completed
&
\texttt{O7901009}: 6 Actual rows, 4 Accrual rows, EUR 439,132.89 Actual,
EUR 308,913.00 Accrual and EUR 748,045.89 Base Forecast; CRM remains independent
\\

Top-10 Project forecast ranking
&
Completed
&
Horizontal bar chart ranks the ten largest \texttt{ProjectID} values by
\texttt{Base Full-Year Forecast} in descending order
\\

Project forecast composition visual
&
Next
&
Show Actual Net Sales versus Project Accrual Forecast for the same Top-10 projects
\\

Dashboard implementation
&
In progress
&
Validation pages and forecast-control visuals exist; final executive layout,
Project ranking and presentation polish remain
\\

Project forecast composition status update
&
Completed
&
Top-10 companion stacked bar chart shows Actual Net Sales versus remaining
Project Accrual Forecast using the validated Project dimension
\\

Executive Dashboard KPI status
&
Completed
&
Eight EUR-formatted KPI cards reconcile to the governed forecast and CRM
treatment measures
\\

Executive Dashboard Business Unit slicer
&
Completed
&
Four governed Business Unit members are available from
\texttt{DimBusinessUnit[BusinessUnitName]}
\\

Executive monthly forecast visual
&
Correction required
&
January--December order is correct and Actual/Accrual period separation is visible,
but the current 100\% stacked chart must be changed to a standard stacked column
chart to display financial magnitude correctly
\\

Executive monthly forecast correction -- status update
&
Completed
&
Standard stacked column chart now uses a monetary axis and preserves the validated
January--August Actual / September--December Accrual pattern
\\

Business Unit slicer interactive validation
&
Completed
&
All four governed Business Unit members correctly filter core financial KPIs,
CRM governance KPIs and the monthly Actual-versus-Accrual visual
\\

Business Unit forecast-composition visual
&
Next
&
Create the executive comparison of Actual Net Sales versus Project Accrual Forecast
by \texttt{DimBusinessUnit[BusinessUnitName]}
\\


Business Unit forecast-composition status -- final
&
Completed
&
Executive stacked bar now responds correctly to all four governed Business Unit selections
\\

Forecast Governance table interaction validation
&
Completed
&
Five governed forecast buckets filter consistently by Business Unit without creating
a physical relationship from the disconnected presentation table
\\

Final semantic relationship audit
&
Completed
&
9 relationships reviewed: 8 active and 1 intentionally inactive CRM date-role relationship;
all fact-to-dimension relationships use controlled many-to-one / Single-direction filtering
\\

Executive Dashboard end-to-end Business Unit validation
&
Completed
&
ExEd, Qualification Programs, University Programs and Unassigned / Missing BU were
re-tested against KPI, governance, monthly, Business Unit and Top-10 Project visuals
\\

Technical field visibility cleanup
&
Completed
&
Surrogate keys, sort helpers and selected data-quality / diagnostic fields are hidden
from business users while retained in the semantic model for governance
\\

Final financial reconciliation sign-off
&
Next
&
Add the final reconciliation controls and verify zero variance for Base Forecast and
Forecast Including Ready CRM before final layout polish and interview delivery
\\

\bottomrule

\end{tabularx}

\subsection{15.1 Final Semantic Relationship Audit}

The semantic model was re-checked in the Power BI
\texttt{Manage relationships} dialog after the Executive Dashboard and field-visibility
cleanup were completed. The final relationship inventory is:

\begin{tabularx}{\textwidth}{
    >{\ttfamily\raggedright\arraybackslash}p{4.1cm}
    >{\ttfamily\raggedright\arraybackslash}p{4.1cm}
    >{\centering\arraybackslash}p{1.7cm}
    >{\centering\arraybackslash}p{1.7cm}
    X
}

\toprule
Fact Side
&
Dimension Side
&
Status
&
Cardinality
&
Filter
\\
\midrule

Fact\_ActualSales[BusinessUnitKey]
&
DimBusinessUnit[BusinessUnitKey]
&
Active
&
*:1
&
Single
\\

Fact\_ActualSales[MonthStart]
&
DimDate[Date]
&
Active
&
*:1
&
Single
\\

Fact\_ActualSales[ProjectID]
&
DimProject[ProjectID]
&
Active
&
*:1
&
Single
\\

Fact\_Opportunities[BusinessUnitKey]
&
DimBusinessUnit[BusinessUnitKey]
&
Active
&
*:1
&
Single
\\

Fact\_Opportunities[ClosingDate]
&
DimCRMDate[Date]
&
Active
&
*:1
&
Single
\\

Fact\_Opportunities[SimulatedProgramStartDate]
&
DimCRMDate[Date]
&
Inactive
&
*:1
&
Single
\\

Fact\_ProjectAccrual[BusinessUnitKey]
&
DimBusinessUnit[BusinessUnitKey]
&
Active
&
*:1
&
Single
\\

Fact\_ProjectAccrual[ForecastMonth]
&
DimDate[Date]
&
Active
&
*:1
&
Single
\\

Fact\_ProjectAccrual[ProjectID]
&
DimProject[ProjectID]
&
Active
&
*:1
&
Single
\\

\bottomrule
\end{tabularx}

\begin{inferencebox}
\textbf{Inference -- Final Relationship Audit}

The final semantic model retains a controlled star-schema pattern. Eight relationships
are active and provide the intended Business Unit, Date, CRM closing-date and Project
filter paths. The second CRM date role remains deliberately inactive so the same
opportunity fact can support an alternative program-start analysis through explicit DAX
without introducing ambiguous simultaneous date filtering.
\end{inferencebox}

\subsection{15.2 Executive Dashboard End-to-End Business Unit Validation}

After the Forecast Governance table and the remaining executive visuals were added,
the governed Business Unit slicer was re-tested against the complete page. The four
available selections produced the following control values:

\begin{tabularx}{\textwidth}{
    >{\raggedright\arraybackslash}p{3.6cm}
    >{\raggedleft\arraybackslash}p{2.3cm}
    >{\raggedleft\arraybackslash}p{2.3cm}
    >{\raggedleft\arraybackslash}p{2.3cm}
    >{\raggedleft\arraybackslash}X
}

\toprule
Business Unit
&
Actual
&
Accrual
&
Base Forecast
&
Ready CRM
\\
\midrule

ExEd Programs BU
&
EUR 6.04M
&
EUR 3.01M
&
EUR 9.05M
&
EUR 527.88K
\\

Qualification Programs BU
&
EUR 4.87M
&
EUR 1.77M
&
EUR 6.64M
&
EUR 35.25K
\\

University Programs BU
&
EUR 3.72M
&
EUR 1.79M
&
EUR 5.52M
&
EUR 244.77K
\\

Unassigned / Missing BU
&
--
&
EUR 42.55K
&
EUR 42.55K
&
--
\\

\bottomrule
\end{tabularx}

The CRM governance values also remain consistent with the source controls:
ExEd carries approximately EUR 1.23M of at-risk weighted pipeline, Qualification
approximately EUR 615.63K, University approximately EUR 143.94K, while the
EUR 13.50K review amount and EUR 50.00K missing-Probability exposure appear under
University Programs BU. The Unassigned member contains only the known EUR 42,550.00
Accrual amount and no fabricated Actual or CRM contribution.

\begin{inferencebox}
\textbf{Inference -- End-to-End Business Unit Interaction}

The complete Executive Dashboard now demonstrates that one conformed Business Unit
dimension controls all business-facing financial and CRM-governance visuals consistently.
The explicit Unassigned member behaves as a governance exception rather than leaking
into a real Business Unit, while the three populated Business Units preserve their
expected financial and CRM risk profiles.
\end{inferencebox}

\subsection{15.3 Technical Field Visibility Cleanup}

The model was then cleaned for business-facing use. Technical and control fields needed
for joins, sorting, reconciliation and data-quality diagnostics remain physically present
in the semantic model, but selected fields are hidden from report authors and consumers.

Confirmed examples include technical Business Unit sort/key helpers, hidden calendar
sorting fields such as \texttt{YearMonth}, the Project surrogate
\texttt{ProjectKey}, and selected row-quality or diagnostic fields in the fact tables.
Business-facing dimensions, descriptive attributes and governed measures remain visible.

\begin{inferencebox}
\textbf{Inference -- Semantic Field Governance}

Hiding technical fields improves usability without weakening auditability. The report
author is guided toward governed dimensions and measures, while the model still retains
the keys, sort columns and diagnostic evidence required for maintenance, troubleshooting
and future extension.
\end{inferencebox}

\subsection{15.4 Current Completion Point and Next Control}

The analytical modelling, relationship construction, main DAX measures, Business Unit
and Project filter tests, executive KPI layer, monthly forecast view, Business Unit
composition, Top-10 Project composition and Forecast Governance table are now complete.

The next technical control is a final financial reconciliation sign-off. Two zero-variance
measures should be added:

\begin{lstlisting}[style=powerquery]
Base Forecast Reconciliation Variance =
    [Base Full-Year Forecast]
    - (
        [Actual Net Sales]
        + [Project Accrual Forecast]
      )
\end{lstlisting}

\begin{lstlisting}[style=powerquery]
Ready CRM Inclusion Variance =
    [Forecast Including Ready CRM]
    - (
        [Base Full-Year Forecast]
        + [CRM Ready Current-Year Value]
      )
\end{lstlisting}

Both controls are expected to return EUR 0.00 at the all-Business-Unit level and under
valid Business Unit selections. After this sign-off, the remaining work is presentation
polish, resetting slicers to the intended default state, saving the final checkpoint and
preparing the short interview narrative.

\begin{inferencebox}
\textbf{Inference -- Final Technical Readiness}

The solution has moved from data preparation into final assurance. No new forecasting
logic is required for the core case; the remaining technical task is to prove that the
already validated components reconcile algebraically in the finished report context.
A zero-variance sign-off provides a compact final control before presentation.
\end{inferencebox}


% ============================================================
% GOVERNANCE VALUE
% ============================================================

\section{16. Governance and Process Value}

The implemented approach provides:

\begin{itemize}[leftmargin=7mm,itemsep=3pt]
    \item transparent source-to-report lineage;
    \item repeatable transformations during every refresh;
    \item explicit row-level data-quality controls;
    \item documented handling of exceptional financial values;
    \item validation of the intended analytical grain;
    \item separation of raw, analytical and control queries;
    \item controlled handling of source Total rows;
    \item prevention of accidental forecast double-counting;
    \item distinction between exact duplicates and legitimate forecast
    components;
    \item explicit CRM data types with zero conversion errors;
    \item CRM text cleaning, standardized analytical names and complete
    row-level quality classification;
    \item independent weighted-value recalculation with an explicit financial
    tolerance;
    \item explicit reconciliation of \texttt{OpportunityValue} and
    \texttt{EstimatedValue};
    \item independent current-year weighted-value calculation using probability
    and annual allocation;
    \item null-safe distinction between numerical mismatches and records that
    cannot yet be tested because a business input is missing;
    \item documented correction of Power Query dependency order without changing
    source data;
    \item a dedicated current-year weighted-status summary that reconciles all
    126 opportunities and separates 85 testable matches from 41 records with
    missing required inputs;
    \item explicit cross-year allocation validation that separates complete
    100\% allocations, incomplete allocation data and true over-allocation
    exceptions without imputing missing values;
    \item an auditable grouped allocation-status summary that reconciles all
    126 opportunities;
    \item targeted isolation of four non-100\% allocation records together with
    closing and simulated program-start dates for business interpretation;
    \item horizon analysis of all 52 missing-allocation records, showing that
    41 close in 2026;
    \item a dedicated 41-row near-term review set for missing-allocation
    opportunities;
    \item financial-exposure quantification for the 41 near-term allocation gaps;
    \item a Business Unit risk summary for the 30 opportunities that cannot yet be
    allocated reliably into the current financial year;
    \item a 2026 forecast-readiness classification across all 106 near-term CRM
    opportunities;
    \item an auditable four-status summary separating EUR 807,900.00 of Ready
    current-year CRM value from review and unresolved pipeline exposures;
    \item a governed CRM forecast bridge that keeps forecast inclusion, held value
    and unresolved pipeline in separate financial columns;
    \item a cross-source Business Unit vocabulary audit covering all three analytical
    facts;
    \item a governed \texttt{DimBusinessUnit} with stable numeric keys and an explicit
    unassigned member;
    \item preservation of source Business Unit labels alongside canonical keys for
    lineage;
    \item completed canonical BU-key mapping for all 795 Actual Sales records;
    \item completed canonical BU-key mapping for all 313 Project Accrual records,
    including 4 explicit Unassigned rows;
    \item completed canonical BU-key mapping for all 126 CRM Opportunity records;
    \item removal of staging and DQ helper tables from the visible semantic model
    without deleting their Power Query controls;
    \item three active one-to-many Business Unit relationships with single-direction
    filtering and no fact-to-fact relationship;
    \item explicit DAX-based validation that the shared Business Unit dimension
    propagates filters correctly to all three facts;
    \item exact reconciliation of Business Unit row counts to 795 Actual,
    313 Accrual and 126 CRM Opportunity records;
    \item a complete, unique and error-free 2026 Date dimension with daily,
    monthly and quarterly attributes;
    \item two active Single-direction Date relationships from the shared calendar
    to Actual Sales and Project Accrual;
    \item explicit monthly validation that January--August contains only Actual
    rows while September--December contains only Accrual rows at the current
    source coverage;
    \item chronological MonthName sorting governed by the numeric MonthNumber;
    \item explicit DAX measures for Actual Net Sales, Project Accrual Forecast and
    the combined Base Full-Year Forecast;
    \item exact source-level reconciliation of the base financial totals to
    EUR 14,635,067.81 Actual + EUR 6,616,025.78 Accrual =
    EUR 21,251,093.59;
    \item formal designation of \texttt{DimDate} as the model Date table;
    \item financial-value validation by Business Unit, including the explicit
    EUR 42,550.00 Unassigned accrual amount;
    \item a governed Ready-CRM DAX layer that includes only 75 technically cleared
    opportunities and EUR 807,900.00 of reconciled current-year value;
    \item a current governed annual forecast of EUR 22,058,993.59 after Ready CRM
    is added to the source-aligned base forecast;
    \item semantic CRM governance measures that keep Review, At-Risk and
    Missing-Probability exposures separate from the cleared forecast;
    \item an optional review-approved scenario of EUR 22,072,493.59 without
    automatically including unresolved pipeline;
    \item a disconnected forecast-bucket presentation dimension that introduces
    no additional fact-filter path;
    \item a complete 2024--2028 CRM calendar and explicit active/inactive
    role-playing relationships for Closing Date and Program Start;
    \item DAX validation proving that both CRM date roles reconcile to all
    126 opportunities while retaining their different business timing;
    \item a governed Project dimension containing 252 clean, unique ProjectIDs;
    \item two active Single-direction Project relationships from the dimension to
    Actual Sales and Project Accrual using the validated natural \texttt{ProjectID};
    \item Project-level reconciliation proving the new dimension preserves all
    795 Actual rows, 313 Accrual rows and the complete EUR 21,251,093.59 base forecast;
    \item deliberate exclusion of CRM from the Project relationship path because the
    supplied opportunity source does not contain a confirmed ProjectID mapping;
    \item an interactive Project slicer test proving that one ProjectID filters both
    project-based facts without changing the independent CRM governance layer;
    \item a Top-10 Project forecast ranking that converts the validated Project dimension
    into an executive prioritization view based on the full-year base forecast; and
    \item a foundation for reconciliation, Business Unit harmonization, shared
    time filtering, controlled CRM integration, Project analysis and exception reporting.
\end{itemize}

% ============================================================
% NEXT IMPLEMENTATION PHASES
% ============================================================

\section{17. Next Implementation Phases}

The implementation sequence identified earlier is preserved below. Completed
items are marked as completed rather than removed from the report:

\begin{enumerate}[leftmargin=9mm,itemsep=4pt]
    \item validate the 12 excess exact-duplicate records with Finance
    (\textit{business validation pending});
    \item validate the four records with missing Business Unit values
    (\textit{business validation pending});
    \item investigate the single CRM weighted-value mismatch and the two
    records with missing Probability
    (\textit{technical isolation completed; business validation pending});
    \item reconcile \texttt{OpportunityValue} with \texttt{EstimatedValue}
    (\textit{completed: all 126 records are Match within EUR 0.01});
    \item validate current-year and next-year allocation percentages
    (\textit{technical control completed: 70 records total 100\%, 29 have both
    allocations missing, 12 miss current-year allocation, 11 miss next-year
    allocation, 3 are below 100\%, and 1 is above 100\%; business interpretation
    remains open});
    \item summarize the exact record counts of the three
    \texttt{CurrentYearWeightedCheck} statuses
    (\textit{completed: 85 Match, 39 Missing Current-Year Allocation,
    2 Missing Probability});
    \item summarize the six allocation-validation statuses
    (\textit{completed: 70, 29, 12, 11, 3 and 1 records});
    \item isolate and review the four complete-allocation exceptions
    (\textit{completed technically: totals of 110\%, 50\%, 50\% and 30\%;
    business interpretation pending});
    \item analyze the 52 records with one or both annual allocations missing by
    closing year and simulated program-start year
    (\textit{completed: 41 close in 2026, 8 in 2027, 1 in 2028 and 2 in 2024--2025});
    \item isolate the near-term 2026 missing-allocation records
    (\textit{completed: 41 rows, comprising 21 both missing, 9 missing current-year
    and 11 missing next-year allocation});
    \item quantify the financial exposure of the 41 2026 missing-allocation records
    using estimated value and probability-weighted value
    (\textit{completed: EUR 4,208,699.23 estimated value and
    EUR 2,317,982.74 weighted value});
    \item isolate 2026 current-year allocation risk by Business Unit
    (\textit{completed: 30 records, EUR 3,572,367.98 estimated value and
    EUR 1,989,178.99 weighted value});
    \item classify 2026 CRM opportunities by forecast readiness and quantify
    reconciled current-year value versus unresolved pipeline exposure
    (\textit{completed: 75 Ready, 28 Missing Current-Year Allocation,
    2 Missing Probability and 1 Cross-Year Allocation Review});
    \item summarize the 2026 forecast-readiness statuses
    (\textit{completed: EUR 807,900.00 Ready current-year value,
    EUR 13,500.00 review value and EUR 1,989,178.99 unresolved weighted pipeline});
    \item create a governed 2026 CRM forecast bridge that keeps Ready forecast,
    Review value and unresolved pipeline exposure separate
    (\textit{completed: EUR 807,900.00 Ready, EUR 13,500.00 Hold,
    EUR 1,989,178.99 weighted risk and EUR 50,000.00 missing-Probability exposure});
    \item audit Business Unit source labels across all three analytical facts
    (\textit{completed: Actual Sales uses 3 shorter labels; CRM and Accrual share
    3 canonical labels; 4 Accrual rows remain missing});
    \item create the canonical \texttt{DimBusinessUnit} table with an explicit
    \texttt{Unassigned / Missing BU} member
    (\textit{completed: keys 1--3 for the canonical BUs and key 0 for Unassigned});
    \item add the canonical \texttt{BusinessUnitKey} to \texttt{Fact\_ActualSales}
    while retaining \texttt{BusinessUnitSource}
    (\textit{completed: all 795 rows mapped with 0 errors and 0 empty keys});
    \item add the canonical \texttt{BusinessUnitKey} to
    \texttt{Fact\_ProjectAccrual}, including key 0 for the 4 missing-BU rows
    (\textit{completed: 313 rows, 100\% valid and 4 distinct keys});
    \item add the canonical \texttt{BusinessUnitKey} to
    \texttt{Fact\_Opportunities}
    (\textit{completed: 126 rows, 100\% valid and 3 distinct keys});
    \item disable model load for staging and DQ queries while retaining their
    transformation and audit logic
    (\textit{completed: only the dimension and three facts remain visible});
    \item standardize Business Unit filtering across all datasets using the canonical
    dimension
    (\textit{relationship structure completed: three active 1:* Single-direction relationships});
    \item validate Business Unit relationship filter propagation with row-count
    measures
    (\textit{completed: ExEd 346/162/92, Qualification 120/43/7,
    University 329/104/27 and Unassigned 0/4/0 across Actual/Accrual/CRM});
    \item create the shared 2026 \texttt{DimDate}
    (\textit{completed technically: 365 unique dates, 12 months and 4 quarters});
    \item move \texttt{DimDate} from \texttt{01\_Staging} to
    \texttt{04\_Dimensions}
    (\textit{completed});
    \item create Date relationships from \texttt{DimDate[Date]} to
    \texttt{Fact\_ActualSales[MonthStart]} and
    \texttt{Fact\_ProjectAccrual[ForecastMonth]}
    (\textit{completed: both active, many-to-one from fact perspective,
    Single cross-filter direction});
    \item validate monthly Date-filter propagation across Actual Sales and
    Project Accrual
    (\textit{completed: 795 Actual rows in January--August and 313 Accrual rows
    in September--December});
    \item sort \texttt{DimDate[MonthName]} by \texttt{MonthNumber}
    (\textit{completed: January--December chronological order confirmed});
    \item mark \texttt{DimDate} as the model Date table using
    \texttt{DimDate[Date]}
    (\textit{completed});
    \item create the first financial DAX measures for Actual Net Sales,
    Project Accrual Forecast and their combined base full-year forecast
    (\textit{completed: EUR 14,635,067.81 + EUR 6,616,025.78 =
    EUR 21,251,093.59});
    \item validate the base financial measures by
    \texttt{DimBusinessUnit[BusinessUnitName]}
    (\textit{completed: BU rows reconcile exactly to the annual totals});
    \item create governed Ready-CRM measures and validate them by Business Unit
    (\textit{completed: 75 opportunities and EUR 807,900.00 current-year value});
    \item create \texttt{Forecast Including Ready CRM}
    (\textit{completed: EUR 22,058,993.59});
    \item surface the remaining CRM governance buckets as semantic measures:
    EUR 13,500.00 Review current-year value, EUR 1,989,178.99 At-Risk weighted
    pipeline and EUR 50,000.00 Missing-Probability estimated exposure
    (\textit{immediate next technical control});
    \item[] \textit{Status update: completed. The three measures are validated and
    the optional review-approved scenario equals EUR 22,072,493.59.}
    \item define and validate the CRM monthly date-role strategy before placing
    current-year CRM value on the shared Date axis;
    \item[] \textit{Status update: completed. \texttt{DimCRMDate} covers
    2024--2028; Closing Date is active, Program Start is inactive, and both roles
    reconcile to 126 records through DAX validation.}
    \item create the Project dimension;
    \item[] \textit{Status update: completed. \texttt{DimProject} contains 252 clean,
    distinct and unique ProjectIDs and is governed under \texttt{04\_Dimensions}.}
    \item implement and validate Project relationships;
    \item[] \textit{Status update: completed. Actual Sales and Project Accrual use
    active Single-direction *:1 relationships to \texttt{DimProject[ProjectID]};
    totals reconcile to 795 Actual rows, 313 Accrual rows and EUR 21,251,093.59.}
    \item continue the Power BI semantic model and relationships;
    \item[] \textit{Status update: the core Business Unit, Date, CRM Date-role and
    Project relationship layers are now technically complete.}
    \item create forecast scenarios and governed CRM-inclusion measures;
    \item[] \textit{Status update: completed for the current case, including Ready CRM,
    Review, At-Risk and Missing-Probability treatments.}
    \item complete the final interactive Project filter validation;
    \item[] \textit{Status update: completed. Selecting \texttt{O7901009} produces
    EUR 439,132.89 Actual Net Sales + EUR 308,913.00 Project Accrual Forecast =
    EUR 748,045.89 Base Full-Year Forecast, while CRM remains outside the Project path.}
    \item create the Top-10 Project forecast-ranking visual;
    \item[] \textit{Status update: completed. The report now contains a descending
    horizontal ranking of the ten projects with the largest Base Full-Year Forecast.}
    \item build executive, detailed and data-quality dashboard pages
    (\textit{current immediate implementation phase: Top-10 Actual-versus-Accrual
    composition and final presentation polish});
    \item[] \textit{Status update: the Top-10 composition visual, Executive KPI layer,
    monthly-chart correction and four-member Business Unit slicer validation are
    completed. The immediate next dashboard visual is Business Unit forecast
    composition.}
    \item document the future automated architecture; and
    \item prepare the final 15-minute interview presentation.
\end{enumerate}

\subsection{17.1 Immediate Weighted-Value Mismatch Investigation}

The step originally documented as the immediate mismatch investigation has now
been completed technically. The planned control query has been created as a
reference from \texttt{Fact\_Opportunities}:

\begin{lstlisting}[style=powerquery]
DQ_Opportunity_WeightedValueMismatch
\end{lstlisting}

The control retains only records where:

\begin{lstlisting}[style=powerquery]
WeightedValueCheck = "Mismatch"
\end{lstlisting}

The reconciliation documented in Section 14.7 identifies one record with a
source-minus-calculated variance of EUR -468,750.00. The following fields remain
the review set for business validation:

\begin{itemize}[leftmargin=7mm,itemsep=2pt]
    \item \texttt{OpportunityName};
    \item \texttt{BusinessUnitSource};
    \item \texttt{ClosingDate};
    \item \texttt{OpportunityQualityStatus};
    \item \texttt{OpportunityValue};
    \item \texttt{EstimatedValue};
    \item \texttt{Probability};
    \item \texttt{WeightedEstimatedValue};
    \item \texttt{CalculatedWeightedEstimatedValue};
    \item \texttt{WeightedValueVariance}; and
    \item \texttt{WeightedValueCheck}.
\end{itemize}

The query is diagnostic only. The source value remains unchanged until the
underlying opportunity and its business calculation are confirmed.

\subsection{17.2 Outstanding Business Validation}

The following exceptions remain open for business confirmation even though the
corresponding technical controls have been implemented:

\begin{tabularx}{\textwidth}{
    >{\raggedright\arraybackslash\bfseries\color{AaltoNavy}}p{7cm}
    X
}

\toprule

Business Validation Item
&
Current Result
\\

\midrule

Exact duplicate excess rows
&
12 rows requiring Finance confirmation
\\

Potential duplicate overstatement
&
EUR 175,318.63 diagnostic estimate
\\

Missing Business Unit classification
&
4 rows with EUR 42,550.00 forecast value
\\

CRM weighted-value mismatch
&
1 row with EUR -468,750.00 source-minus-calculated variance
\\

CRM records not testable for weighted value
&
2 rows with missing Probability
\\

CRM allocation above 100\%
&
Open custom opportunity 111: 30\% + 80\% = 110\%
\\

CRM allocation below 100\%
&
3 records: totals of 50\%, 50\% and 30\% requiring business interpretation
\\

Below-100\% date patterns
&
2026/2026, 2028/2028 and 2026/2027 closing/program-start combinations
\\

CRM records missing both annual allocations
&
29 records
\\

CRM records missing only current-year allocation
&
12 records
\\

CRM records missing only next-year allocation
&
11 records
\\

Missing-allocation records closing in 2026
&
41 of 52 records (78.8\%)
\\

2026 missing-allocation composition
&
21 both missing, 9 missing current-year and 11 missing next-year
\\

2026 missing-allocation weighted exposure
&
EUR 2,317,982.74 across 41 records
\\

2026 missing-current-year-allocation weighted risk
&
EUR 1,989,178.99 across 30 records
\\

Largest current-year allocation risk BU
&
ExEd Programs BU: EUR 1,229,613.75 weighted value across 18 records
\\

2026 Ready CRM opportunities
&
75 records with EUR 807,900.00 reconciled current-year value
\\

2026 unresolved current-year allocation pipeline
&
28 records with EUR 1,989,178.99 weighted value
\\

2026 missing-Probability exposure
&
2 records with EUR 50,000.00 estimated value; source weighted value is not reliable
\\

2026 cross-year allocation review value
&
1 record with EUR 13,500.00 calculated current-year value held for review
\\

2026 Ready CRM forecast bridge value
&
EUR 807,900.00 cleared for inclusion
\\

Business Unit label mismatch
&
Actual Sales uses shorter labels than CRM and Accrual; explicit mapping required
\\

Accrual records with missing Business Unit
&
4 rows remain Unassigned pending business validation
\\

Current-year weighted status distribution
&
85 Match, 39 Missing Current-Year Allocation and 2 Missing Probability
\\

Automatic source-record deletion
&
None performed
\\

\bottomrule

\end{tabularx}


\subsection{17.3 Completed Current-Year Status Summary}

The previously documented immediate technical step has now been completed.
A reference query was created from \texttt{Fact\_Opportunities} and named:

\begin{lstlisting}[style=powerquery]
DQ_Opportunity_CurrentYearWeightedStatus
\end{lstlisting}

In Advanced Editor, use:

\begin{lstlisting}[style=powerquery]
let
    Source = Fact_Opportunities,

    GroupedRows =
        Table.Group(
            Source,
            {"CurrentYearWeightedCheck"},
            {
                {
                    "RecordCount",
                    each Table.RowCount(_),
                    Int64.Type
                }
            }
        ),

    SortedRows =
        Table.Sort(
            GroupedRows,
            {
                {"RecordCount", Order.Descending}
            }
        )

in
    SortedRows
\end{lstlisting}

The output contains three status rows and the \texttt{RecordCount} values sum
exactly to 126:

\begin{itemize}[leftmargin=7mm,itemsep=2pt]
    \item \texttt{Match}: 85 records;
    \item \texttt{Not Testable - Missing Current-Year Allocation}: 39 records; and
    \item \texttt{Not Testable - Missing Probability}: 2 records.
\end{itemize}

No \texttt{Mismatch} status is present. The query is diagnostic only: no source
record was removed or changed. The result means that all 85 records with the
required inputs pass the current-year weighted-value reconciliation, while the
remaining 41 records are explicitly excluded from numerical testing because an
input is missing.

\subsection{17.4 Completed Allocation Validation}

The broader allocation review has now been implemented in
\texttt{DQ\_Opportunity\_AllocationValidation}. The control calculates a
two-year allocation total only when both inputs are present and classifies all
126 opportunities into six observed statuses. The result separates complete
100\% allocations from missing-input cases and from the four records whose
two-year total differs from 100\%.

No allocation value was imputed, replaced or overwritten. The control remains
diagnostic and preserves the source-aligned fact.

\subsection{17.5 Completed Allocation Status Summary}

The visual allocation distribution has now been converted into an explicit
grouped control table. A reference query from
\texttt{DQ\_Opportunity\_AllocationValidation} was created and named:

\begin{lstlisting}[style=powerquery]
DQ_Opportunity_AllocationStatusSummary
\end{lstlisting}

In Advanced Editor, use:

\begin{lstlisting}[style=powerquery]
let
    Source = DQ_Opportunity_AllocationValidation,

    GroupedRows =
        Table.Group(
            Source,
            {"AllocationValidationStatus"},
            {
                {
                    "RecordCount",
                    each Table.RowCount(_),
                    Int64.Type
                }
            }
        ),

    SortedRows =
        Table.Sort(
            GroupedRows,
            {
                {"RecordCount", Order.Descending}
            }
        )

in
    SortedRows
\end{lstlisting}

The completed output contains six status rows whose \texttt{RecordCount} values
sum exactly to 126. The confirmed counts are 70, 29, 12, 11, 3 and 1 in
descending order. This grouped query provides an auditable count table that can
be used directly in the report and later in a data-quality dashboard.

\begin{inferencebox}

\textbf{Inference -- Grouped Allocation Control}

The grouped status table reconciles exactly to the source opportunity count and
therefore confirms that the row-level allocation classification is complete.
It also provides a compact control table suitable for future dashboard KPI
cards and exception monitoring.

\end{inferencebox}

\subsection{17.6 Completed Allocation Exception Isolation}

The four complete-allocation records whose two-year total differs from 100\%
have now been isolated in:

\begin{lstlisting}[style=powerquery]
DQ_Opportunity_AllocationExceptions
\end{lstlisting}

The completed exception set contains one 110\% over-allocation and three
below-100\% totals of 50\%, 50\% and 30\%. Closing and simulated program-start
dates have been retained for each exception to support business interpretation.

\begin{inferencebox}

\textbf{Inference -- Exception Isolation}

The exceptions are few and individually reviewable. Their date patterns differ,
so a single automatic correction rule would be unsafe. The dedicated exception
query keeps the fact table unchanged while making the four records easy to
demonstrate and validate with Finance.

\end{inferencebox}

\subsection{17.7 Completed Missing Allocation Horizon Analysis}

The 52 records where one or both annual allocation inputs are missing have now
been analyzed by closing year and simulated program-start year. The purpose was
to determine whether missing allocation was concentrated in distant
opportunities rather than to fill the missing values.

Create a new reference query from
\texttt{DQ\_Opportunity\_AllocationValidation} and name it:

\begin{lstlisting}[style=powerquery]
DQ_Opportunity_MissingAllocationByYear
\end{lstlisting}

In Advanced Editor, use:

\begin{lstlisting}[style=powerquery]
let
    Source = DQ_Opportunity_AllocationValidation,

    FilteredRows =
        Table.SelectRows(
            Source,
            each
                [AllocationValidationStatus] =
                    "Review - Both Allocations Missing"
                or [AllocationValidationStatus] =
                    "Review - Missing Current-Year Allocation"
                or [AllocationValidationStatus] =
                    "Review - Missing Next-Year Allocation"
        ),

    AddedClosingYear =
        Table.AddColumn(
            FilteredRows,
            "ClosingYear",
            each
                if [ClosingDate] = null
                then null
                else Date.Year([ClosingDate]),
            Int64.Type
        ),

    AddedProgramStartYear =
        Table.AddColumn(
            AddedClosingYear,
            "ProgramStartYear",
            each
                if [SimulatedProgramStartDate] = null
                then null
                else Date.Year([SimulatedProgramStartDate]),
            Int64.Type
        ),

    GroupedRows =
        Table.Group(
            AddedProgramStartYear,
            {
                "AllocationValidationStatus",
                "ClosingYear",
                "ProgramStartYear"
            },
            {
                {
                    "RecordCount",
                    each Table.RowCount(_),
                    Int64.Type
                }
            }
        ),

    SortedRows =
        Table.Sort(
            GroupedRows,
            {
                {"AllocationValidationStatus", Order.Ascending},
                {"ClosingYear", Order.Ascending},
                {"ProgramStartYear", Order.Ascending}
            }
        )

in
    SortedRows
\end{lstlisting}

The completed grouped result reconciles exactly to 52 records. Forty-one of them
close in 2026, which represents 78.8\% of the missing-allocation population.
Eight close in 2027, one in 2028 and two in 2024--2025.

\begin{inferencebox}

\textbf{Inference -- Missing Allocation Horizon Control}

The result disproves the assumption that missing allocation is mainly caused by
long-dated pipeline opportunities. Most affected records are near-term 2026
opportunities and therefore deserve direct forecast-governance attention.

\end{inferencebox}

\subsection{17.8 Completed 2026 Missing-Allocation Detail}

A dedicated reference query was then created to isolate the 41 records with
\texttt{ClosingDate} in 2026 and at least one missing annual allocation:

\begin{lstlisting}[style=powerquery]
DQ_Opportunity_2026MissingAllocationDetail
\end{lstlisting}

The completed detail set contains:

\begin{itemize}[leftmargin=7mm,itemsep=2pt]
    \item 21 records with both annual allocations missing;
    \item 9 records missing only the current-year allocation; and
    \item 11 records missing only the next-year allocation.
\end{itemize}

The 41-row total reconciles exactly to the 2026 segment identified by the
horizon analysis.

\begin{inferencebox}

\textbf{Inference -- Near-Term Allocation Detail}

The near-term review set is now explicit and auditable. Because 41 records are
affected, record count alone is no longer sufficient for prioritization; the
next control should quantify the financial value represented by these
opportunities.

\end{inferencebox}

\subsection{17.9 Completed 2026 Missing-Allocation Financial Exposure}

The financial exposure of the 41 near-term missing-allocation records has now
been quantified using a reference query from \texttt{Fact\_Opportunities} named:

\begin{lstlisting}[style=powerquery]
DQ_Opportunity_2026MissingAllocationExposure
\end{lstlisting}

In Advanced Editor, use:

\begin{lstlisting}[style=powerquery]
let
    Source = Fact_Opportunities,

    FilteredRows =
        Table.SelectRows(
            Source,
            each
                [ClosingDate] <> null
                and Date.Year([ClosingDate]) = 2026
                and (
                    [CurrentYearAllocationPct] = null
                    or [NextYearAllocationPct] = null
                )
        ),

    AddedMissingAllocationStatus =
        Table.AddColumn(
            FilteredRows,
            "MissingAllocationStatus",
            each
                if [CurrentYearAllocationPct] = null
                    and [NextYearAllocationPct] = null then
                    "Both Allocations Missing"
                else if [CurrentYearAllocationPct] = null then
                    "Missing Current-Year Allocation"
                else
                    "Missing Next-Year Allocation",
            type text
        ),

    GroupedRows =
        Table.Group(
            AddedMissingAllocationStatus,
            {
                "MissingAllocationStatus",
                "BusinessUnitSource"
            },
            {
                {
                    "RecordCount",
                    each Table.RowCount(_),
                    Int64.Type
                },
                {
                    "TotalEstimatedValue",
                    each List.Sum(
                        List.RemoveNulls([EstimatedValue])
                    ),
                    type nullable number
                },
                {
                    "TotalWeightedEstimatedValue",
                    each List.Sum(
                        List.RemoveNulls([WeightedEstimatedValue])
                    ),
                    type nullable number
                }
            }
        ),

    SortedRows =
        Table.Sort(
            GroupedRows,
            {
                {"TotalWeightedEstimatedValue", Order.Descending},
                {"RecordCount", Order.Descending}
            }
        )

in
    SortedRows
\end{lstlisting}

The completed result reconciles exactly to 41 records. Their combined
\texttt{EstimatedValue} is EUR 4,208,699.23 and their combined
\texttt{WeightedEstimatedValue} is EUR 2,317,982.74.

\begin{inferencebox}

\textbf{Inference -- Financial Exposure Control}

The missing-allocation population carries enough financial value to affect the
forecast materially. Prioritization should therefore be driven by weighted
financial exposure in addition to record counts.

\end{inferencebox}

\subsection{17.10 Completed 2026 Current-Year Allocation Risk}

A second control was created from \texttt{Fact\_Opportunities} to retain only
2026-closing opportunities whose current-year allocation is missing:

\begin{lstlisting}[style=powerquery]
DQ_Opportunity_2026CurrentYearAllocationRisk
\end{lstlisting}

The implemented logic is:

\begin{lstlisting}[style=powerquery]
let
    Source = Fact_Opportunities,

    FilteredRows =
        Table.SelectRows(
            Source,
            each
                [ClosingDate] <> null
                and Date.Year([ClosingDate]) = 2026
                and [CurrentYearAllocationPct] = null
        ),

    GroupedRows =
        Table.Group(
            FilteredRows,
            {
                "BusinessUnitSource"
            },
            {
                {
                    "RecordCount",
                    each Table.RowCount(_),
                    Int64.Type
                },
                {
                    "TotalEstimatedValue",
                    each List.Sum(
                        List.RemoveNulls([EstimatedValue])
                    ),
                    type nullable number
                },
                {
                    "TotalWeightedEstimatedValue",
                    each List.Sum(
                        List.RemoveNulls([WeightedEstimatedValue])
                    ),
                    type nullable number
                }
            }
        ),

    SortedRows =
        Table.Sort(
            GroupedRows,
            {
                {
                    "TotalWeightedEstimatedValue",
                    Order.Descending
                }
            }
        )

in
    SortedRows
\end{lstlisting}

The completed result contains 30 records in total and EUR 1,989,178.99 of
probability-weighted pipeline whose current-year allocation is not yet available.
ExEd Programs BU contributes EUR 1,229,613.75 of this amount, Qualification
Programs BU EUR 615,625.00 and University Programs BU EUR 143,940.24.

\begin{inferencebox}

\textbf{Inference -- 2026 Current-Year Allocation Risk Control}

The unresolved current-year allocation issue is concentrated in ExEd Programs
BU and represents a material portion of the near-term CRM pipeline. These
opportunities should remain visible as an explicit at-risk bucket until Finance
or the opportunity owner confirms the annual split.

\end{inferencebox}

\subsection{17.11 Completed 2026 Opportunity Forecast Readiness}

All CRM opportunities closing in 2026 have now been classified according to
whether their current-year forecast contribution can be used, requires review,
or cannot yet be calculated reliably.

Create a new reference query from \texttt{Fact\_Opportunities} and name it:

\begin{lstlisting}[style=powerquery]
DQ_Opportunity_2026ForecastReadinessSummary
\end{lstlisting}

In Advanced Editor, use:

\begin{lstlisting}[style=powerquery]
let
    Source = Fact_Opportunities,

    FilteredRows =
        Table.SelectRows(
            Source,
            each
                [ClosingDate] <> null
                and Date.Year([ClosingDate]) = 2026
        ),

    AddedForecastReadinessStatus =
        Table.AddColumn(
            FilteredRows,
            "ForecastReadinessStatus",
            each
                if [Probability] = null then
                    "Not Ready - Missing Probability"
                else if [CurrentYearAllocationPct] = null then
                    "Not Ready - Missing Current-Year Allocation"
                else if [OpportunityQualityStatus] =
                    "Critical - Allocation Above 100%" then
                    "Review - Cross-Year Allocation Above 100%"
                else if [CurrentYearWeightedCheck] = "Match" then
                    "Ready - Reconciled Current-Year Value"
                else
                    "Review - Current-Year Reconciliation",
            type text
        ),

    GroupedRows =
        Table.Group(
            AddedForecastReadinessStatus,
            {
                "ForecastReadinessStatus",
                "BusinessUnitSource"
            },
            {
                {
                    "RecordCount",
                    each Table.RowCount(_),
                    Int64.Type
                },
                {
                    "TotalEstimatedValue",
                    each List.Sum(
                        List.RemoveNulls([EstimatedValue])
                    ),
                    type nullable number
                },
                {
                    "TotalWeightedEstimatedValue",
                    each List.Sum(
                        List.RemoveNulls([WeightedEstimatedValue])
                    ),
                    type nullable number
                },
                {
                    "TotalCalculatedCurrentYearValue",
                    each
                        let
                            Values =
                                List.RemoveNulls(
                                    [CalculatedCurrentYearWeightedValue]
                                )
                        in
                            if List.Count(Values) = 0 then
                                null
                            else
                                List.Sum(Values),
                    type nullable number
                }
            }
        ),

    SortedRows =
        Table.Sort(
            GroupedRows,
            {
                {"ForecastReadinessStatus", Order.Ascending},
                {"TotalWeightedEstimatedValue", Order.Descending}
            }
        )

in
    SortedRows
\end{lstlisting}

The completed query returns eight Business Unit rows and reconciles to
106 2026-closing opportunities:

\begin{itemize}[leftmargin=7mm,itemsep=2pt]
    \item 75 Ready records with EUR 807,900.00 of reconciled current-year value;
    \item 28 Not Ready records with EUR 1,989,178.99 of weighted pipeline because
    current-year allocation is missing;
    \item 2 Not Ready records with EUR 50,000.00 of estimated value because
    Probability is missing; and
    \item 1 Review record with EUR 13,500.00 of calculated current-year value because
    its two-year allocation totals 110\%.
\end{itemize}

\begin{inferencebox}

\textbf{Inference -- Forecast Readiness Control}

The query establishes a clear inclusion boundary for CRM forecast use. Ready
records have both the required business inputs and a reconciled current-year
calculation. Not Ready and Review records remain visible but are not silently
folded into the forecast. This preserves both financial transparency and data-
quality accountability.

\end{inferencebox}

\subsection{17.12 Completed Forecast Readiness Status Summary}

The Business Unit detail was summarized into one row per forecast-readiness
status using:

\begin{lstlisting}[style=powerquery]
DQ_Opportunity_2026ForecastReadinessStatusSummary
\end{lstlisting}

The implemented Power Query logic is:

\begin{lstlisting}[style=powerquery]
let
    Source = DQ_Opportunity_2026ForecastReadinessSummary,

    GroupedRows =
        Table.Group(
            Source,
            {
                "ForecastReadinessStatus"
            },
            {
                {
                    "RecordCount",
                    each List.Sum([RecordCount]),
                    Int64.Type
                },
                {
                    "TotalEstimatedValue",
                    each List.Sum(
                        List.RemoveNulls([TotalEstimatedValue])
                    ),
                    type nullable number
                },
                {
                    "TotalWeightedEstimatedValue",
                    each List.Sum(
                        List.RemoveNulls([TotalWeightedEstimatedValue])
                    ),
                    type nullable number
                },
                {
                    "TotalCalculatedCurrentYearValue",
                    each
                        let
                            Values =
                                List.RemoveNulls(
                                    [TotalCalculatedCurrentYearValue]
                                )
                        in
                            if List.Count(Values) = 0 then
                                null
                            else
                                List.Sum(Values),
                    type nullable number
                }
            }
        ),

    SortedRows =
        Table.Sort(
            GroupedRows,
            {
                {"RecordCount", Order.Descending}
            }
        )

in
    SortedRows
\end{lstlisting}

The four status rows reconcile exactly to 106 records. The Ready status contains
EUR 807,900.00 of reconciled current-year value. The cross-year allocation
review contains a further EUR 13,500.00 of calculated current-year value that is
not yet cleared for inclusion.

\begin{inferencebox}

\textbf{Inference -- Readiness Status Summary Control}

The grouped query converts the detailed control into a compact executive view.
It makes a crucial distinction between forecast value that is currently usable,
value that is calculated but under review, and financially material pipeline that
cannot yet be allocated into 2026.

\end{inferencebox}

\subsection{17.13 Completed 2026 CRM Forecast Bridge}

The four-row forecast bridge has now been created and assigns an explicit
treatment to each readiness status without mixing values that have different
financial meanings.

Create a new reference query from
\texttt{DQ\_Opportunity\_2026ForecastReadinessStatusSummary} and name it:

\begin{lstlisting}[style=powerquery]
DQ_Opportunity_2026ForecastBridge
\end{lstlisting}

In Advanced Editor, use:

\begin{lstlisting}[style=powerquery]
let
    Source =
        DQ_Opportunity_2026ForecastReadinessStatusSummary,

    AddedForecastTreatment =
        Table.AddColumn(
            Source,
            "ForecastTreatment",
            each
                if [ForecastReadinessStatus] =
                    "Ready - Reconciled Current-Year Value" then
                    "Include - Ready Current-Year Forecast"

                else if [ForecastReadinessStatus] =
                    "Review - Cross-Year Allocation Above 100%" then
                    "Hold - Current-Year Value Under Review"

                else if [ForecastReadinessStatus] =
                    "Not Ready - Missing Current-Year Allocation" then
                    "At Risk - Weighted Pipeline Not Allocated to 2026"

                else
                    "At Risk - Estimated Pipeline Missing Probability",
            type text
        ),

    AddedReadyCurrentYearForecast =
        Table.AddColumn(
            AddedForecastTreatment,
            "ReadyCurrentYearForecastValue",
            each
                if [ForecastReadinessStatus] =
                    "Ready - Reconciled Current-Year Value"
                then [TotalCalculatedCurrentYearValue]
                else null,
            type nullable number
        ),

    AddedReviewCurrentYearValue =
        Table.AddColumn(
            AddedReadyCurrentYearForecast,
            "ReviewCurrentYearValue",
            each
                if [ForecastReadinessStatus] =
                    "Review - Cross-Year Allocation Above 100%"
                then [TotalCalculatedCurrentYearValue]
                else null,
            type nullable number
        ),

    AddedAtRiskWeightedPipeline =
        Table.AddColumn(
            AddedReviewCurrentYearValue,
            "AtRiskWeightedPipeline",
            each
                if [ForecastReadinessStatus] =
                    "Not Ready - Missing Current-Year Allocation"
                then [TotalWeightedEstimatedValue]
                else null,
            type nullable number
        ),

    AddedMissingProbabilityEstimatedValue =
        Table.AddColumn(
            AddedAtRiskWeightedPipeline,
            "MissingProbabilityEstimatedValue",
            each
                if [ForecastReadinessStatus] =
                    "Not Ready - Missing Probability"
                then [TotalEstimatedValue]
                else null,
            type nullable number
        ),

    SelectedColumns =
        Table.SelectColumns(
            AddedMissingProbabilityEstimatedValue,
            {
                "ForecastReadinessStatus",
                "ForecastTreatment",
                "RecordCount",
                "ReadyCurrentYearForecastValue",
                "ReviewCurrentYearValue",
                "AtRiskWeightedPipeline",
                "MissingProbabilityEstimatedValue"
            }
        )

in
    SelectedColumns
\end{lstlisting}

The expected control values are:

\begin{itemize}[leftmargin=7mm,itemsep=2pt]
    \item Ready current-year forecast: EUR 807,900.00;
    \item current-year value held for review: EUR 13,500.00;
    \item at-risk weighted pipeline with missing current-year allocation:
    EUR 1,989,178.99; and
    \item estimated pipeline with missing Probability: EUR 50,000.00.
\end{itemize}

These four values remain in separate columns. They are not four additive
components of the same accounting measure. The completed bridge therefore
communicates forecast treatment clearly before Business Unit harmonization and
semantic-model implementation.

\begin{inferencebox}

\textbf{Inference -- CRM Forecast Bridge Control}

The bridge creates a governed boundary between forecast inclusion and unresolved
pipeline. It prevents financially different concepts from being collapsed into
one number and gives the future semantic model four explicit treatments rather
than forcing missing information into zero-valued forecast rows.

\end{inferencebox}

\subsection{17.14 Completed Cross-Source Business Unit Vocabulary Audit}

A dedicated control query was created to inventory the Business Unit values used
by \texttt{Fact\_ActualSales}, \texttt{Fact\_ProjectAccrual} and
\texttt{Fact\_Opportunities}:

\begin{lstlisting}[style=powerquery]
DQ_BusinessUnit_SourceValues
\end{lstlisting}

The audit confirms:

\begin{itemize}[leftmargin=7mm,itemsep=2pt]
    \item Actual Sales: 346 \texttt{ExEd}, 120 \texttt{Qualification Programs}
    and 329 \texttt{University Programs};
    \item CRM Opportunities: 92 \texttt{ExEd Programs BU},
    7 \texttt{Qualification Programs BU} and
    27 \texttt{University Programs BU};
    \item Project Accrual: 162 \texttt{ExEd Programs BU},
    43 \texttt{Qualification Programs BU},
    104 \texttt{University Programs BU} and 4 missing Business Unit values.
\end{itemize}

\begin{inferencebox}

\textbf{Inference -- Business Unit Source Audit}

The three datasets describe the same three populated Business Units, but Actual
Sales uses a shorter naming convention. CRM and Accrual already agree. The four
missing Accrual Business Units must remain unassigned until validated. An
explicit shared dimension is therefore safer and more auditable than replacing
labels ad hoc inside each source.

\end{inferencebox}

\subsection{17.15 Completed DimBusinessUnit}

Create a new query group named:

\begin{lstlisting}[style=powerquery]
04_Dimensions
\end{lstlisting}

Then create a Blank Query named:

\begin{lstlisting}[style=powerquery]
DimBusinessUnit
\end{lstlisting}

Use the following Power Query code:

\begin{lstlisting}[style=powerquery]
let
    Source =
        #table(
            type table [
                BusinessUnitKey = Int64.Type,
                BusinessUnitName = text,
                BusinessUnitShortName = text,
                BusinessUnitSortOrder = Int64.Type
            ],
            {
                {
                    1,
                    "ExEd Programs BU",
                    "ExEd",
                    1
                },
                {
                    2,
                    "Qualification Programs BU",
                    "Qualification Programs",
                    2
                },
                {
                    3,
                    "University Programs BU",
                    "University Programs",
                    3
                },
                {
                    0,
                    "Unassigned / Missing BU",
                    "Unassigned",
                    99
                }
            }
        )

in
    Source
\end{lstlisting}

The completed dimension contains exactly four rows. The fourth row is a technical
member for records whose Business Unit is genuinely missing; it does not infer
or fabricate a business classification. All four dimension columns are fully
populated and the four keys are unique.

\begin{inferencebox}

\textbf{Inference -- Canonical Business Unit Dimension}

The Business Unit vocabulary is now governed independently from the source
systems. Keys 1--3 represent the three actual Business Units and key 0 provides
a controlled destination for genuinely unassigned records. This design avoids
guessing the Business Unit of incomplete accrual rows while enabling consistent
cross-source filtering.

\end{inferencebox}

\subsection{17.16 Completed Actual Sales Business Unit Key Mapping}

A \texttt{BusinessUnitKey} column was added to
\texttt{Fact\_ActualSales}. The original \texttt{BusinessUnitSource} was not
removed.

The implemented mapping is:

\begin{lstlisting}[style=powerquery]
if [BusinessUnitSource] = "ExEd" then
    1
else if [BusinessUnitSource] = "Qualification Programs" then
    2
else if [BusinessUnitSource] = "University Programs" then
    3
else
    0
\end{lstlisting}

The field was typed as \texttt{Int64.Type}. The result contains 795 rows,
100\% valid keys, 0 errors, 0 empty keys and 3 distinct key values. The
distribution is 346 rows on key 1, 120 rows on key 2 and 329 rows on key 3.
No Actual Sales row maps to key 0.

\begin{inferencebox}

\textbf{Inference -- Actual Sales BU-Key Integration}

The shorter Actual Sales source labels no longer prevent consistent Business
Unit modelling. Every row now carries a stable relationship key, but the
original source text remains available for audit and lineage. The mapping also
reconciles exactly to the earlier Business Unit source audit.

\end{inferencebox}

\subsection{17.17 Completed Project Accrual Business Unit Key}

The \texttt{Fact\_ProjectAccrual} mapping has now been completed. A Custom Column
named \texttt{BusinessUnitKey} was added:

\begin{lstlisting}[style=powerquery]
BusinessUnitKey
\end{lstlisting}

Use:

\begin{lstlisting}[style=powerquery]
if [BusinessUnitSource] = null
    or Text.Trim(Text.From([BusinessUnitSource])) = "" then
    0
else if [BusinessUnitSource] = "ExEd Programs BU" then
    1
else if [BusinessUnitSource] = "Qualification Programs BU" then
    2
else if [BusinessUnitSource] = "University Programs BU" then
    3
else
    null
\end{lstlisting}

Then set the column type to Whole Number / \texttt{Int64.Type}.

Based on the completed source audit, the expected distribution is:

\begin{itemize}[leftmargin=7mm,itemsep=2pt]
    \item key 1: 162 records;
    \item key 2: 43 records;
    \item key 3: 104 records; and
    \item key 0: 4 records with missing Business Unit.
\end{itemize}

The completed result remains at 313 rows and the key column is 100\% valid,
with 0 errors, 0 empty values and 4 distinct keys. The earlier source-value
audit reconciles those keys to 162 / 43 / 104 / 4 records on keys
1 / 2 / 3 / 0 respectively. The \texttt{else null} branch remains
deliberately in the transformation so an unexpected future BU label will
surface as an unmapped value rather than being silently assigned.

\begin{inferencebox}

\textbf{Inference -- Project Accrual BU-Key Integration}

Project Accrual is now dimension-ready. The explicit key-0 treatment preserves
the four source records with missing Business Unit as an auditable exception
while still allowing them to participate in the model under the Unassigned
dimension member.

\end{inferencebox}

\subsection{17.18 Completed CRM Opportunity Business Unit Key}

The same canonical mapping was then applied to
\texttt{Fact\_Opportunities}:

\begin{lstlisting}[style=powerquery]
if [BusinessUnitSource] = null
    or Text.Trim(Text.From([BusinessUnitSource])) = "" then
    0
else if [BusinessUnitSource] = "ExEd Programs BU" then
    1
else if [BusinessUnitSource] = "Qualification Programs BU" then
    2
else if [BusinessUnitSource] = "University Programs BU" then
    3
else
    null
\end{lstlisting}

The key was typed as \texttt{Int64.Type}. The completed profile contains
126 rows, 100\% valid values, 0 errors, 0 empty values and 3 distinct keys.
The distribution reconciles to 92 / 7 / 27 records on keys 1 / 2 / 3, with
no records on key 0.

\begin{inferencebox}

\textbf{Inference -- CRM Opportunity BU-Key Integration}

All three analytical facts now carry the same stable Business Unit relationship
key. Source-specific Business Unit text remains preserved for lineage, while
the semantic model can use one shared dimension without relying on text joins.

\end{inferencebox}

\subsection{17.19 Completed Model-Load Cleanup and Business Unit Relationships}

Before relationship validation, loading was disabled for all staging and
data-quality queries. They remain in Power Query with their applied steps and
continue to support transformation, reconciliation and audit logic.

The visible model now contains only:

\begin{itemize}[leftmargin=7mm,itemsep=2pt]
    \item \texttt{DimBusinessUnit};
    \item \texttt{Fact\_ActualSales};
    \item \texttt{Fact\_ProjectAccrual}; and
    \item \texttt{Fact\_Opportunities}.
\end{itemize}

Three relationships have been created on \texttt{BusinessUnitKey}. In each
case, \texttt{DimBusinessUnit} is on the one side and the fact table is on the
many side. The relationships are active and use Single cross-filter direction.
The selected Actual Sales relationship is explicitly shown by Power BI as
\texttt{Many to one (*:1)} from fact to dimension, which is the same star-schema
cardinality viewed from the fact side.

\begin{inferencebox}

\textbf{Inference -- Business Unit Relationship Model}

The Business Unit model is now structurally clean: there are no loaded staging
or DQ helper tables and no fact-to-fact Business Unit relationships. One
dimension filters all three facts through stable keys. This limits ambiguous
filter paths and establishes the correct star-schema direction before time and
project dimensions are added.

\end{inferencebox}

\subsection{17.20 Completed Business Unit Filter-Propagation Validation}

The three Business Unit relationships have now been validated with row-count
measures to confirm that filters propagate exactly as intended.

Create the following three measures:

\begin{lstlisting}[style=powerquery]
Actual Row Count =
COALESCE(
    COUNTROWS(Fact_ActualSales),
    0
)
\end{lstlisting}

\begin{lstlisting}[style=powerquery]
Accrual Row Count =
COALESCE(
    COUNTROWS(Fact_ProjectAccrual),
    0
)
\end{lstlisting}

\begin{lstlisting}[style=powerquery]
Opportunity Row Count =
COALESCE(
    COUNTROWS(Fact_Opportunities),
    0
)
\end{lstlisting}

Create a Table or Matrix visual with
\texttt{DimBusinessUnit[BusinessUnitName]} on rows and the three measures as
values. The expected validation result is:

\begin{tabularx}{\textwidth}{
    >{\raggedright\arraybackslash}X
    >{\centering\arraybackslash}p{2.3cm}
    >{\centering\arraybackslash}p{2.3cm}
    >{\centering\arraybackslash}p{2.3cm}
}

\toprule

Business Unit
&
Actual Rows
&
Accrual Rows
&
Opportunity Rows
\\

\midrule

ExEd Programs BU
&
346
&
162
&
92
\\

Qualification Programs BU
&
120
&
43
&
7
\\

University Programs BU
&
329
&
104
&
27
\\

Unassigned / Missing BU
&
0
&
4
&
0
\\

\midrule

Total
&
795
&
313
&
126
\\

\bottomrule

\end{tabularx}

The completed visual reconciles exactly to the expected values shown above.
The totals are 795 Actual rows, 313 Accrual rows and 126 Opportunity rows.
The Unassigned member contains 0 Actual rows, 4 Accrual rows and 0 Opportunity
rows, which confirms that the four missing-BU accrual records are isolated
correctly and no other fact record is misclassified.

\begin{inferencebox}

\textbf{Inference -- Business Unit Filter Propagation}

The row-count test validates the Business Unit relationships functionally, not
only structurally. Each Business Unit produces the expected row counts in all
three facts, and the grand totals reconcile exactly to the fact-table row
counts. This confirms that \texttt{DimBusinessUnit} can be used safely as a
shared slicer across Actual Sales, Project Accrual and CRM Opportunities.

\end{inferencebox}

\subsection{17.21 Completed 2026 DimDate Creation}

A new query named \texttt{DimDate} was created with a complete daily calendar
from 1 January through 31 December 2026.

The resulting table contains 365 rows and seven columns:

\begin{itemize}[leftmargin=7mm,itemsep=2pt]
    \item \texttt{Date};
    \item \texttt{Year};
    \item \texttt{MonthNumber};
    \item \texttt{MonthName};
    \item \texttt{YearMonth};
    \item \texttt{MonthStart}; and
    \item \texttt{Quarter}.
\end{itemize}

The profile confirms 365 unique dates, zero errors, zero empty dates,
12 MonthNumber values, 12 MonthName values, 12 YearMonth values,
12 MonthStart values and 4 Quarter values.

\begin{warningbox}

The current Power Query layout places \texttt{DimDate} under
\texttt{01\_Staging}. The data itself is correct, but the query should be moved
to \texttt{04\_Dimensions} before Date relationships are created. This is a
query-organization correction only and does not require changing the M code.

\end{warningbox}

\begin{inferencebox}

\textbf{Inference -- Date Dimension Readiness}

The calendar has a unique daily key and complete 2026 coverage, so it is ready
to act as the one side of Date relationships. Actual Sales covers January--
August through \texttt{MonthStart}, while Project Accrual covers September--
December through \texttt{ForecastMonth}; a shared Date dimension will therefore
support a continuous full-year 2026 view without appending the two facts.

\end{inferencebox}

\subsection{17.22 Completed DimDate Placement and Date Relationships}

\texttt{DimDate} is now in the governed dimension layer. The Date relationships
were then created in Model View:

\begin{tabularx}{\textwidth}{
    >{\ttfamily\raggedright\arraybackslash}p{5.1cm}
    >{\ttfamily\raggedright\arraybackslash}p{5.1cm}
    >{\centering\arraybackslash}p{1.7cm}
    >{\centering\arraybackslash}X
}

\toprule

Dimension Side
&
Fact Side
&
Cardinality
&
Filter
\\

\midrule

DimDate[Date]
&
Fact\_ActualSales[MonthStart]
&
1:*
&
Single
\\

DimDate[Date]
&
Fact\_ProjectAccrual[ForecastMonth]
&
1:*
&
Single
\\

\bottomrule

\end{tabularx}

Both relationships are active. Power BI displays them from the fact-table
perspective as \texttt{Many to one (*:1)}, with \texttt{DimDate} on the one
side. Both use \texttt{Cross-filter direction = Single}.

\texttt{Fact\_Opportunities} remains intentionally disconnected from
\texttt{DimDate}. CRM contains multiple date roles
(\texttt{EstimatedClosingDate}, \texttt{ClosingDate} and
\texttt{SimulatedProgramStartDate}) and records outside 2026. The CRM Date
strategy will therefore be designed separately.

\begin{inferencebox}

\textbf{Inference -- Date Relationship Structure}

The semantic model now has two clean conformed dimensions:
\texttt{DimBusinessUnit} and \texttt{DimDate}. Date filters flow from one
calendar to Actual Sales and Project Accrual without creating fact-to-fact
relationships or bidirectional ambiguity. CRM remains deliberately outside this
Date path until its multiple date roles are modelled explicitly.

\end{inferencebox}

\subsection{17.23 Completed Date Filter-Propagation Validation}

The existing row-count measures were reused with
\texttt{DimDate[MonthName]} to verify the Date relationships operationally.

The completed month-level result is:

\begin{tabularx}{\textwidth}{
    >{\raggedright\arraybackslash}p{4.2cm}
    >{\centering\arraybackslash}p{3cm}
    >{\centering\arraybackslash}p{3cm}
}

\toprule

Month
&
Actual Row Count
&
Accrual Row Count
\\

\midrule

January & 105 & 0 \\
February & 96 & 0 \\
March & 112 & 0 \\
April & 113 & 0 \\
May & 115 & 0 \\
June & 110 & 0 \\
July & 50 & 0 \\
August & 94 & 0 \\
September & 0 & 87 \\
October & 0 & 76 \\
November & 0 & 84 \\
December & 0 & 66 \\

\midrule

Total & 795 & 313 \\

\bottomrule

\end{tabularx}

The January--August Actual counts sum exactly to 795, while the September--
December Accrual counts sum exactly to 313. No Accrual row appears in the
January--August period and no Actual row appears in the September--December
period.

\begin{inferencebox}

\textbf{Inference -- Date Filter Propagation}

The Date relationships are now validated functionally, not only visually.
The shared calendar reproduces the known source horizons exactly:
Actual Sales supplies the first eight months of 2026 and Project Accrual
supplies the last four months. This proves that a single Date axis can support a
continuous full-year 2026 view without appending the two facts or introducing
double-counting through the relationship layer.

\end{inferencebox}

\subsection{17.24 Completed Month Sorting and Base Financial Measures}

The month-order control has now been completed. \texttt{DimDate[MonthName]} is
displayed chronologically from January through December using
\texttt{MonthNumber} as its sort column.

The first three financial measures have also been created:

\begin{lstlisting}[style=powerquery]
Actual Net Sales =
SUM(Fact_ActualSales[ActualNetSales])
\end{lstlisting}

\begin{lstlisting}[style=powerquery]
Project Accrual Forecast =
SUM(Fact_ProjectAccrual[ProjectAccrualForecast])
\end{lstlisting}

\begin{lstlisting}[style=powerquery]
Base Full-Year Forecast =
COALESCE([Actual Net Sales], 0)
    + COALESCE([Project Accrual Forecast], 0)
\end{lstlisting}

The completed month-level visual shows:

\begin{tabularx}{\textwidth}{
    >{\raggedright\arraybackslash}p{2.6cm}
    >{\raggedleft\arraybackslash}p{3.4cm}
    >{\raggedleft\arraybackslash}p{3.4cm}
    >{\raggedleft\arraybackslash}X
}

\toprule

Month
&
Actual Net Sales
&
Project Accrual Forecast
&
Base Full-Year Forecast
\\

\midrule

January   & EUR 1,727,161.80 & -- & EUR 1,727,161.80 \\
February  & EUR 1,682,476.45 & -- & EUR 1,682,476.45 \\
March     & EUR 1,977,142.41 & -- & EUR 1,977,142.41 \\
April     & EUR 2,514,361.93 & -- & EUR 2,514,361.93 \\
May       & EUR 2,710,177.12 & -- & EUR 2,710,177.12 \\
June      & EUR 1,736,234.71 & -- & EUR 1,736,234.71 \\
July      & EUR 429,048.56   & -- & EUR 429,048.56 \\
August    & EUR 1,858,464.84 & -- & EUR 1,858,464.84 \\
September & -- & EUR 1,866,125.46 & EUR 1,866,125.46 \\
October   & -- & EUR 1,706,279.10 & EUR 1,706,279.10 \\
November  & -- & EUR 1,807,452.49 & EUR 1,807,452.49 \\
December  & -- & EUR 1,236,168.73 & EUR 1,236,168.73 \\

\midrule

\textbf{Total}
&
\textbf{EUR 14,635,067.81}
&
\textbf{EUR 6,616,025.78}
&
\textbf{EUR 21,251,093.59}
\\

\bottomrule

\end{tabularx}

The grand totals reconcile exactly to the source-level controls already
documented in the report:

\begin{lstlisting}[style=powerquery]
EUR 14,635,067.81
+ EUR 6,616,025.78
= EUR 21,251,093.59
\end{lstlisting}

The displayed monthly values are formatted to two decimals independently.
Accordingly, a one-cent difference can arise when manually summing rounded
monthly display values versus evaluating the DAX grand total over the underlying
source precision. The authoritative control is the Power BI grand total, which
matches the previously reconciled source totals.

\begin{warningbox}

\textbf{Source-Aligned Accrual Basis}

The EUR 6,616,025.78 Project Accrual Forecast remains the source-aligned amount.
The earlier duplicate analysis identified EUR 175,318.63 of potential duplicate
overstatement, but no accrual rows have been removed because Finance validation
is still pending. Therefore, EUR 21,251,093.59 is the controlled
\emph{source-aligned base forecast}, not a Finance-approved duplicate-adjusted
forecast.

\end{warningbox}

\begin{inferencebox}

\textbf{Inference -- Base Full-Year Forecast}

The first two financial layers are now operational on one shared 2026 calendar.
Actual Net Sales supplies the realized January--August value and Project Accrual
supplies the September--December ongoing-project forecast. The two measures
combine into a continuous 12-month base forecast of EUR 21,251,093.59 before
CRM opportunity contribution is added. This validates both the relationship
design and the financial aggregation logic at the annual and monthly levels.

\end{inferencebox}

\subsection{17.25 Completed Business Unit Financial Validation}

The three base financial measures were validated through
\texttt{DimBusinessUnit[BusinessUnitName]}. The completed result is:

\begin{tabularx}{\textwidth}{
    >{\raggedright\arraybackslash}p{4.2cm}
    >{\raggedleft\arraybackslash}p{3.2cm}
    >{\raggedleft\arraybackslash}p{3.2cm}
    >{\raggedleft\arraybackslash}X
}

\toprule

Business Unit
&
Actual Net Sales
&
Project Accrual Forecast
&
Base Full-Year Forecast
\\

\midrule

ExEd Programs BU
&
EUR 6,044,545.34
&
EUR 3,009,546.88
&
EUR 9,054,092.22
\\

Qualification Programs BU
&
EUR 4,866,104.98
&
EUR 1,772,854.87
&
EUR 6,638,959.85
\\

University Programs BU
&
EUR 3,724,417.49
&
EUR 1,791,074.03
&
EUR 5,515,491.52
\\

Unassigned / Missing BU
&
EUR 0.00
&
EUR 42,550.00
&
EUR 42,550.00
\\

\midrule

\textbf{Total}
&
\textbf{EUR 14,635,067.81}
&
\textbf{EUR 6,616,025.78}
&
\textbf{EUR 21,251,093.59}
\\

\bottomrule

\end{tabularx}

The Business Unit rows reconcile exactly to the previously validated annual
totals. The explicit Unassigned member contains no Actual Net Sales and exactly
EUR 42,550.00 of Project Accrual Forecast, matching the four known missing-BU
accrual records.

\begin{inferencebox}

\textbf{Inference -- Business Unit Financial Validation}

The Business Unit relationship layer is now validated for both record counts
and financial values. The same dimension that reconciled 795 Actual, 313
Accrual and 126 CRM rows also propagates the financial measures without loss or
misallocation. Keeping EUR 42,550.00 in the explicit Unassigned member preserves
the source exception instead of forcing it into a real Business Unit.

\end{inferencebox}

\subsection{17.26 Completed Date-Table Designation}

\texttt{DimDate} has now been formally marked as the Power BI model Date table
using:

\begin{lstlisting}[style=powerquery]
Table tools
-> Mark as date table
-> Date column: DimDate[Date]
\end{lstlisting}

This completes the model metadata control for the shared 2026 calendar.

\begin{inferencebox}

\textbf{Inference -- Date-Table Designation}

The calendar is no longer only a technically valid dimension with active
relationships; it is also explicitly designated as the model Date table. This
provides a clear semantic time backbone for current and future time-intelligence
measures while preserving the deliberate decision not to force CRM onto a
single Date role yet.

\end{inferencebox}

\subsection{17.27 Completed Ready CRM Measures and Forecast Inclusion}

Two governed CRM measures were created to expose only the technically cleared
2026 current-year CRM contribution.

The Ready opportunity count is:

\begin{lstlisting}[style=powerquery]
CRM Ready Opportunity Count =
CALCULATE(
    COUNTROWS(Fact_Opportunities),
    FILTER(
        Fact_Opportunities,
        YEAR(Fact_Opportunities[ClosingDate]) = 2026
            && Fact_Opportunities[CurrentYearWeightedCheck] = "Match"
            && Fact_Opportunities[OpportunityQualityStatus]
                <> "Critical - Allocation Above 100%"
    )
)
\end{lstlisting}

The Ready current-year value is:

\begin{lstlisting}[style=powerquery]
CRM Ready Current-Year Value =
CALCULATE(
    SUM(Fact_Opportunities[CurrentYearWeightedEstimatedValue]),
    FILTER(
        Fact_Opportunities,
        YEAR(Fact_Opportunities[ClosingDate]) = 2026
            && Fact_Opportunities[CurrentYearWeightedCheck] = "Match"
            && Fact_Opportunities[OpportunityQualityStatus]
                <> "Critical - Allocation Above 100%"
    )
)
\end{lstlisting}

The completed Business Unit validation is:

\begin{tabularx}{\textwidth}{
    >{\raggedright\arraybackslash}p{5.0cm}
    >{\centering\arraybackslash}p{2.6cm}
    >{\raggedleft\arraybackslash}X
}

\toprule

Business Unit
&
Ready Opportunity Count
&
Ready Current-Year Value
\\

\midrule

ExEd Programs BU
&
60
&
EUR 527,877.50
\\

Qualification Programs BU
&
2
&
EUR 35,250.00
\\

University Programs BU
&
13
&
EUR 244,772.50
\\

Unassigned / Missing BU
&
0
&
EUR 0.00
\\

\midrule

\textbf{Total}
&
\textbf{75}
&
\textbf{EUR 807,900.00}
\\

\bottomrule

\end{tabularx}

The Ready CRM contribution was then added to the source-aligned base forecast:

\begin{lstlisting}[style=powerquery]
Forecast Including Ready CRM =
COALESCE([Base Full-Year Forecast], 0)
    + COALESCE([CRM Ready Current-Year Value], 0)
\end{lstlisting}

The completed Business Unit result is:

\begin{tabularx}{\textwidth}{
    >{\raggedright\arraybackslash}p{4.2cm}
    >{\raggedleft\arraybackslash}p{3.2cm}
    >{\raggedleft\arraybackslash}p{3.2cm}
    >{\raggedleft\arraybackslash}X
}

\toprule

Business Unit
&
Base Full-Year Forecast
&
Ready CRM Value
&
Forecast Including Ready CRM
\\

\midrule

ExEd Programs BU
&
EUR 9,054,092.22
&
EUR 527,877.50
&
EUR 9,581,969.72
\\

Qualification Programs BU
&
EUR 6,638,959.85
&
EUR 35,250.00
&
EUR 6,674,209.85
\\

University Programs BU
&
EUR 5,515,491.52
&
EUR 244,772.50
&
EUR 5,760,264.02
\\

Unassigned / Missing BU
&
EUR 42,550.00
&
EUR 0.00
&
EUR 42,550.00
\\

\midrule

\textbf{Total}
&
\textbf{EUR 21,251,093.59}
&
\textbf{EUR 807,900.00}
&
\textbf{EUR 22,058,993.59}
\\

\bottomrule

\end{tabularx}

\begin{warningbox}

\textbf{Annual / Business Unit CRM Inclusion Only}

The Ready CRM measure is intentionally validated at annual and Business Unit
level. It must not yet be interpreted as a month-level CRM forecast through
\texttt{DimDate}, because \texttt{Fact\_Opportunities} has multiple candidate
date roles and is deliberately not connected to the 2026 Date dimension at this
stage.

\end{warningbox}

\begin{inferencebox}

\textbf{Inference -- Governed Forecast Including Ready CRM}

The source-aligned base forecast of EUR 21,251,093.59 can now be combined with
EUR 807,900.00 of technically cleared current-year CRM value without mixing in
unresolved pipeline. The current governed annual forecast is therefore
EUR 22,058,993.59. The CRM contribution is also fully reconciled by Business
Unit: ExEd contributes EUR 527,877.50, Qualification Programs EUR 35,250.00 and
University Programs EUR 244,772.50.

\end{inferencebox}

\subsection{17.28 Immediate Next Step -- Semantic CRM Governance Measures}

The next step is to surface the three remaining CRM forecast-treatment buckets
as semantic-model measures while keeping them separate from the Ready forecast.

Create:

\begin{lstlisting}[style=powerquery]
CRM Review Current-Year Value =
CALCULATE(
    SUM(Fact_Opportunities[CurrentYearWeightedEstimatedValue]),
    FILTER(
        Fact_Opportunities,
        YEAR(Fact_Opportunities[ClosingDate]) = 2026
            && NOT ISBLANK(Fact_Opportunities[Probability])
            && NOT ISBLANK(Fact_Opportunities[CurrentYearAllocationPct])
            && NOT ISBLANK(Fact_Opportunities[NextYearAllocationPct])
            && Fact_Opportunities[CurrentYearAllocationPct]
                + Fact_Opportunities[NextYearAllocationPct] > 100
    )
)
\end{lstlisting}

Expected result:

\begin{lstlisting}[style=powerquery]
EUR 13,500.00
\end{lstlisting}

Then create:

\begin{lstlisting}[style=powerquery]
CRM At-Risk Weighted Pipeline =
CALCULATE(
    SUM(Fact_Opportunities[WeightedEstimatedValue]),
    FILTER(
        Fact_Opportunities,
        YEAR(Fact_Opportunities[ClosingDate]) = 2026
            && NOT ISBLANK(Fact_Opportunities[Probability])
            && ISBLANK(Fact_Opportunities[CurrentYearAllocationPct])
    )
)
\end{lstlisting}

Expected result:

\begin{lstlisting}[style=powerquery]
EUR 1,989,178.99
\end{lstlisting}

Then create:

\begin{lstlisting}[style=powerquery]
CRM Missing Probability Estimated Value =
CALCULATE(
    SUM(Fact_Opportunities[EstimatedValue]),
    FILTER(
        Fact_Opportunities,
        YEAR(Fact_Opportunities[ClosingDate]) = 2026
            && ISBLANK(Fact_Opportunities[Probability])
    )
)
\end{lstlisting}

Expected result:

\begin{lstlisting}[style=powerquery]
EUR 50,000.00
\end{lstlisting}

Optionally create a review-approved scenario:

\begin{lstlisting}[style=powerquery]
Forecast If Review Approved =
[Forecast Including Ready CRM]
    + [CRM Review Current-Year Value]
\end{lstlisting}

Expected result:

\begin{lstlisting}[style=powerquery]
EUR 22,072,493.59
\end{lstlisting}

Use separate Card visuals for the governed values. Do not add the At-Risk
weighted pipeline or the Missing-Probability exposure into
\texttt{Forecast Including Ready CRM}; they remain unresolved pipeline, not
cleared current-year forecast.

\begin{inferencebox}

\textbf{Inference -- Why the Risk Buckets Stay Separate}

The next semantic layer should communicate forecast confidence rather than
collapse all pipeline into one number. EUR 807,900.00 is cleared for inclusion,
EUR 13,500.00 is a calculable amount held for review, EUR 1,989,178.99 is
probability-weighted pipeline whose 2026 allocation is unknown, and
EUR 50,000.00 is estimated pipeline whose Probability is missing. Keeping these
buckets separate prevents false precision and creates an interview-ready
forecast bridge.

\end{inferencebox}

\subsection{17.29 Completed Semantic CRM Governance Measures and KPI Presentation}

The three CRM governance measures documented previously as the immediate next
semantic-model control have now been created and validated in the report model.
The completed measures are:

\begin{lstlisting}[style=powerquery]
CRM Review Current-Year Value =
CALCULATE(
    SUM(Fact_Opportunities[CurrentYearWeightedEstimatedValue]),
    FILTER(
        Fact_Opportunities,
        YEAR(Fact_Opportunities[ClosingDate]) = 2026
            && NOT ISBLANK(Fact_Opportunities[Probability])
            && NOT ISBLANK(Fact_Opportunities[CurrentYearAllocationPct])
            && NOT ISBLANK(Fact_Opportunities[NextYearAllocationPct])
            && Fact_Opportunities[CurrentYearAllocationPct]
                + Fact_Opportunities[NextYearAllocationPct] > 100
    )
)
\end{lstlisting}

\begin{lstlisting}[style=powerquery]
CRM At-Risk Weighted Pipeline =
CALCULATE(
    SUM(Fact_Opportunities[WeightedEstimatedValue]),
    FILTER(
        Fact_Opportunities,
        YEAR(Fact_Opportunities[ClosingDate]) = 2026
            && NOT ISBLANK(Fact_Opportunities[Probability])
            && ISBLANK(Fact_Opportunities[CurrentYearAllocationPct])
    )
)
\end{lstlisting}

\begin{lstlisting}[style=powerquery]
CRM Missing Probability Estimated Value =
CALCULATE(
    SUM(Fact_Opportunities[EstimatedValue]),
    FILTER(
        Fact_Opportunities,
        YEAR(Fact_Opportunities[ClosingDate]) = 2026
            && ISBLANK(Fact_Opportunities[Probability])
    )
)
\end{lstlisting}

\begin{lstlisting}[style=powerquery]
Forecast If Review Approved =
[Forecast Including Ready CRM]
    + [CRM Review Current-Year Value]
\end{lstlisting}

The resulting governed values are:

\begin{tabularx}{\textwidth}{
    >{\raggedright\arraybackslash}X
    >{\raggedleft\arraybackslash}p{4.2cm}
}

\toprule

Semantic Measure
&
Validated Result
\\

\midrule

CRM Review Current-Year Value
&
EUR 13,500.00
\\

CRM At-Risk Weighted Pipeline
&
EUR 1,989,178.99
\\

CRM Missing Probability Estimated Value
&
EUR 50,000.00
\\

Forecast If Review Approved
&
EUR 22,072,493.59
\\

\bottomrule

\end{tabularx}

The KPI cards have also been changed from the earlier dollar-symbol display to
EUR formatting. The current Card visual still abbreviates large numbers using
\texttt{K} and \texttt{M} display units. This is a presentation setting only;
it does not change the underlying DAX values shown above.

\begin{inferencebox}

\textbf{Inference -- Completed CRM Governance Measures}

The semantic layer now distinguishes four different confidence states instead
of collapsing them into one pipeline number. The EUR 807,900.00 Ready CRM value
remains the only CRM amount automatically included in the governed forecast.
EUR 13,500.00 is calculable but held for business review, EUR 1,989,178.99 is
probability-weighted pipeline that cannot yet be allocated to the current year,
and EUR 50,000.00 represents estimated pipeline whose Probability is missing.
This prevents unresolved CRM data from being presented as recognized or cleared
forecast revenue.

\end{inferencebox}

\subsection{17.30 Completed Disconnected Forecast-Bucket Dimension}

A dedicated table named \texttt{DimForecastBucket} is now present in the
semantic model. It contains the fields:

\begin{itemize}[leftmargin=7mm,itemsep=2pt]
    \item \texttt{BucketKey};
    \item \texttt{ForecastBucket}; and
    \item \texttt{SortOrder}.
\end{itemize}

The dimension is intentionally disconnected from the fact tables. The
forecast-bridge measures use the selected bucket label to present the governed
forecast states without adding another filter path between the analytical
facts.

The bridge currently exposes the following business treatments:

\begin{tabularx}{\textwidth}{
    >{\raggedright\arraybackslash}p{5.0cm}
    >{\raggedleft\arraybackslash}p{3.0cm}
    X
}

\toprule

Forecast Bucket
&
Amount
&
Treatment
\\

\midrule

Base Full-Year Forecast
&
EUR 21,251,093.59
&
Included -- Core Forecast
\\

Ready CRM -- Included
&
EUR 807,900.00
&
Included -- Ready CRM
\\

Review CRM -- Held
&
EUR 13,500.00
&
Held -- Business Review
\\

At-Risk Pipeline -- Not Included
&
EUR 1,989,178.99
&
Not Included -- Allocation Missing
\\

Missing Probability -- Not Included
&
EUR 50,000.00
&
Not Included -- Probability Missing
\\

\bottomrule

\end{tabularx}

\begin{inferencebox}

\textbf{Inference -- Forecast Bucket Dimension}

Keeping \texttt{DimForecastBucket} disconnected is deliberate. It acts as a
presentation and scenario-control dimension rather than a business filter on
Actual Sales, Accrual or CRM. This avoids ambiguous relationships while giving
the dashboard an auditable bridge from the source-aligned base forecast to
Ready, Review and unresolved CRM states.

\end{inferencebox}

\subsection{17.31 Completed CRM Calendar Dimension}

A separate CRM calendar named \texttt{DimCRMDate} has now been created in the
\texttt{04\_Dimensions} layer. It covers the full CRM horizon from 1 January
2024 through 31 December 2028.

The Power Query implementation is:

\begin{lstlisting}[style=powerquery]
let
    StartDate = #date(2024, 1, 1),
    EndDate = #date(2028, 12, 31),

    DateList =
        List.Dates(
            StartDate,
            Duration.Days(EndDate - StartDate) + 1,
            #duration(1, 0, 0, 0)
        ),

    ToTable =
        Table.FromList(
            DateList,
            Splitter.SplitByNothing(),
            {"Date"}
        ),

    ChangedType =
        Table.TransformColumnTypes(
            ToTable,
            {{"Date", type date}}
        ),

    AddedYear =
        Table.AddColumn(
            ChangedType,
            "Year",
            each Date.Year([Date]),
            Int64.Type
        ),

    AddedMonthNumber =
        Table.AddColumn(
            AddedYear,
            "MonthNumber",
            each Date.Month([Date]),
            Int64.Type
        ),

    AddedMonthName =
        Table.AddColumn(
            AddedMonthNumber,
            "MonthName",
            each Date.MonthName([Date]),
            type text
        ),

    AddedYearMonth =
        Table.AddColumn(
            AddedMonthName,
            "YearMonth",
            each Date.Year([Date]) * 100 + Date.Month([Date]),
            Int64.Type
        ),

    AddedMonthStart =
        Table.AddColumn(
            AddedYearMonth,
            "MonthStart",
            each Date.StartOfMonth([Date]),
            type date
        )

in
    AddedMonthStart
\end{lstlisting}

Full-dataset profiling confirms:

\begin{tabularx}{\textwidth}{
    >{\raggedright\arraybackslash\bfseries\color{AaltoNavy}}p{7cm}
    X
}

\toprule

CRM Date Metric
&
Result
\\

\midrule

Calendar period
&
1 January 2024 -- 31 December 2028
\\

Rows
&
1,827
\\

Date errors
&
0
\\

Date empty values
&
0
\\

Distinct / unique Date values
&
1,827 / 1,827
\\

Distinct Year values
&
5
\\

Distinct MonthNumber / MonthName values
&
12 / 12
\\

Distinct YearMonth values
&
60
\\

Distinct MonthStart values
&
60
\\

\bottomrule

\end{tabularx}

The row count of 1,827 is correct for the inclusive five-year period because
both 2024 and 2028 are leap years.

\begin{inferencebox}

\textbf{Inference -- CRM Calendar Readiness}

The CRM time horizon is now covered by a complete, unique and error-free daily
calendar that is independent of the 2026 operational calendar used by Actual
Sales and Project Accrual. This prevents the multi-year CRM pipeline from being
forced into the narrower 2026 Date dimension and creates a clean foundation for
explicit CRM date roles.

\end{inferencebox}

\subsection{17.32 Completed CRM Date-Role Relationships}

The CRM Date strategy has now been implemented using one active relationship
and one inactive role-playing relationship between \texttt{DimCRMDate} and
\texttt{Fact\_Opportunities}.

The relationship configuration is:

\begin{tabularx}{\textwidth}{
    >{\ttfamily\raggedright\arraybackslash}p{4.3cm}
    >{\ttfamily\raggedright\arraybackslash}p{4.8cm}
    >{\centering\arraybackslash}p{1.5cm}
    >{\centering\arraybackslash}p{1.5cm}
    X
}

\toprule

Dimension Side
&
Fact Side
&
Cardinality
&
Active
&
Filter
\\

\midrule

DimCRMDate[Date]
&
Fact\_Opportunities[ClosingDate]
&
1:*
&
Yes
&
Single
\\

DimCRMDate[Date]
&
Fact\_Opportunities[SimulatedProgramStartDate]
&
1:*
&
No
&
Single
\\

\bottomrule

\end{tabularx}

Power BI displays these relationships from the fact-table perspective as
\texttt{Many to one (*:1)}. The solid relationship line represents the active
\texttt{ClosingDate} role; the dotted relationship line represents the inactive
\texttt{SimulatedProgramStartDate} role. \texttt{EstimatedClosingDate} remains
unconnected at this stage.

The existing Business Unit relationships remain unchanged: one governed
\texttt{DimBusinessUnit} continues to filter all three facts with active,
single-direction relationships. The 2026 \texttt{DimDate} relationships to
Actual Sales and Project Accrual also remain unchanged.

\begin{inferencebox}

\textbf{Inference -- CRM Date-Role Model}

The relationship structure is correct and avoids competing active date paths.
Closing Date provides the default CRM time context, while Program Start Date is
available explicitly through DAX when required. This lets the same opportunity
be analysed by sales-closing timing or delivery/program-start timing without
creating ambiguous bidirectional filters or duplicating the CRM fact table.

\end{inferencebox}

\subsection{17.33 Immediate Next Step -- Validate CRM Date Roles with DAX}

The next control is to prove that the active and inactive CRM Date relationships
behave differently and both reconcile to the same source population.

First create a measure that uses the active \texttt{ClosingDate} relationship:

\begin{lstlisting}[style=powerquery]
CRM Opportunity Count by Closing Date =
COUNTROWS(Fact_Opportunities)
\end{lstlisting}

Then create a second measure that explicitly activates the inactive Program
Start relationship:

\begin{lstlisting}[style=powerquery]
CRM Opportunity Count by Program Start =
CALCULATE(
    COUNTROWS(Fact_Opportunities),
    USERELATIONSHIP(
        Fact_Opportunities[SimulatedProgramStartDate],
        DimCRMDate[Date]
    )
)
\end{lstlisting}

Create a Table or Matrix with:

\begin{itemize}[leftmargin=7mm,itemsep=2pt]
    \item \texttt{DimCRMDate[Year]} as the first validation row field;
    \item optionally \texttt{DimCRMDate[YearMonth]} for monthly detail;
    \item \texttt{CRM Opportunity Count by Closing Date}; and
    \item \texttt{CRM Opportunity Count by Program Start}.
\end{itemize}

At the grand-total level, both measures should reconcile to all 126 CRM
opportunities because both source date fields are populated and the CRM calendar
covers the complete 2024--2028 horizon. The year and month distributions should
not necessarily be identical, because an opportunity can close in one period
and start its simulated program in another.

\begin{warningbox}

Do not create a second active relationship between \texttt{DimCRMDate} and
\texttt{Fact\_Opportunities}. The Program Start relationship should remain
inactive and be activated only inside measures that intentionally analyse that
date role.

\end{warningbox}

\begin{inferencebox}

\textbf{Inference -- What This Validation Proves}

If both measures reconcile to 126 at total level while producing different
time distributions, the role-playing Date design is working as intended. The
model can then answer two different business questions---when opportunities are
expected to close and when their programs are expected to start---without
changing the underlying opportunity rows or creating ambiguous relationship
paths.

\end{inferencebox}


\subsection{17.34 Completed CRM Date-Role DAX Validation Results}

The planned CRM Date-role control has now been executed in the report view.
The validation matrix uses \texttt{DimCRMDate[Year]} together with the two
previously defined measures:

\begin{lstlisting}[style=powerquery]
CRM Opportunity Count by Closing Date
CRM Opportunity Count by Program Start
\end{lstlisting}

The completed result is:

\begin{tabularx}{\textwidth}{
    >{\centering\arraybackslash}p{2.5cm}
    >{\centering\arraybackslash}p{4.0cm}
    >{\centering\arraybackslash}X
}

\toprule

Year
&
CRM Opportunity Count by Closing Date
&
CRM Opportunity Count by Program Start
\\

\midrule

2024 & 1 & 1 \\
2025 & 1 & 1 \\
2026 & 106 & 56 \\
2027 & 15 & 64 \\
2028 & 3 & 4 \\

\midrule

\textbf{Total} & \textbf{126} & \textbf{126} \\

\bottomrule

\end{tabularx}

The two measures reconcile exactly to all 126 CRM opportunities at grand-total
level. Their yearly distributions are intentionally different. In particular,
106 opportunities close in 2026 while only 56 have simulated program start in
2026; in 2027 the pattern reverses to 15 closings and 64 program starts.

\begin{successbox}

\textbf{\color{SuccessGreen}{CRM Date-Role Validation Result}}

The active Closing Date relationship and the inactive Program Start relationship
both resolve the complete CRM population. The inactive relationship is therefore
working correctly through \texttt{USERELATIONSHIP}; it is not a broken or
orphaned path.

\end{successbox}

\begin{inferencebox}

\textbf{Inference -- CRM Date-Role Validation}

The result proves that Closing Date and Program Start represent genuinely
different business timelines rather than duplicate date fields. Using one active
and one inactive relationship preserves both interpretations without introducing
two competing active filter paths. The CRM temporal model is therefore
functionally validated, not merely structurally configured.

\end{inferencebox}

\subsection{17.35 Completed Project Dimension Construction and Profiling}

A Project dimension has now been created by combining the project identifiers
from Actual Sales and Project Accrual. CRM Opportunities are not included in
this dimension because the supplied CRM source contains an opportunity
description rather than a confirmed ProjectID mapping; no project assignment is
invented without source evidence.

The implemented Power Query logic is:

\begin{lstlisting}[style=powerquery]
let
    ActualProjects =
        Table.SelectColumns(
            Fact_ActualSales,
            {"ProjectID"}
        ),

    AccrualProjects =
        Table.SelectColumns(
            Fact_ProjectAccrual,
            {"ProjectID"}
        ),

    CombinedProjects =
        Table.Combine(
            {
                ActualProjects,
                AccrualProjects
            }
        ),

    CleanedProjects =
        Table.TransformColumns(
            CombinedProjects,
            {
                {
                    "ProjectID",
                    each
                        if _ = null then
                            null
                        else
                            Text.Trim(Text.Clean(Text.From(_))),
                    type text
                }
            }
        ),

    RemovedBlankProjects =
        Table.SelectRows(
            CleanedProjects,
            each
                [ProjectID] <> null
                and [ProjectID] <> ""
        ),

    DistinctProjects =
        Table.Distinct(
            RemovedBlankProjects,
            {"ProjectID"}
        ),

    SortedProjects =
        Table.Sort(
            DistinctProjects,
            {
                {"ProjectID", Order.Ascending}
            }
        ),

    AddedProjectKey =
        Table.AddIndexColumn(
            SortedProjects,
            "ProjectKey",
            1,
            1,
            Int64.Type
        ),

    ReorderedColumns =
        Table.ReorderColumns(
            AddedProjectKey,
            {
                "ProjectKey",
                "ProjectID"
            }
        )

in
    ReorderedColumns
\end{lstlisting}

The completed Power Query profile shows:

\begin{tabularx}{\textwidth}{
    >{\raggedright\arraybackslash\bfseries\color{AaltoNavy}}p{7cm}
    X
}

\toprule

Project-Dimension Metric
&
Result
\\

\midrule

Rows
&
252
\\

ProjectID valid values
&
100\%
\\

ProjectID errors
&
0
\\

ProjectID empty values
&
0
\\

Distinct ProjectID values
&
252
\\

Unique ProjectID values
&
252
\\

ProjectKey range
&
1--252
\\

\bottomrule

\end{tabularx}

The dimension was subsequently moved from \texttt{01\_Staging} into the
governed \texttt{04\_Dimensions} query group with model load retained. This
completes the query-governance placement that was still pending at the earlier
checkpoint.

\begin{inferencebox}

\textbf{Inference -- Project Dimension Readiness}

The combined Actual/Accrual project vocabulary contains 252 clean and unique
ProjectIDs, so the dimension provides a valid one-side business key. The later
semantic implementation confirmed that a separate fact-table surrogate-key merge
was not necessary for this case: both project facts already contain clean
\texttt{ProjectID} values that resolve to the unique dimension. The table has now
become a conformed Project dimension and its relationships have been validated.

\end{inferencebox}

\subsection{17.36 Immediate Next Step -- Govern Project Keys and Relationships}

\begin{warningbox}

\textbf{Implementation outcome note.} This subsection is retained as the original
planned control for audit history. During implementation, the proposed Left Outer
merge of \texttt{ProjectKey} into both facts was found to be unnecessary for the
current case because \texttt{DimProject[ProjectID]} is unique and both facts already
contain clean matching \texttt{ProjectID} values. The final implemented relationships
therefore use \texttt{ProjectID} directly. The completed state is documented in
Section 17.37.

\end{warningbox}

The next technical control is not to create another dimension. It is to finish
the semantic implementation of the Project dimension already created.

First, move \texttt{DimProject} from \texttt{01\_Staging} to
\texttt{04\_Dimensions}. Keep model load enabled.

Then add \texttt{ProjectKey} to both project-based facts using a Left Outer
merge on \texttt{ProjectID}:

\begin{enumerate}[leftmargin=9mm,itemsep=2pt]
    \item merge \texttt{Fact\_ActualSales[ProjectID]} with
    \texttt{DimProject[ProjectID]} and expand only \texttt{ProjectKey};
    \item confirm all 795 Actual Sales rows receive a non-null ProjectKey;
    \item merge \texttt{Fact\_ProjectAccrual[ProjectID]} with
    \texttt{DimProject[ProjectID]} and expand only \texttt{ProjectKey};
    \item confirm all 313 Accrual rows receive a non-null ProjectKey; and
    \item apply the changes and create the semantic relationships.
\end{enumerate}

The target relationships are:

\begin{tabularx}{\textwidth}{
    >{\ttfamily\raggedright\arraybackslash}p{5.0cm}
    >{\ttfamily\raggedright\arraybackslash}p{5.0cm}
    >{\centering\arraybackslash}p{1.7cm}
    >{\centering\arraybackslash}X
}

\toprule

Dimension Side
&
Fact Side
&
Cardinality
&
Filter
\\

\midrule

DimProject[ProjectKey]
&
Fact\_ActualSales[ProjectKey]
&
1:*
&
Single
\\

DimProject[ProjectKey]
&
Fact\_ProjectAccrual[ProjectKey]
&
1:*
&
Single
\\

\bottomrule

\end{tabularx}

Both relationships should be active. No Project relationship should be created
to \texttt{Fact\_Opportunities} unless a real CRM-to-Project identifier becomes
available from the source or business owner.

\begin{inferencebox}

\textbf{Inference -- Why ProjectKey Mapping Is the Next Control}

Creating a unique Project dimension is only the structural half of the task.
Propagating the surrogate key into the two project-based facts converts the
clean project vocabulary into a governed star-schema filter path. Validating
795 Actual rows and 313 Accrual rows after the merge will prove that no project
is lost during key standardization. Excluding CRM from this relationship avoids
fabricating a project assignment that the supplied source data does not prove.

\end{inferencebox}


\subsection{17.37 Completed Project Relationships and Functional Validation}

The Project semantic layer has now been completed. Two relationships were created
using the validated natural project identifier rather than introducing an
unnecessary additional merge step:

\begin{tabularx}{\textwidth}{
    >{\ttfamily\raggedright\arraybackslash}p{5.0cm}
    >{\ttfamily\raggedright\arraybackslash}p{5.0cm}
    >{\centering\arraybackslash}p{1.7cm}
    >{\centering\arraybackslash}X
}

\toprule

Dimension Side
&
Fact Side
&
Cardinality
&
Filter
\\

\midrule

DimProject[ProjectID]
&
Fact\_ActualSales[ProjectID]
&
1:*
&
Single
\\

DimProject[ProjectID]
&
Fact\_ProjectAccrual[ProjectID]
&
1:*
&
Single
\\

\bottomrule

\end{tabularx}

Both relationships are active. In the Power BI relationship Properties pane the
same structure is displayed from the fact perspective as
\texttt{Many to one (*:1)}. No Project relationship has been created to
\texttt{Fact\_Opportunities}, because the supplied CRM data does not provide a
confirmed ProjectID that would justify such a mapping.

The functional validation table uses \texttt{DimProject[ProjectID]} together with
existing fact measures. The completed grand totals are:

\begin{tabularx}{\textwidth}{
    >{\raggedright\arraybackslash\bfseries\color{AaltoNavy}}p{7.2cm}
    X
}

\toprule

Project-Relationship Validation Metric
&
Result
\\

\midrule

Actual Row Count
&
795
\\

Accrual Row Count
&
313
\\

Actual Net Sales
&
EUR 14,635,067.81
\\

Project Accrual Forecast
&
EUR 6,616,025.78
\\

Base Full-Year Forecast
&
EUR 21,251,093.59
\\

\bottomrule

\end{tabularx}

The Project-level rows also demonstrate combined fact filtering. For example,
project \texttt{O7901009} shows EUR 439,132.89 of Actual Net Sales and
EUR 308,913.00 of Project Accrual Forecast, producing EUR 748,045.89 of Base
Full-Year Forecast for the same ProjectID.

\begin{successbox}

\textbf{\color{SuccessGreen}{Project Relationship Validation Result}}

The Project dimension preserves all project-based source records and financial
values. The row counts remain 795 Actual and 313 Accrual, and the financial
totals remain identical to the previously validated base forecast. This proves
that the new Project dimension filters both facts correctly without losing or
duplicating value.

\end{successbox}

\begin{inferencebox}

\textbf{Inference -- Project Semantic Layer}

The shared Project dimension is now functionally complete. A user can select one
ProjectID and analyze recognized Actual revenue and remaining Project Accrual
forecast through the same governed project filter. Because the totals reconcile
exactly before and after introducing the dimension, the relationship layer adds
analytical capability without changing the financial baseline. The model now has
validated shared dimensions for Business Unit, calendar time and Project, while
CRM retains its separate governed date-role and opportunity logic.

\end{inferencebox}

\subsection{17.38 Immediate Next Step -- Project Filter Test and Dashboard Refinement}

Before final dashboard styling, perform one final interactive Project filter test:

\begin{enumerate}[leftmargin=9mm,itemsep=2pt]
    \item add a slicer using \texttt{DimProject[ProjectID]};
    \item select \texttt{O7901009};
    \item confirm \texttt{Actual Net Sales = EUR 439,132.89};
    \item confirm \texttt{Project Accrual Forecast = EUR 308,913.00};
    \item confirm \texttt{Base Full-Year Forecast = EUR 748,045.89}; and
    \item clear the slicer and confirm the grand total returns to
    \texttt{EUR 21,251,093.59}.
\end{enumerate}

After this final interaction test, the immediate presentation step is to create a
Project forecast-ranking visual for the executive dashboard: rank the top projects
by \texttt{Base Full-Year Forecast} and expose Actual versus remaining Accrual in
tooltips or a companion table.

\begin{inferencebox}

\textbf{Inference -- Transition to Dashboard Phase}

The remaining work is now mainly presentation and communication rather than core
data modelling. A Project slicer test closes the last shared-dimension functional
control; the model can then be presented through executive KPIs, Business Unit and
Project rankings, monthly forecast views and clearly separated CRM risk buckets.

\end{inferencebox}


\subsection{17.39 Completed Interactive Project Filter Validation}

The final interactive Project-dimension control has now been executed in Report
View using a slicer based on \texttt{DimProject[ProjectID]}. Project
\texttt{O7901009} was selected as the test case.

The report returned:

\begin{tabularx}{\textwidth}{
    >{\raggedright\arraybackslash\bfseries\color{AaltoNavy}}p{7.2cm}
    X
}

\toprule

Project Filter Validation Metric
&
Result
\\

\midrule

Selected ProjectID
&
\texttt{O7901009}
\\

Actual Row Count
&
6
\\

Accrual Row Count
&
4
\\

Actual Net Sales
&
EUR 439,132.89
\\

Project Accrual Forecast
&
EUR 308,913.00
\\

Base Full-Year Forecast
&
EUR 748,045.89
\\

\bottomrule

\end{tabularx}

The arithmetic reconciles exactly:

\begin{lstlisting}[style=powerquery]
EUR 439,132.89 + EUR 308,913.00 = EUR 748,045.89
\end{lstlisting}

The same interaction also confirms an intentional modelling boundary: Project
selection affects Actual Sales, Project Accrual and the Base Full-Year Forecast,
but it does not create a fabricated filter path into CRM Opportunities. CRM
forecast-governance values therefore remain independent unless a real source
mapping between opportunities and ProjectID becomes available.

\begin{successbox}

\textbf{\color{SuccessGreen}{Project Filter Validation Result}}

The Project slicer now behaves as intended across both project-based facts. The
selected ProjectID reconciles row counts and financial values exactly, proving
that the two active Project relationships are operational in Report View and not
merely structurally present in Model View.

\end{successbox}

\begin{inferencebox}

\textbf{Inference -- Interactive Project Filtering}

The Project dimension has passed both aggregate and interactive validation. A
single business ProjectID can now be used to analyze realized revenue and
remaining accrual forecast together, while CRM remains governed by its own
opportunity logic. This preserves semantic integrity and avoids implying a
Project-to-CRM relationship that is not present in the supplied source data.

\end{inferencebox}

\subsection{17.40 Completed Top-10 Project Forecast Ranking}

The first Project-level executive visual has now been created as a horizontal
bar chart titled:

\begin{lstlisting}[style=powerquery]
Base Full-Year Forecast by ProjectID
\end{lstlisting}

The visual uses:

\begin{itemize}[leftmargin=7mm,itemsep=2pt]
    \item category: \texttt{DimProject[ProjectID]};
    \item value: \texttt{Base Full-Year Forecast};
    \item visual-level Project filter: Top 10 by
    \texttt{Base Full-Year Forecast}; and
    \item descending value order.
\end{itemize}

The current ranking displayed in the report is:

\begin{enumerate}[leftmargin=9mm,itemsep=2pt]
    \item \texttt{D8001031};
    \item \texttt{L9109051};
    \item \texttt{D8001028};
    \item \texttt{D8001021};
    \item \texttt{O7901009};
    \item \texttt{D8001036};
    \item \texttt{D8001062};
    \item \texttt{O7401158};
    \item \texttt{D8001035}; and
    \item \texttt{P8404024}.
\end{enumerate}

The ranking uses the same validated Base Full-Year Forecast measure that
reconciles globally to EUR 21,251,093.59. No new financial calculation was
introduced solely for the chart.

\begin{successbox}

\textbf{\color{SuccessGreen}{Top-10 Project Ranking Result}}

The report can now identify the largest project-level contributors to the
source-aligned annual forecast. The visual is driven by the governed Project
dimension and existing financial measures, so project prioritization remains
consistent with the validated semantic model.

\end{successbox}

\begin{inferencebox}

\textbf{Inference -- Project Forecast Concentration}

Project-level ranking adds management value without changing the forecast
baseline. It turns the Project dimension from a navigation control into a
prioritization tool: users can immediately identify the projects with the
largest full-year financial contribution and then drill into their Actual and
Accrual components.

\end{inferencebox}

\subsection{17.41 Immediate Next Step -- Top-10 Actual versus Accrual Composition}

The next dashboard control is to explain \emph{why} each Top-10 project is large
by separating recognized Actual Net Sales from the remaining Project Accrual
Forecast.

Create a second horizontal \textbf{Stacked bar chart}. Use:

\begin{itemize}[leftmargin=7mm,itemsep=2pt]
    \item Y-axis: \texttt{DimProject[ProjectID]};
    \item X-axis / Values: \texttt{Actual Net Sales};
    \item X-axis / Values: \texttt{Project Accrual Forecast};
    \item visual filter on \texttt{ProjectID}: Top 10;
    \item By value: \texttt{Base Full-Year Forecast}; and
    \item keep the same Project slicer interaction as the existing Project visuals.
\end{itemize}

A suggested title is:

\begin{lstlisting}[style=powerquery]
Top 10 Projects - Actual vs Remaining Accrual
\end{lstlisting}

If the stacked chart does not retain the preferred descending order after the
second measure is added, keep the existing ranking chart as the authoritative
ordering visual and use the stacked chart as the companion composition view.
Do not change the underlying measures merely to force presentation order.

\begin{inferencebox}

\textbf{Inference -- Why Composition Is the Next Visual}

The ranking answers which projects are financially largest; the composition
view answers whether each project is already largely realized as Actual revenue
or still depends on future Accrual. Together the two visuals communicate both
scale and forecast maturity without introducing a new business assumption.

\end{inferencebox}


\subsection{17.42 Completed Top-10 Actual versus Accrual Composition}

The companion Project composition visual planned in Section 17.41 has now been
implemented as a horizontal stacked bar chart titled:

\begin{lstlisting}[style=powerquery]
Top 10 Projects - Actual vs Remaining Accrual
\end{lstlisting}

The visual uses:

\begin{itemize}[leftmargin=7mm,itemsep=2pt]
    \item Y-axis: \texttt{DimProject[ProjectID]};
    \item values: \texttt{Actual Net Sales} and
    \texttt{Project Accrual Forecast};
    \item Top-10 filtering by \texttt{Base Full-Year Forecast}; and
    \item the same governed Project relationship path already validated for the
    ranking chart and Project slicer.
\end{itemize}

No new financial measure was introduced. The visual decomposes the same
source-aligned Base Full-Year Forecast into its realized Actual component and its
remaining Accrual component.

\begin{successbox}

\textbf{\color{SuccessGreen}{Project Composition Visual Result}}

The report now contains both a Top-10 Project ranking and a companion composition
view. The first visual identifies the largest project-level forecast contributors;
the second explains how much of each project's annual value has already been
recognized as Actual Net Sales and how much still depends on remaining Accrual.

\end{successbox}

\begin{inferencebox}

\textbf{Inference -- Project Forecast Maturity}

Project size and project maturity can now be interpreted together. A large bar in
the ranking view does not by itself reveal whether the project is already realized
or still forecast-dependent; the composition chart adds that distinction without
changing the validated financial baseline.

\end{inferencebox}


\subsection{17.43 Executive Dashboard KPI and Business Unit Control Layer}

A dedicated report page named \texttt{Executive Dashboard} has now been created.
The first dashboard layer contains eight KPI cards:

\begin{tabularx}{\textwidth}{
    >{\raggedright\arraybackslash\bfseries\color{AaltoNavy}}X
    >{\raggedleft\arraybackslash}p{4.2cm}
}

\toprule

Executive KPI
&
Displayed Value
\\

\midrule

Actual Net Sales
&
EUR 14.64M
\\

Project Accrual Forecast
&
EUR 6.62M
\\

Base Full-Year Forecast
&
EUR 21.25M
\\

Forecast Including Ready CRM
&
EUR 22.06M
\\

CRM Ready Current-Year Value
&
EUR 807.90K
\\

CRM Review Current-Year Value
&
EUR 13.50K
\\

CRM At-Risk Weighted Pipeline
&
EUR 1.99M
\\

CRM Missing Probability Estimated Value
&
EUR 50.00K
\\

\bottomrule

\end{tabularx}

All eight cards now use EUR formatting consistently. A report slicer based on:

\begin{lstlisting}[style=powerquery]
DimBusinessUnit[BusinessUnitName]
\end{lstlisting}

has also been added and exposes the four governed members:

\begin{itemize}[leftmargin=7mm,itemsep=2pt]
    \item ExEd Programs BU;
    \item Qualification Programs BU;
    \item University Programs BU; and
    \item Unassigned / Missing BU.
\end{itemize}

The existing semantic model already places \texttt{DimBusinessUnit} on the one
side of active Single-direction relationships to all three analytical facts, so
this slicer uses the validated conformed-dimension path rather than source-specific
Business Unit text fields.

\begin{inferencebox}

\textbf{Inference -- Executive Control Layer}

The Executive Dashboard now starts from governed totals rather than raw columns.
The KPI row separates the committed base forecast, cleared Ready CRM contribution
and unresolved CRM risk categories, while the shared Business Unit slicer provides
one consistent management filter across Actual Sales, Project Accrual and CRM.

\end{inferencebox}


\subsection{17.44 Executive Monthly Forecast Visual -- Presentation Correction}

A 12-month visual titled:

\begin{lstlisting}[style=powerquery]
Monthly Base Forecast - Actual vs Accrual
\end{lstlisting}

has been added with chronological January--December ordering. The underlying
fields are correct:

\begin{itemize}[leftmargin=7mm,itemsep=2pt]
    \item X-axis: \texttt{DimDate[MonthName]};
    \item value: \texttt{Actual Net Sales};
    \item value: \texttt{Project Accrual Forecast}; and
    \item month order governed by \texttt{DimDate[MonthNumber]}.
\end{itemize}

The visual correctly separates the source periods: Actual Net Sales is present in
January--August and Project Accrual Forecast in September--December.

However, the current screenshot shows a percentage y-axis from 0\% to 100\%.
This identifies the selected visual as a \textbf{100\% stacked column chart}.
That chart type normalizes every populated month to 100\% and therefore hides the
financial difference between, for example, a EUR 0.4M month and a EUR 2.7M month.

The chart must therefore be changed to the standard
\textbf{Stacked column chart}. No field or measure should be changed.

\begin{warningbox}

\textbf{Do not keep the 100\% stacked version for the final dashboard.}

The business question is monthly forecast magnitude, not percentage composition.
The corrected chart should display a monetary y-axis and preserve the same
January--December Actual/Accrual split.

\end{warningbox}

\begin{inferencebox}

\textbf{Inference -- Monthly Forecast Presentation}

The Date model and financial measures are functioning correctly; the remaining
issue is visual encoding rather than data logic. Replacing the 100\% stacked chart
with a standard stacked column chart will make monthly magnitude immediately
comparable while retaining the validated Actual-versus-Accrual structure.

\end{inferencebox}


\subsection{17.45 Immediate Next Step -- Correct Monthly Chart and Add BU Forecast View}

The immediate next action is:

\begin{enumerate}[leftmargin=9mm,itemsep=2pt]
    \item select the current monthly chart on the \texttt{Executive Dashboard};
    \item in Visualizations, choose \textbf{Stacked column chart}, not
    \textbf{100\% stacked column chart};
    \item keep \texttt{DimDate[MonthName]} on the X-axis;
    \item keep \texttt{Actual Net Sales} and
    \texttt{Project Accrual Forecast} as values;
    \item verify that the y-axis changes from percentages to monetary values;
    \item select each Business Unit in the slicer and confirm the KPI cards and
    monthly chart respond consistently; and
    \item after that validation, add a Business Unit financial-composition visual
    using \texttt{DimBusinessUnit[BusinessUnitName]} with
    \texttt{Actual Net Sales}, \texttt{Project Accrual Forecast} and
    \texttt{Base Full-Year Forecast}.
\end{enumerate}

\begin{inferencebox}

\textbf{Inference -- Next Dashboard Control}

This sequence closes the only visible issue in the current Executive Dashboard
and then proves the shared Business Unit dimension at presentation level. Once
that interaction is validated, the remaining work is primarily final layout,
titles, conditional emphasis and interview storytelling rather than core model
construction.

\end{inferencebox}



\subsection{17.46 Completed Executive Monthly Forecast Correction}

The presentation correction planned in Sections 17.44--17.45 has now been
completed. The visual titled:

\begin{lstlisting}[style=powerquery]
Monthly Base Forecast - Actual vs Accrual
\end{lstlisting}

now uses a standard stacked column chart rather than the earlier 100\% stacked
version. The field configuration remains unchanged:

\begin{itemize}[leftmargin=7mm,itemsep=2pt]
    \item X-axis: \texttt{DimDate[MonthName]};
    \item Y-axis value: \texttt{Actual Net Sales};
    \item Y-axis value: \texttt{Project Accrual Forecast}; and
    \item tooltip: \texttt{Base Full-Year Forecast}.
\end{itemize}

The corrected chart now displays monetary magnitude on the y-axis instead of
percentage share. January--August continues to display Actual Net Sales and
September--December continues to display Project Accrual Forecast. The month
order remains chronological.

The explicit \texttt{Unassigned / Missing BU} test also produces only
September--December Accrual bars, which is consistent with the known
EUR 42,550.00 unassigned accrual exposure and the absence of Actual Sales for
that member.

\begin{successbox}

\textbf{\color{SuccessGreen}{Monthly Forecast Visual Correction Result}}

The monthly executive visual now communicates both timing and financial magnitude
correctly. No DAX measure, Date relationship or source transformation needed to be
changed; only the chart encoding was corrected.

\end{successbox}

\begin{inferencebox}

\textbf{Inference -- Correct Monthly Magnitude}

The successful correction confirms that the earlier issue was purely a
presentation-layer problem. The same validated Date model and measures now produce
an executive visual in which month-to-month EUR differences are directly
comparable without percentage normalization.

\end{inferencebox}


\subsection{17.47 Completed Executive Business Unit Slicer Validation}

The \texttt{Executive Dashboard} Business Unit slicer was tested interactively
against all four governed members. The test confirms that the same selection
propagates to Actual Sales, Project Accrual and CRM governance measures.

The core financial result by selected Business Unit is:

\begin{tabularx}{\textwidth}{
    >{\raggedright\arraybackslash}p{3.5cm}
    >{\raggedleft\arraybackslash}p{2.7cm}
    >{\raggedleft\arraybackslash}p{2.7cm}
    >{\raggedleft\arraybackslash}p{2.7cm}
    >{\raggedleft\arraybackslash}X
}

\toprule

Business Unit
&
Actual
&
Accrual
&
Base Forecast
&
Forecast + Ready CRM
\\

\midrule

ExEd Programs BU
&
EUR 6,044,545.34
&
EUR 3,009,546.88
&
EUR 9,054,092.22
&
EUR 9,581,969.72
\\

Qualification Programs BU
&
EUR 4,866,104.98
&
EUR 1,772,854.87
&
EUR 6,638,959.85
&
EUR 6,674,209.85
\\

University Programs BU
&
EUR 3,724,417.49
&
EUR 1,791,074.03
&
EUR 5,515,491.52
&
EUR 5,760,264.02
\\

Unassigned / Missing BU
&
EUR 0.00
&
EUR 42,550.00
&
EUR 42,550.00
&
EUR 42,550.00
\\

\bottomrule

\end{tabularx}

The CRM governance measures under the same slicer context are:

\begin{tabularx}{\textwidth}{
    >{\raggedright\arraybackslash}p{3.5cm}
    >{\raggedleft\arraybackslash}p{2.7cm}
    >{\raggedleft\arraybackslash}p{2.7cm}
    >{\raggedleft\arraybackslash}p{2.7cm}
    >{\raggedleft\arraybackslash}X
}

\toprule

Business Unit
&
Ready CRM
&
Review
&
At-Risk
&
Missing Probability
\\

\midrule

ExEd Programs BU
&
EUR 527,877.50
&
EUR 0.00
&
EUR 1,229,613.75
&
EUR 0.00
\\

Qualification Programs BU
&
EUR 35,250.00
&
EUR 0.00
&
EUR 615,625.00
&
EUR 0.00
\\

University Programs BU
&
EUR 244,772.50
&
EUR 13,500.00
&
EUR 143,940.24
&
EUR 50,000.00
\\

Unassigned / Missing BU
&
EUR 0.00
&
EUR 0.00
&
EUR 0.00
&
EUR 0.00
\\

\bottomrule

\end{tabularx}

The card visuals currently render some zero-row filtered results as
\texttt{--}. This is an expected DAX \texttt{BLANK()} display when the selected
Business Unit has no record in that CRM treatment bucket. It is not evidence of
a broken relationship. If a later presentation preference requires explicit
EUR 0.00 cards, display-specific \texttt{COALESCE} measures can be added without
changing the validated business logic.

\begin{successbox}

\textbf{\color{SuccessGreen}{Business Unit Slicer Validation Result}}

All four selections produce the expected financial and CRM distributions.
ExEd, Qualification and University each return their previously reconciled
Actual, Accrual, Ready CRM and At-Risk values. The Unassigned member returns
only EUR 42,550.00 of Project Accrual Forecast and no Actual or CRM contribution,
exactly matching the known four missing-BU accrual records.

\end{successbox}

\begin{inferencebox}

\textbf{Inference -- Report-Level Relationship Validation}

The Business Unit model is now validated not only through row-count matrices and
standalone financial tables, but also through live Executive Dashboard
interaction. One dimension selection filters all three analytical facts in a
business-consistent way, while the explicit Unassigned member prevents missing
source classifications from leaking into a real Business Unit.

\end{inferencebox}


\subsection{17.48 Immediate Next Step -- Business Unit Forecast Composition}

The next executive visual should compare the composition of the base forecast
across Business Units without creating any new measure.

Create a standard horizontal \textbf{Stacked bar chart} with:

\begin{itemize}[leftmargin=7mm,itemsep=2pt]
    \item Y-axis: \texttt{DimBusinessUnit[BusinessUnitName]};
    \item X-axis / Values: \texttt{Actual Net Sales};
    \item X-axis / Values: \texttt{Project Accrual Forecast};
    \item Tooltip: \texttt{Base Full-Year Forecast};
    \item Tooltip: \texttt{Forecast Including Ready CRM}; and
    \item Tooltip: \texttt{CRM Ready Current-Year Value}.
\end{itemize}

Sort the visual by:

\begin{lstlisting}[style=powerquery]
Base Full-Year Forecast
\end{lstlisting}

in descending order. A suggested title is:

\begin{lstlisting}[style=powerquery]
Base Forecast by Business Unit - Actual vs Remaining Accrual
\end{lstlisting}

With the Business Unit slicer cleared, the expected Base Full-Year Forecast
ordering is:

\begin{enumerate}[leftmargin=9mm,itemsep=2pt]
    \item ExEd Programs BU -- EUR 9,054,092.22;
    \item Qualification Programs BU -- EUR 6,638,959.85;
    \item University Programs BU -- EUR 5,515,491.52; and
    \item Unassigned / Missing BU -- EUR 42,550.00.
\end{enumerate}

Keep the existing measures unchanged. The purpose of this visual is presentation,
not a new calculation.

\begin{inferencebox}

\textbf{Inference -- Why Business Unit Composition Is Next}

The dashboard can already answer the annual total, monthly timing and Project
concentration questions. The Business Unit composition chart adds the management
question of where the source-aligned forecast sits organizationally and whether
each Business Unit is already realized as Actual revenue or still depends on
remaining Accrual. It reuses the same validated semantic model and therefore adds
explanation without adding forecast assumptions.

\end{inferencebox}




\subsection{17.49 Completed Business Unit Forecast Composition}

The Business Unit composition visual planned in Section 17.48 has now been
implemented on the \texttt{Executive Dashboard}. The visual is a standard
horizontal stacked bar chart with:

\begin{itemize}[leftmargin=7mm,itemsep=2pt]
    \item Y-axis: \texttt{DimBusinessUnit[BusinessUnitName]};
    \item X-axis / Values: \texttt{Actual Net Sales};
    \item X-axis / Values: \texttt{Project Accrual Forecast}; and
    \item the existing base-forecast measures available for context and tooltip use.
\end{itemize}

With the Business Unit slicer cleared, the chart displays the expected descending
organizational pattern: ExEd Programs BU first, followed by Qualification Programs
BU, University Programs BU and the explicit \texttt{Unassigned / Missing BU}
member. The stacked composition also preserves the distinction between recognized
Actual revenue and remaining Accrual forecast.

\begin{successbox}

\textbf{\color{SuccessGreen}{Business Unit Composition Result}}

The Executive Dashboard can now answer not only the annual total and monthly timing
questions, but also where the base forecast sits organizationally and how much of
each Business Unit's annual value has already been realized versus remaining in
Accrual.

\end{successbox}

\begin{inferencebox}

\textbf{Inference -- Business Unit Forecast Composition}

The shared Business Unit dimension is now demonstrated at presentation level as
well as model level. The same governed dimension that reconciles row counts and KPI
cards also produces a coherent financial composition across all four Business Unit
members, including the explicit unassigned accrual exposure.

\end{inferencebox}


\subsection{17.50 Completed Forecast Governance Table}

A compact executive control table titled:

\begin{lstlisting}[style=powerquery]
Forecast Governance - Included vs At Risk
\end{lstlisting}

has now been added to the Executive Dashboard. It uses:

\begin{itemize}[leftmargin=7mm,itemsep=2pt]
    \item \texttt{DimForecastBucket[ForecastBucket]};
    \item measure \texttt{Forecast Bucket Amount}; and
    \item measure \texttt{Forecast Bucket Treatment}.
\end{itemize}

The rendered order is the intended governance order and the displayed values are:

\begin{tabularx}{\textwidth}{
    >{\raggedright\arraybackslash}p{5.3cm}
    >{\raggedleft\arraybackslash}p{3.0cm}
    X
}

\toprule

Forecast Bucket
&
Amount
&
Treatment
\\

\midrule

Base Full-Year Forecast
&
EUR 21,251,093.59
&
Included -- Core Forecast
\\

Ready CRM -- Included
&
EUR 807,900.00
&
Included -- Ready CRM
\\

Review CRM -- Held
&
EUR 13,500.00
&
Held -- Business Review
\\

At-Risk Pipeline -- Not Included
&
EUR 1,989,178.99
&
Not Included -- Allocation Missing
\\

Missing Probability -- Not Included
&
EUR 50,000.00
&
Not Included -- Probability Missing
\\

\bottomrule

\end{tabularx}

The table therefore keeps the cleared forecast, held amount and unresolved pipeline
visibly separate instead of presenting all CRM value as one additive forecast.

\begin{successbox}

\textbf{\color{SuccessGreen}{Forecast Governance Result}}

The executive page now contains a direct control explaining which values are included,
which are held for business review and which remain outside the governed forecast.
The financial meaning of each CRM treatment is visible beside its amount.

\end{successbox}

\begin{inferencebox}

\textbf{Inference -- Included versus At-Risk Forecast}

The dashboard now communicates forecast confidence as well as forecast magnitude.
EUR 21.251M forms the source-aligned core forecast and EUR 807.9K of Ready CRM is
cleared for inclusion, while EUR 13.5K remains held and the EUR 1.989M allocation-
missing pipeline plus EUR 50K missing-Probability exposure remain explicitly outside
the committed forecast.

\end{inferencebox}


\subsection{17.51 Executive Dashboard Analytical Composition Checkpoint}

The Executive Dashboard now contains the following complementary management views:

\begin{itemize}[leftmargin=7mm,itemsep=2pt]
    \item eight annual financial and CRM-governance KPI cards;
    \item one governed Business Unit slicer;
    \item the \texttt{Forecast Governance -- Included vs At Risk} control table;
    \item Business Unit Actual-versus-Accrual forecast composition;
    \item monthly Actual-versus-Accrual timing across January--December; and
    \item Top-10 Project Actual-versus-remaining-Accrual composition.
\end{itemize}

This page now covers five distinct executive questions: total forecast magnitude,
CRM readiness and risk, organizational concentration, monthly timing and Project
concentration. No new source transformations or forecast assumptions were required
to create the two latest visuals.

\begin{inferencebox}

\textbf{Inference -- Executive Dashboard Coverage}

The report has moved from model construction into final validation and presentation.
The remaining work is no longer core semantic-model development; it is primarily
interaction verification, field cleanup, layout refinement and preparation of a
concise interview narrative explaining the governed forecast.

\end{inferencebox}


\subsection{17.52 Immediate Next Step -- Final Interaction Validation}

Because the Business Unit composition and Forecast Governance visuals were added
after the earlier slicer test, the next control is to re-run the Business Unit
interaction test across the complete Executive Dashboard.

Start with \texttt{ExEd Programs BU}. The expected KPI context is:

\begin{itemize}[leftmargin=7mm,itemsep=2pt]
    \item Actual Net Sales: EUR 6,044,545.34;
    \item Project Accrual Forecast: EUR 3,009,546.88;
    \item Base Full-Year Forecast: EUR 9,054,092.22;
    \item Forecast Including Ready CRM: EUR 9,581,969.72;
    \item CRM Ready Current-Year Value: EUR 527,877.50;
    \item CRM Review Current-Year Value: EUR 0.00 / blank display;
    \item CRM At-Risk Weighted Pipeline: EUR 1,229,613.75; and
    \item CRM Missing Probability Estimated Value: EUR 0.00 / blank display.
\end{itemize}

During this selection, verify that:

\begin{enumerate}[leftmargin=9mm,itemsep=2pt]
    \item all eight KPI cards respond;
    \item the Forecast Governance table recalculates its bucket amounts;
    \item the Business Unit composition chart reflects the selected Business Unit;
    \item the monthly chart recalculates the selected Business Unit's monthly profile;
    \item the Top-10 Project chart re-ranks projects inside the selected Business Unit;
    \item clearing the slicer restores the full-company totals; and
    \item no visual remains unintentionally disconnected because of an Edit
    Interactions setting.
\end{enumerate}

Repeat the same interaction check for Qualification Programs BU, University Programs
BU and \texttt{Unassigned / Missing BU}. After this control passes, the next work
package is final semantic-field hiding and visual polish rather than another major
modelling step.

\begin{inferencebox}

\textbf{Inference -- Why Final Interaction Validation Is Next}

The new visuals reuse already validated measures, but a dashboard is only fully
validated when filter propagation is confirmed at report level. Re-testing all four
Business Unit selections proves that the complete executive page behaves as one
coherent analytical product rather than as a collection of individually correct
visuals.

\end{inferencebox}


% ============================================================
% INTERVIEW MESSAGE
% ============================================================

\section{18. Key Interview Message}

\begin{summarybox}

The implementation does not treat the provided Excel files as simple
reporting inputs.

It establishes a controlled forecasting process with raw-data staging,
repeatable transformations, documented quality rules, grain validation,
financial reconciliation and traceable analytical outputs.

The source Total is preserved as a control value instead of being silently
deleted. Exceptional zero, negative and missing-classification records are
retained transparently for Finance review rather than being hidden through
automatic filtering.

The exact-duplicate control quantifies 8 groups, 20 involved source
records and 12 excess records. It also estimates EUR 175,318.63 of
potential overstatement without altering the source-aligned forecast
before business validation.

The broader repeated-key control identifies 10 project-month groups
containing 32 source records. It distinguishes four same-amount groups
from six groups containing multiple amounts, demonstrating why repeated
records must be classified before any removal decision is made.

The completed \texttt{RepeatedKeyStatus} classification further separates
the results into 4 same-amount groups, 4 mixed groups and 2 groups containing
only multiple distinct components. This preserves the financial meaning of
the source while directing Finance review to the correct exceptions.

The CRM opportunity source has also been profiled across all 126 records and
14 columns. The profile confirms complete opportunity descriptions, three
Business Unit values and reliable numeric value fields, while exposing missing
probability and annual-allocation values that require explicit quality rules.

Explicit analytical data types have now been assigned to all 14 CRM columns.
The conversion retained all 126 records and produced zero errors. This provides
a technically stable basis for date analysis, probability calculations and
financial validation while preserving missing business values for transparent
quality classification.

The four CRM text columns were standardized using Trim and Clean, and all 14
source fields were assigned consistent analytical names. The subsequent
\texttt{OpportunityQualityStatus} control classified every opportunity with
zero errors and zero empty statuses. It identified 95 valid records, 26 records
without annual allocation, 2 zero-value records, 2 records with missing
probability and 1 record whose combined allocation exceeds 100\%. No source row
was deleted automatically; the findings remain visible for controlled business
review.

An independent CRM weighted-value reconciliation now confirms 123 matches
within a EUR 0.01 tolerance. Two records remain not testable because Probability
is missing, while one material mismatch has been isolated with a source-minus-
calculated variance of EUR -468,750.00. The source value has not been changed;
the exception is preserved for targeted business validation.

The comparison between \texttt{OpportunityValue} and \texttt{EstimatedValue}
has also been converted from a visual observation into an explicit control. All
126 CRM records return \texttt{Match} within the EUR 0.01 tolerance, with zero
errors and zero empty classifications.

The current-year weighted amount is now independently recalculated from
estimated value, probability and the explicit current-year allocation. The
resulting \texttt{CurrentYearWeightedVariance} contains zero errors, and the
final \texttt{CurrentYearWeightedCheck} is 100\% valid with zero errors and zero
empty classifications. The dedicated status-summary query confirms 85
\texttt{Match} records, 39 records that are not testable because the current-year
allocation is missing, and 2 records that are not testable because Probability
is missing. No current-year \texttt{Mismatch} is present among the 85 fully
testable records.

A Power Query dependency-order issue was also resolved transparently: the
current-year variance is now created before the final check that references it.
No source value was changed to solve the technical issue.

The broader annual-allocation control now classifies all 126 CRM opportunities.
Seventy records have both current-year and next-year allocations and sum to
exactly 100\%. Twenty-nine records have both allocations missing, 12 are missing
only the current-year allocation and 11 are missing only the next-year
allocation. Three complete records sum to less than 100\%, while one record
exceeds 100\% and remains a targeted business exception. Missing allocations
were not replaced with zero.

The grouped allocation-status query independently confirms all six status counts
and reconciles to 126 records. A dedicated exception query further isolates the
four non-100\% cases. The over-allocation is \texttt{Open custom opportunity 111}
in University Programs BU, where 30\% current-year plus 80\% next-year equals
110\%. The three below-100\% cases total 50\%, 50\% and 30\% and have different
closing/program-start date patterns, so they remain business-review items rather
than being normalized automatically.

The missing-allocation horizon analysis further shows that 41 of the 52 records
with one or both annual allocations missing close in 2026. The dedicated
near-term detail query confirms that these 41 records comprise 21 with both
allocations missing, 9 missing only the current-year allocation and 11 missing
only the next-year allocation. This means the allocation-completeness issue is
primarily relevant to the near-term forecast rather than only to distant
pipeline opportunities.

Those 41 opportunities represent EUR 4,208,699.23 of estimated pipeline and
EUR 2,317,982.74 of probability-weighted pipeline. Thirty of the 41 are missing
the current-year allocation itself and therefore cannot yet be allocated
reliably into the 2026 forecast; these 30 represent EUR 1,989,178.99 of
weighted pipeline. ExEd Programs BU is the largest concentration, with
EUR 1,229,613.75 across 18 affected records.

The 2026 forecast-readiness control now evaluates all 106 opportunities closing
in 2026. Seventy-five are technically Ready and contribute EUR 807,900.00 of
reconciled current-year CRM value. Twenty-eight remain Not Ready because the
current-year allocation is missing, two remain Not Ready because Probability is
missing, and one EUR 13,500.00 calculated current-year amount is held for review
because its cross-year allocation totals 110\%. The model therefore separates
usable forecast value from unresolved pipeline instead of forcing missing inputs
to zero or treating every weighted opportunity amount as current-year revenue.

The completed CRM forecast bridge formalizes those treatments into separate
financial columns: EUR 807,900.00 Ready forecast, EUR 13,500.00 held for review,
EUR 1,989,178.99 of at-risk weighted pipeline and EUR 50,000.00 of estimated
pipeline with missing Probability.

A cross-source Business Unit audit also confirms that CRM and Accrual already
share the same three populated Business Unit labels, while Actual Sales uses
shorter names for the same categories. Four Accrual rows still have no Business
Unit and remain explicitly unassigned. The model can therefore standardize BU
labels through one governed dimension without guessing the missing classifications.

That dimension has now been implemented as \texttt{DimBusinessUnit} with keys
1--3 for the three canonical Business Units and key 0 for
\texttt{Unassigned / Missing BU}. \texttt{Fact\_ActualSales} has also been
mapped successfully: all 795 rows have a valid BusinessUnitKey, with 346 rows
on key 1, 120 on key 2 and 329 on key 3, while the original source label remains
available for lineage.

The same governed key has now been added to Project Accrual and CRM
Opportunities. All 313 accrual rows have valid keys, including the 4 missing-BU
records mapped explicitly to key 0, while all 126 CRM opportunities map to the
three populated Business Units. Staging and data-quality query loading has been
disabled for the report model without deleting their Power Query logic.

The visible semantic model is therefore reduced to one Business Unit dimension
and three analytical facts. Three active one-to-many relationships connect
\texttt{DimBusinessUnit} to Actual Sales, Project Accrual and CRM Opportunities
using Single cross-filter direction.

Those Business Unit relationships have now been validated functionally with
three row-count measures. The resulting matrix reconciles exactly to
795 Actual rows, 313 Accrual rows and 126 CRM Opportunity rows, including the
expected 0/4/0 pattern for the Unassigned member. This confirms the dimension
filters all three facts correctly rather than merely appearing structurally
valid in Model View.

A complete 2026 \texttt{DimDate} has also been created with 365 unique daily
dates, 12 months, 12 YearMonth values, 12 MonthStart values and 4 quarters.
The calendar is now connected actively to Actual Sales through
\texttt{MonthStart} and to Project Accrual through \texttt{ForecastMonth}, with
Single cross-filter direction.

The Date relationships have also been validated at month level. Actual Sales
appears only in January--August and reconciles to all 795 rows, while Project
Accrual appears only in September--December and reconciles to all 313 rows.
This confirms one shared Date axis can support the continuous 2026 base forecast
without appending the facts. CRM remains intentionally outside the Date path
until its multiple date roles are modelled explicitly.

Chronological month sorting is now complete and the first financial DAX
measures are operational. Actual Net Sales reconciles to EUR 14,635,067.81,
Project Accrual Forecast reconciles to EUR 6,616,025.78 and the source-aligned
Base Full-Year Forecast is therefore EUR 21,251,093.59. The monthly visual shows
one continuous 2026 series: realized Actual value in January--August and
ongoing-project Accrual value in September--December.

The accrual layer remains deliberately source-aligned. The previously identified
EUR 175,318.63 potential duplicate overstatement has not been removed because
Finance validation is still pending. This preserves reconciliation and makes the
difference between a controlled source forecast and a later business-approved
adjusted scenario explicit.

The three base financial measures have now also been validated by Business Unit.
ExEd contributes EUR 9,054,092.22 to the base forecast, Qualification Programs
EUR 6,638,959.85, University Programs EUR 5,515,491.52 and the explicit
Unassigned member EUR 42,550.00. The rows reconcile exactly to the annual
EUR 21,251,093.59 base forecast.

\texttt{DimDate} has also been formally marked as the model Date table using its
\texttt{Date} column, completing the semantic calendar control.

The Ready CRM layer is now implemented in DAX. Seventy-five opportunities
contribute EUR 807,900.00 of reconciled current-year value: EUR 527,877.50 from
ExEd Programs, EUR 35,250.00 from Qualification Programs and EUR 244,772.50
from University Programs. Adding only this cleared CRM amount to the
source-aligned base forecast produces a current governed annual forecast of
EUR 22,058,993.59.

The next semantic-model control is to surface the remaining CRM treatment
buckets separately: EUR 13,500.00 held for review, EUR 1,989,178.99 of
At-Risk weighted pipeline and EUR 50,000.00 of estimated pipeline with missing
Probability. Those amounts should remain separate from the cleared forecast
until their business inputs are resolved. CRM also remains deliberately outside
the shared monthly Date path until its multiple date roles are modelled
explicitly.

That previously planned control has now been completed. The three governance
measures are visible as EUR-formatted KPIs, and the optional review-approved
scenario reconciles to EUR 22,072,493.59. A disconnected
\texttt{DimForecastBucket} now presents the five forecast treatments without
adding a new relationship path.

The CRM time model has also been completed with a dedicated 2024--2028
\texttt{DimCRMDate}. Closing Date is the active relationship and Simulated
Program Start Date is the inactive role-playing relationship. DAX validation
reconciles both roles to all 126 opportunities while preserving materially
different yearly timing: 106/56 in 2026 and 15/64 in 2027 for Closing Date
versus Program Start.

The Project dimension is now constructed from Actual Sales and Project Accrual.
It contains 252 clean, distinct and unique ProjectIDs with zero errors and zero
empty values and has been moved into \texttt{04\_Dimensions}. Two active
Single-direction relationships now connect \texttt{DimProject[ProjectID]} to
Actual Sales and Project Accrual. Functional validation preserves all 795 Actual
rows, 313 Accrual rows and the complete EUR 21,251,093.59 base forecast. The
technical \texttt{ProjectKey} remains available in the dimension, but no fact-table
merge was required because the natural ProjectID already forms a clean validated
relationship key.

The Project layer has now also passed its final interactive slicer test. Selecting
\texttt{O7901009} produces EUR 439,132.89 of Actual Net Sales and EUR 308,913.00
of Project Accrual Forecast, reconciling exactly to EUR 748,045.89 of Base Full-Year
Forecast. The Top-10 Project ranking visual is now present in the report and converts
that validated Project model into an executive prioritization view. CRM remains
intentionally outside the Project filter path because no confirmed CRM-to-Project
identifier exists in the supplied source.

The core semantic model is now technically mature enough to support the final
presentation layer. The next implementation focus is executive dashboard refinement,
Project ranking and final interview storytelling.

The Project presentation layer now includes both Top-10 forecast ranking and
Actual-versus-Accrual composition. A dedicated Executive Dashboard also exposes
the eight governed financial and CRM-risk KPIs together with the canonical
Business Unit slicer. The 12-month Actual-versus-Accrual visual has been added and
its underlying Date logic is correct; the remaining presentation correction is to
replace the current 100\% stacked column rendering with a standard stacked column
chart so monthly bar height communicates EUR magnitude.

That previously documented presentation correction has now been completed. The
monthly chart uses a monetary axis and preserves the validated January--August
Actual / September--December Accrual pattern. The Business Unit slicer has also
been tested interactively across ExEd, Qualification Programs, University Programs
and the explicit Unassigned member. The selected Business Unit consistently filters
the core financial measures, CRM governance measures and monthly forecast visual.
In particular, the Unassigned member isolates exactly EUR 42,550.00 of Accrual
with no Actual or CRM contribution, providing a strong report-level demonstration
that the governed Business Unit relationships operate as intended.


The Business Unit forecast-composition visual has now also been completed. The
Executive Dashboard can therefore show the same EUR 21,251,093.59 base forecast
by organizational ownership while separating realized Actual Net Sales from the
remaining Project Accrual Forecast. This reuses the validated conformed Business
Unit dimension and introduces no new forecast assumption.

The Forecast Governance table is also now present on the executive page. It keeps
the five forecast treatments in one auditable view: the EUR 21.251M core base
forecast, EUR 807.9K Ready CRM inclusion, EUR 13.5K held Review amount,
EUR 1.989M At-Risk weighted pipeline and EUR 50K missing-Probability exposure.
This makes the distinction between included, held and not-included value explicit
for an executive audience.

The analytical composition of the Executive Dashboard is therefore close to final.
The next control is an end-to-end interaction pass across the complete page, followed
by technical-field hiding, layout polish and preparation of the short interview
story that explains how source reconciliation, data quality and forecast governance
connect to the final management view.

The goal is to move from a manual Excel-based forecasting process toward
a trusted, governed and scalable decision-support solution.



The final end-to-end Business Unit interaction pass has now also been completed.
All four governed Business Unit selections were re-tested after the Forecast Governance,
Business Unit composition and Top-10 Project visuals were in place. The complete page
responds consistently, including the explicit Unassigned / Missing BU case.

A final relationship audit was also completed in \texttt{Manage relationships}. The
model contains nine governed relationships: eight active relationships and one
intentionally inactive CRM program-start date role. The active paths remain
single-directional from conformed dimensions to facts, preserving the intended
star-schema behavior.

Technical-field cleanup has also been performed. Helper keys, sort columns and selected
data-quality diagnostics remain in the semantic model for governance but are hidden from
the business-facing field list where appropriate.

The next and final technical control is the financial reconciliation sign-off: verify zero
variance for the Base Full-Year Forecast equation and for the Forecast Including Ready CRM
equation. After that check, only presentation polish, default-filter reset, final save and
the interview narrative remain.

\end{summarybox}

\vfill

\begin{tcolorbox}[
    colback=AaltoNavy,
    colframe=AaltoNavy,
    boxrule=0pt,
    arc=2mm,
    left=4mm,
    right=4mm,
    top=3mm,
    bottom=3mm
]

\centering
\small
\color{white}

\textbf{Confidential Interview Preparation Material}
\\[2pt]

Prepared solely for the Aalto EE Data Analyst/Engineer recruitment
process. Source datasets and internal process information should not be
distributed or published.

\end{tcolorbox}

\end{document}
